% DER_VI-a: Grothendieck's text with the commentary set in, then the overview, inside a group that applies the chunk's own definitions
\dgrchapter{VI}{Chapter VI}{HOT}
\dgrpart{0}{VI-a}{Part 1 of 3, manuscript pages 1--70}{VI-a\enspace\textperiodcentered\enspace Chapter VI, part 1 of 3\enspace\textperiodcentered\enspace manuscript pages 1--70}{Chapter VI \textperiodcentered\ \textup{VI-a}}
\dgrchunkbegin
% DER_VI-a: the definitions of the chunk's preamble, applied inside its group
\let\printedpage\relax\newcommand{\printedpage}[1]{\marginpar{\footnotesize p.~#1}}
\let\msp\relax\newcommand{\msp}[1]{\marginpar{\footnotesize [#1]}}
\let\deq\relax\newcommand{\deq}{\overset{\text{déf}}{=}}
\let\ul\relax\newcommand{\ul}[1]{\underline{\mathrm{#1}}}
\let\xline\relax\newcommand{\xline}[1]{\mathrel{\overset{#1}{\rule[0.6ex]{2em}{0.4pt}}}}
\let\cat\relax\newcommand{\cat}[1]{\underline{\mathrm{#1}}}
\let\circled\relax\newcommand{\circled}[1]{\tikz[baseline=(c.base)]\node[draw,circle,inner sep=1pt](c){#1};}
\let\und\relax\newcommand{\und}[1]{\underline{\mathrm{#1}}}
\let\Cat\relax\newcommand{\Cat}{\mathrm{Cat}}
\let\uu\relax\newcommand{\uu}[1]{\underline{\mathrm{#1}}}
\let\lline\relax\newcommand{\lline}{\mathrel{\relbar\joinrel\relbar\joinrel\relbar}}
\let\uc\relax\newcommand{\uc}[1]{\underline{\mathrm{#1}}}
\let\HOT\relax\newcommand{\HOT}{\mathrm{HOT}}
\let\defeq\relax\newcommand{\defeq}{\overset{\text{déf}}{=}}
\let\cercle\relax\newcommand{\cercle}[1]{\tikz[baseline=(c.base)]\node[draw,circle,inner sep=1pt](c){#1};}
\let\Ob\relax\newcommand{\Ob}{\mathrm{Ob}}
\let\Fl\relax\newcommand{\Fl}{\mathrm{Fl}}
\let\id\relax\newcommand{\id}{\mathrm{id}}
\let\cA\relax\newcommand{\cA}{\mathcal{A}}

% ---- from Der-ChVI-345_p001-005.tex
% ===================== page 1 of the PDF (title page) =====================
\begin{center}
\printedpage{[sans numéro]}% no page number printed on this page
% head: (none)

\vspace*{6cm}

{\Large\textsf{LES DÉRIVATEURS}}

\bigskip

\textbf{ALEXANDRE GROTHENDIECK}

\vspace{2cm}

\textbf{Chapitre VI}

\medskip

\textbf{HOT}
\end{center}

\clearpage

% ===================== page 2 of the PDF (editorial note) =====================
\vspace*{6cm}

\noindent\printedpage{[sans numéro]}% no page number printed on this page
% head: (none)
Ce texte a été déchiffré et transcrit en \LaTeXe{} par M. Künzer. Il a été édité par M. Künzer, J. Malgoire, et G. Maltsiniotis. La transcription est aussi fidèle que possible au manuscrit. Pour les quelques corrections évidentes, ou rares commentaires des éditeurs, ainsi que pour la numérotation originale des pages du manuscrit, les caractères de machine à écrire \texttt{[typewriter]} entre crochets sont utilisés. Un point d'interrogation entre crochets signifie que l'on n'est pas sûr du mot qui précède.

\noindent Cette édition est provisoire. Les remarques, commentaires et corrections sont bienvenus. Envoyer un message à :

\noindent\textbf{[e-mail address removed]}

\hfill G. Maltsiniotis\hspace*{2cm}

\clearpage

% ===================== printed page 1 =====================
\noindent\printedpage{1}
% head: A. Grothendieck : Dérivateurs, Ch. VI, Édition des Universités de Montpellier II et Paris VII
% foot: Version 9 novembre 2001
\msp{page 1}

\section*{\normalsize Le dérivateur \textmd{HOT}}

\begin{flushright}
14.11.90
\end{flushright}

\begin{flushright}
(Les premières feuilles, d'il y a deux semaines, ont été larguées, dépassées par VII)
\end{flushright}

\noindent Je définis Hot à partir de la catégorie de modèles
\[
\mathcal{M} = \mathrm{Cat},\ W \subseteq \mathrm{Fl}(\mathrm{Cat}),
\]
$W$ étant en principe les Hot-équivalences. Mais je vais dans la mesure du possible axiomatiser plus, en supposant donné un ensemble $W$ plus général, appelé `localisateur fondamental' (basic localiser), p.~ex. $W = W_{\mathbf{D}}$, où $\mathbf{D}$ est un dérivateur déjà donné (mais en supposant que son domaine soit (Cat) tout entier). Nous allons noter au fur et à mesure de quelles propriétés de $W$ nous aurons besoin. Aux fins de référence ultérieure, je vais rappeler ici celles que j'ai dégagées dans Pursuing Stacks.

\noindent\textbf{L1}\ \ Saturation de $W$

\begin{itemize}[leftmargin=2.5em]
\item[a)] \textit{Les identités} \texttt{[sont]} \textit{dans $W$.}
\item[b)] \textit{Si deux parmi $gf$, $f$, $g$} \texttt{[sont]} \textit{dans $W$, le troisième aussi, $X \xrightarrow{f} Y \xrightarrow{g} Z$.}
\item[c)] \textit{Si $gf$ et $fg$ sont dans $W$, alors $f$ et $g$ aussi, $X \overset{f}{\underset{g}{\rightleftarrows}} Y$ .}
\end{itemize}

\noindent\textbf{L2}\ \ $\mathbf{\Delta}^1$ \textit{est universellement $W$-asphérique, i.e. $\mathbf{\Delta}^1 \times X \longrightarrow X$ est dans $W$ pour tout $X$ dans} Cat. Cela implique

\noindent\textbf{L2 bis}\ \ $\mathbf{\Delta}^1$ \textit{est $W$-asphérique.} (C'était L3$'$ dans Pursuing Stacks.)

\noindent\textbf{L3}\ \ \textit{Si $X$ a un objet final, $X$ est $W$-asphérique.} (\footnote[1]{Il y a un L3$'$ dual de L3.})

\noindent Je note ici \texttt{[des]} variantes renforcées.

\noindent\textbf{L3 bis}\ \ \textit{Si $X$ est} Hot-\textit{asphérique, il est $W$-asphérique.}

\noindent\textbf{L3 ter}\ \ \textit{Si $X$ est} Hot-\textit{asphérique, il est universellement $W$-asphérique, i.e. $X \times Y \xrightarrow{\mathrm{pr}_2} Y$ est dans $W$ pour tout $Y$ dans} Cat. (\footnote[2]{Suffit pour entraîner L2, L3.})

\clearpage

% ===================== printed page 2 =====================
\noindent\printedpage{2}
% head: A. Grothendieck : Dérivateurs, Ch. VI, Édition des Universités de Montpellier II et Paris VII
% foot: Version 9 novembre 2001
\textbf{L3 quater}\ \ \textit{Toute} Hot-\textit{équivalence est dans $W$.}

\noindent L3 quater $\Longrightarrow$ L3 ter $\Longrightarrow$ L3 bis $\Longrightarrow$ L3.

\noindent L3 ter $\Longrightarrow$ L2 $\Longrightarrow$ L2 bis.

\noindent\msp{page 2}

\dgrblock{dgrb193}{The axioms are being found}
\begin{dcomment}{context}{p1}{The axioms are being found}
Grothendieck wants to construct the homotopy theory of diagrams from categories themselves. The expression “catégorie de modèles” here initially means a category supplied with distinguished weak equivalences; no classes of Quillen cofibrations and fibrations have yet been specified. We write W\ensuremath{\infty} for functors inducing weak homotopy equivalences of nerves, and Q : Cat \ensuremath{\to}Cat[W \ensuremath{-}1] for ordinary localization. Categories and localizations are understood in fixed, sufficiently large universes. This avoids mistaking the construction of Q for a smallness theorem about its hom sets.

The changing list of axioms is part of the argument. In particular, the L4 of manuscript page 2 is the absolute assertion that an aspherical functor is a weak equivalence. The relative assertion for a map X \ensuremath{\to}Y over an arbitrary third category S appears only as L4 bis on page 44. Later proofs require this stronger formulation. The modern definition uses the relative statement from the outset; compare Cisinski (2004, Definition 1.1.3 and Remark 1.1.4). Similarly, the L5 on page 2 concerns descent of fibrewise equivalences, whereas the L5 on page 43 concerns recovery of fibrewise equivalences from a total equivalence. They are different assertions, with different strength. The interval [1] explains why these axioms can remember homotopy. A natural transformation F \ensuremath{\Rightarrow}G is a functor A \ensuremath{\times} [1] \ensuremath{\to}B with prescribed endpoint restrictions. L2 makes both endpoints inverse to the projection after localization. Thus a natural transformation becomes an equality between the induced maps of localized categories. This observation will repeatedly turn an elementary adjunction into a weak equivalence without requiring a model structure.

The optional L3 quater also has a precise later history. Under the modern relative axioms, the inclusion W\ensuremath{\infty}\ensuremath{\subseteq}W is a theorem, established by Cisinski (2004, Theorem 2.2.11). That theorem does not license replacing the initial absolute L4 by its relative form without explanation. In these pages the additional hypothesis is explicitly adopted at manuscript page 7, and we retain the distinction when tracing the proofs.

\end{dcomment}


\begin{itemize}[leftmargin=2.5em]
\item[\textbf{L4}] \textit{Si $f : X \longrightarrow Y$ est $W$-asphérique (i.e. les $X/y \longrightarrow Y/y$ dans $W$, i.e. (si on a L3) si les $X/y$ sont $W$-asphériques), alors $f$ aussi.} (Axiome de localisation.) (\footnote[3]{il y a un L4$'$ dual de L4.}) \texttt{[`aussi'} \texttt{= dans} $W$\texttt{]}

\item[\textbf{L5}] (Axiome de fibration). \textit{Si}
\[
\begin{tikzcd}[column sep=small]
X \arrow[rr, "f"] \arrow[dr] & & Y \arrow[dl] \\
& S &
\end{tikzcd}
\]
\textit{avec $X$ et $Y$ fibrées sur $S$ et $f$ cartésien, si pour tout $s \in S$ $f_s : X_s \longrightarrow Y_s$ est dans $W$, $f$ aussi.}
\end{itemize}

\noindent Dans Pursuing Stacks j'ai supposé toujours \fbox{L1, L3, L4} (impliquant L2) - c'étaient ces conditions qui définissaient la notion de `localisateur' (fondamental) dans Cat. (\footnote[4]{Si $W$ \texttt{[est]} un localisateur, alors les catégories contractiles sont $W$-asphériques, les homotopismes sont dans $W$, les applications homotopes sont égales dans (Cat)$W^{-1} \deq \mathrm{Hot}_W$.}) D'autre part, \fbox{L1, L2 bis, L5} impliquent L1,2,3,4,5.

\noindent\textit{On va supposer ici aussi qu'on ait au moins} L1, L3, L4 (i.e. $W$ un localisateur). En appliquant la construction de I, p.~72-\textbf{75,} on voit que $f : X \longrightarrow Y$ est dans $W$ si et seulement si $f^o : X^o \longrightarrow Y^o$ l'est. Donc L1, L3, L4 impliquent L3$'$, L4$'$ (axiomes duaux de L3, L4).

\noindent Il y avait des axiomes complémentaires

\begin{itemize}[leftmargin=2.5em]
\item[\textbf{La}] $f \in W \Longrightarrow \pi_0(f)$ \textit{bijectif.}

\item[\textbf{Lb}] \textit{Le foncteur canonique de localisation}
\[
\mathrm{Cat} \longrightarrow (\mathrm{Cat})W^{-1} \deq \mathrm{Hot}_W
\]
\textit{commute aux sommes finies.}

\item[\textbf{Lc}] \textit{Itou pour produits finis.}
\end{itemize}

\clearpage

% ===================== printed page 3 =====================
\noindent\printedpage{3}
% head: A. Grothendieck : Dérivateurs, Ch. VI, Édition des Universités de Montpellier II et Paris VII
% foot: Version 9 novembre 2001
Il faudrait quand même dégager des conditions concrètes sur $W$ pour la validité de Lb, Lc. Les conditions L3 quater et La s'écrivent respectivement
\[
W_\infty \subseteq W,\ \ W \subseteq W_0.
\]

\noindent\msp{page 3}

\dgrblock{dgrb194}{What the first fibration axiom already gives}
\begin{dcomment}{exposition}{p2}{What the first fibration axiom already gives}
Grothendieck’s short implication from L1, L2 bis and the original L5 can be expanded into a useful collection of elementary arguments. It also clarifies why that first L5 becomes dispensable after relative descent has been imposed. First apply L5 to the cartesian functor [1] \ensuremath{\times} A \ensuremath{\to}A between fibrations over A. Its fibre map is [1] \ensuremath{\to}\{\ensuremath{*}\}, so L2 follows. For the projection p : A \ensuremath{\times} [1] \ensuremath{\to}A and its sections e0, e1, two-out-of-three gives e0, e1 \ensuremath{\in}W. If H : A \ensuremath{\times} [1] \ensuremath{\to}B is a natural transformation between F and G, it follows that F \ensuremath{\in}W if and only if G \ensuremath{\in}W: both are equivalent to H \ensuremath{\in}W. Moreover, Qe0 = Qp\ensuremath{-}1 = Qe1, whence QF = QG. An adjoint pair F \ensuremath{\dashv}G has GF and FG connected to the corresponding identities by its unit and counit. Both composites belong to W, and L1(c) gives F, G \ensuremath{\in}W. Applying this to the terminal-object adjunction proves L3, and the initial-object case follows as well.

There is a particularly concrete way to obtain invariance under opposites without assuming it. Let T(A) have objects the arrows a \ensuremath{\to}b of A, with a morphism from a \ensuremath{\to}b to a\ensuremath{\prime} \ensuremath{\to}b\ensuremath{\prime} given by a\ensuremath{\prime} \ensuremath{\to}a and b \ensuremath{\to}b\ensuremath{\prime} making the evident rectangle commute. The two projections

T(A)op \ensuremath{-}\ensuremath{\to}A,

T(A)op \ensuremath{-}\ensuremath{\to}Aop are fibrations. Their fibres are respectively (a\textbackslash{}A)op and A/b, each with a terminal object. L5, applied to each projection as a cartesian functor to the identity fibration, makes both projections elements of W. Their naturality in A and two-out-of-three now imply

T(f)op \ensuremath{\in}W

f op \ensuremath{\in}W. f \ensuremath{\in}W \ensuremath{\Leftarrow}\ensuremath{\Rightarrow} \ensuremath{\Leftarrow}\ensuremath{\Rightarrow} Here the fibration assertion can be checked directly: a change of the first endpoint is made by precomposition before taking opposites, and a change of the second endpoint by postcomposition. These are the specified cartesian lifts after reversal.

Now let f : X \ensuremath{\to}Y be a map over S such that every f/s : X/s \ensuremath{\to}Y/s belongs to W. The construction \ensuremath{\Psi} on manuscript pages 17--20 produces split opfibrations with these maps as their fibre maps, and \ensuremath{\Psi}(f) preserves the chosen cocartesian lifts. After taking opposites, L5 applies, so \ensuremath{\Psi}(f) \ensuremath{\in}W. The natural maps X \ensuremath{\to}\ensuremath{\Psi}(X) and Y \ensuremath{\to}\ensuremath{\Psi}(Y ) are homotopy equivalences by the adjunction exhibited below; hence two-out-of-three gives f \ensuremath{\in}W. This proves even L4 bis, and therefore L4. Each ingredient is independent of the later, much stronger L5. The argument verifies the implication recorded here while anticipating only constructions that Grothendieck gives explicitly in the next few pages.

\end{dcomment}


\noindent Leur conjonction est l'axiome des plus raisonnables
\[
W_\infty \subseteq W \subseteq W_0.
\]
Bien sûr, l'axiome L5 a maintenant l'air un peu étriqué, i.e. faiblard. On va voir très vite de quoi on aura besoin!

\noindent On veut donc définir un dérivateur
\[
\mathrm{HOT}_W
\]
à partir de la catégorie de modèles
\[
(\mathrm{Cat}, W \subseteq \mathrm{Fl}(\mathrm{Cat})).
\]
Donc pour définir
\[
\mathrm{HOT}_W(S) \qquad\qquad\qquad (S \in \mathrm{Ob}\,\mathrm{Cat}),
\]
la récette générale est de prendre la catégorie % [sic] récette -> recette, misspelling of "recette"
\[
\mathcal{M}(S) = \mathrm{Cat}(S) = \underline{\mathrm{Hom}}(S^o, (\mathrm{Cat})),
\]
d'y introduire l'ensemble de flèches
\[
\left\{
\begin{aligned}
& W_S \subseteq \mathrm{Fl}(\mathcal{M}(S)) \\
& u \in W_S \Longleftrightarrow u_s \in W \ \ \forall\, s \in \mathrm{Ob}\,S.
\end{aligned}
\right.
\]
On doit poser
\[
\mathrm{HOT}_W(S) \deq \underline{\mathrm{Hom}}(S^o, \mathrm{Cat})W_S^{-1}.
\]
Mais la catégorie $\underline{\mathrm{Hom}}(S^o, \mathrm{Cat})$ s'interprète, à isomorphisme canonique près, comme celle
\[
\underline{\mathrm{Fibsc}}\,S
\]
des catégories $X$ Cat-fibrées sur $S$, munies d'un `scindage', i.e. d'un choix de foncteurs changement de base $u^* : X_t \longrightarrow X_s$ (pour $t \xrightarrow{u} s$ dans $S$) \texttt{[plutôt} $u^* : X_s \longrightarrow X_t$\texttt{]}, avec de plus \textit{transitivité stricte}
\[
(vu)^* = u^* v^* \qquad\qquad\qquad \text{si}\ \ t \xrightarrow{u} s \xrightarrow{v} r\ \ \text{dans}\ S
\]

\noindent\msp{page 4}

\dgrblock{dgrb195}{Products and sums survive the localization}
\begin{dcomment}{exposition}{p3}{Products and sums survive the localization}
Grothendieck asks for concrete conditions ensuring Lb and Lc because a homotopy theory should still allow objects to be assembled and compared. A useful answer is available as soon as relative L4 bis is admitted. We give the argument, including the passage from closure of weak equivalences to the universal property of products and coproducts.

Suppose W satisfies L1, L2 and L4 bis. If f : A \ensuremath{\to}B lies in W, then f \ensuremath{\times} idC lies in W for every C. View this map over C. Its comma map at c is f \ensuremath{\times} idC/c. Since C/c has a terminal object, the two projections A \ensuremath{\times} (C/c) \ensuremath{\to}A and B \ensuremath{\times} (C/c) \ensuremath{\to}B are homotopy equivalences: the terminal object defines sections and a contraction natural transformation. They are in W by the preceding block. Two-out-of-three makes the comma map a weak equivalence, and L4 bis gives the assertion. The same axiom, applied over a discrete indexing category I, shows that ` i fi \ensuremath{\in}W whenever every fi \ensuremath{\in}W. The latter closure is also Cisinski (2004, Proposition 1.1.5).

For completeness, localization of a product of relative categories is the product of their localizations. One can see this directly from the generators: (w, id) and (id, w\ensuremath{\prime}) become invertible, their commutation squares still commute after inversion, and a pair of finite zigzags can be interleaved to give a zigzag in the product. Consequently both the diagonal and the binary product functor descend. The adjunction \ensuremath{\Delta}\ensuremath{\dashv}\ensuremath{\times} descends with its original unit and counit, since those are natural transformations between relative functors and the triangle identities remain identities. Its descended right adjoint is therefore a binary product in Cat[W \ensuremath{-}1]. The same argument with ` \ensuremath{\dashv}\ensuremath{\Delta}gives binary coproducts. The adjunctions defining the initial and terminal objects descend without any closure hypothesis beyond identities belonging to W.

Thus finite products and coproducts are preserved by Q under these concrete assumptions. This is an ordinary localization argument, not a claim that arbitrary homotopy limits or colimits agree with ordinary ones. In particular, it does not supply the derived Kan extensions needed for the derivator. Those require the diagram categories whose construction begins here.

\end{dcomment}


% ---- from Der-ChVI-345_p006-010.tex
\printedpage{4}
% head: A. Grothendieck : Dérivateurs, Ch. VI, Édition des Universités de Montpellier II et Paris VII
et non seulement transitivité à isomorphisme canonique près. Ainsi on peut écrire aussi
\[
\mathrm{HOT}_W(S) \xrightarrow{\ \sim\ } (\ul{Fibsc}\,S)W_S^{-1}, \tag*{$(*)$}
\]
où cette fois $W_S$ est interprété comme
\[
W_S \subseteq \mathrm{Fl}(\ul{Fibsc}\,S),
\]
savoir
\[
W_S = \{u \in \mathrm{Fl}(\ul{Fibsc}\,S) \mid u_s \in W \text{ pour tout } s \in \mathrm{Ob}\,S\}
\]
(foncteurs strictement cartésiens, i.e. commutant aux scindages, qui induisent des $W$-équivalences entre les catégories fibres).

La flèche dans $(*)$ est bel et bien un isomorphisme de catégories, pas seulement une équivalence. Mais on va définir d'autres catégories, plus commodes à certains égards, qui seront \textit{équivalentes} (on l'espère), mais sûrement pas isomorphes, à $\mathrm{HOT}_W(S)$. (On s'occupera plus tard de la variance de $\mathrm{HOT}_W(S)$ en $S$.)

Les catégories $\ul{Hom}(S^o,\mathrm{Cat})$, $\ul{Fibsc}\,S$, ainsi que les variantes que nous allons introduire maintenant, sont en fait des 2-catégories. Mais pour décrire $\mathrm{HOT}_W(S)$, seule les structures de 1-catégories sont utilisées. On a l'impression que c'est bien % [sic] seule les structures -> seules les structures, accord du pluriel
\msp{page 5}
grossier, qu'on perd quelque chose ce faisant. Mais ce qui est perdu va se retrouver, et surabondamment, dans la structure des \underline{Hom} internes dans les $\mathrm{HOT}(S)$, si on se `rappelle' que ces \underline{Hom} internes, donnant lieu à des `types d'homotopie relatifs' $\ul{Hom}_{\mathrm{HOT}_W}(S;\xi,\eta)$ (si $\xi$, $\eta$ en sont), peuvent s'interpréter en termes de $\infty$-champs en groupoïdes sur $S$, avec toute la richesse de structure qui leur appartient. Mais ceci est une autre histoire\dots

\dgrblock{dgrb196}{A diagram written as its category of elements}
\begin{dcomment}{exposition}{p4}{A diagram written as its category of elements}
Grothendieck passes to fibrations in order to replace an externally indexed diagram by one category carrying a projection to its index. This is the same organizing construction used in descent; the relevant background is Grothendieck (1971, Exposé VI, \S{}\S{}8--9). It is useful to write out the variance because both directions of the construction will soon be needed.

For F : Sop \ensuremath{\to}Cat, its Grothendieck construction has objects (s, x) with x \ensuremath{\in}F(s). An arrow (s, x) \ensuremath{\to}(t, y) consists of u : s \ensuremath{\to}t and a : x \ensuremath{\to}F(u)y in F(s). Composition is

(v, b) \ensuremath{\circ}(u, a) = (vu, F(u)(b) \ensuremath{\circ}a).

The chosen cartesian lift of u at (t, y) is (u, id) : (s, F(u)y) \ensuremath{\to}(t, y). The strict functoriality of F makes these choices compose strictly. Conversely, a normalized splitting of a fibration gives strict pullback functors u\ensuremath{*}, with (vu)\ensuremath{*}= u\ensuremath{*}v\ensuremath{*}and identity pullbacks equal to identities. These equations, rather than the existence of unspecified isomorphisms, are what is meant by a splitting.

A natural transformation F \ensuremath{\to}G induces a functor of the total categories that preserves the specified lifts. Conversely, such a functor restricts to the components of a strict natural transformation. Therefore the weak equivalences in the diagram category correspond exactly to fibrewise weak equivalences among splitting-preserving functors. With fibrations encoded by their fibre families, this is the identification Grothendieck calls an isomorphism. If arbitrary underlying object sets for total categories are allowed, the usual construction and reconstruction are canonically isomorphic over S, and the comparison is expressed as an equivalence. No homotopical argument depends on this bookkeeping convention.

The subsequent forgetful functor to Fib S changes the morphisms: there one allows all functors over S. Forgetting a splitting is therefore not merely forgetting redundant labels. The natural comparison maps used later must be checked in the category whose morphisms they actually preserve. This is the source of the obstruction on manuscript pages 32--39, and of the strictification work beginning on page 52.

\end{dcomment}


En fait, je vois un diagramme de 13 catégories de `categories sur $S$\,', à examiner à présent % [sic] categories -> catégories, accent manquant
\[
\begin{tikzcd}[column sep=small]
 & \ul{Fibsc}_0S \arrow[dl, hook] \arrow[d, "\mathrm{fid}"] & & \ul{Cofibsc}_0S \arrow[d] \arrow[dr, hook] & \\
\ul{Fibsc}\,S \arrow[d, "\mathrm{fid}"] & \ul{Fib}_0S \arrow[dl, hook] & & \ul{Cofib}_0S \arrow[dr, hook] & \ul{Cofibsc}\,S \arrow[d, "\mathrm{fid}"] \\
\ul{Fib}\,S \arrow[dr, hook] & & \ul{Parf}\,S \arrow[dl, hook] \arrow[d, hook] \arrow[dr, hook] & & \ul{Cofib}\,S \arrow[dl, hook] \\
 & \ul{Liss}\,S \arrow[dr, hook] & \ul{Fibhot}\,S \arrow[d, hook] & \ul{Prop}\,S \arrow[dl, hook] & \\
 & & \mathrm{Cat}/S & &
\end{tikzcd}
\]
% foot: Version 9 novembre 2001

\printedpage{5}
% head: A. Grothendieck : Dérivateurs, Ch. VI, Édition des Universités de Montpellier II et Paris VII
où \underline{Fib}, \underline{Cofib} désignent les catégories Cat-fibrées, Cat-cofibrées sur $S$, \underline{Liss} et \underline{Prop} les catégories
\msp{page 6}
lisses resp. propres sur $S$, \underline{Parf} celles qui sont à la fois lisses et propres (fibrations parfaites), enfin \underline{Fibhot} les catégories Hot-fibrées sur $S$. Toutes les flèches sont des foncteurs strictement pleinement fidèles, à l'exception des deux (plutôt quatre) flèches verticales supérieures, qui sont fidèles sans plus. Je viens d'en rajouter encore 4, de catégories : $\ul{Fibsc}_0$, $\ul{Fib}_0$, $\ul{Cofibsc}_0$, $\ul{Cofib}_0$ sont les sous-catégories pleines des catégories correspondantes sans indice, formées des catégories fibrées resp. cofibrées $X$ sur $S$, pour lesquelles les foncteurs changement de base resp. cochangement de base
\[
u^* : X_t \longrightarrow X_s, \qquad\quad \text{resp.} \qquad\quad u_* : X_s \longrightarrow X_t,
\]
pour $s \xrightarrow{u} t$ dans $S$, sont des Hot-équivalences. En termes de foncteurs
\[
S^o \xrightarrow{F} \mathrm{Cat} \qquad\quad \text{resp.} \qquad\quad S \xrightarrow{F'} \mathrm{Cat},
\]
cela signifie que ce foncteur transforme flèches quelconques dans $S$ en éléments de $W$ dans Cat, i.e. qu'on a un diagramme commutatif
\[
\begin{tikzcd}[row sep=large]
S^o \arrow[r, "F"] \arrow[d, "\mathrm{can}"'] & \mathrm{Cat} \arrow[d, "\mathrm{can}"] \\
\Pi(S^o) \simeq (\Pi S)^o = S^o(\mathrm{Fl}\,S^o)^{-1} \arrow[r, "\bar F"'] & \mathrm{Hot}_W = (\mathrm{Cat})W^{-1}
\end{tikzcd}
\]
(et de même pour $F'$), i.e. qu'on a un `système
\msp{page 7}
local' $\bar F$ sur $S$, à valeurs dans $\mathrm{Hot}_W$, défini à l'aide de $F$. Je rappelle qu'on trouve également un tel système local dans le cas d'une catégorie Hot-fibrée sur $S$ (et même un système local à valeurs dans Hot lui-même, au lieu de $\mathrm{Hot}_W$ - on utilise ici tacitement que $W_\infty \subseteq W$, i.e. \fbox{L3 quater} (\footnote[5]{On va le supposer vrai dans toute la suite.})). Et il en est de même a fortiori pour les objets de $\ul{Parf}\,S$. Pour $\ul{Liss}\,S$ resp. $\ul{Prop}\,S$, on trouve simplement un foncteur
\[
S^o \xrightarrow{\bar F} \mathrm{Hot} \qquad\quad \text{resp.} \qquad\quad S \xrightarrow{\bar F'} \mathrm{Hot},
\]
d'où, en composant avec $\mathrm{Hot} \longrightarrow \mathrm{Hot}_W$, des foncteurs à valeurs dans $\mathrm{Hot}_W$.

Dans chacune de ces 14 \texttt{[13?]} catégories, on introduit un ensemble de flèches, qui serviront à localiser, et on le désignera toujours par la même notation $W_S$. Dans toutes ces catégories sauf seulement $\mathrm{Cat}/S$, $W_S$ désigne la $W$-équivalence `fibre par fibre'. Dans $\mathrm{Cat}/S$ d'autre part, ce sera la $W$-équivalence sans plus. Toutes les flèches du diagramme sont compatibles avec les ensembles localisants de flèches.
% foot: Version 9 novembre 2001

\printedpage{6}
% head: A. Grothendieck : Dérivateurs, Ch. VI, Édition des Universités de Montpellier II et Paris VII
\msp{page 8}
\[
\begin{tikzcd}[column sep=small, row sep=large]
 & \ul{Fibsc}_0S \arrow[dl, hook, "\alpha'"'] \arrow[d] & & \ul{Cofibsc}_0S \arrow[d] \arrow[dr, hook, "\beta'"] & \\
\ul{Fibsc}S \arrow[d] & \ul{Fib}_0S \arrow[dl, hook, "\alpha"'] \arrow[ddr, dotted, bend right=20] & & \ul{Cofib}_0S \arrow[dr, hook, "\beta"] \arrow[ddl, dotted, bend left=20] & \ul{Cofibsc}\,S \arrow[d] \\
\ul{Fib}S \arrow[dr, hook] & & \ul{Parf}\,S \arrow[dl, hook] \arrow[d, hook] \arrow[dr, hook] & & \ul{Cofib}\,S \arrow[dl, hook] \\
 & \ul{Liss}\,S \arrow[dr, hook] & \ul{Fibhot}\,S \arrow[d, hook] \arrow[uuul, dotted, bend left=10, "\Phi_0"] \arrow[uuur, dotted, bend right=10, "\Psi_0"'] & \ul{Prop}\,S \arrow[dl, hook] & \\
 & & \mathrm{Cat}/S \arrow[uuull, dotted, bend left=60, "\Phi"] \arrow[uuurr, dotted, bend right=60, "\Psi"'] & &
\end{tikzcd}
\]
% ?? uncertain: starting point of the dotted arrows labelled Phi_0 and Psi_0 (the dotted lines run alongside Parf S down to Fibhot S) -- Fibhot S or Parf S

\dgrblock{dgrb197}{A local system in the homotopy category misses twisting}
\begin{dcomment}{examples}{p7}{A local system in the homotopy category misses twisting}
Grothendieck’s passage from F to F\textasciimacron{} : \ensuremath{\Pi}(Sop) \ensuremath{\to}HotW records monodromy, but he also signals on page 5 that richer coherence remains. A familiar pair of bundles makes the difference visible. We use ordinary spaces to display the phenomenon; replacing spaces by their simplex categories gives categorical models with the same homotopy information.

Consider the trivial circle bundle S2 \ensuremath{\times} S1 \ensuremath{\to}S2 and the Hopf bundle (Hatcher, 2002, Examples 4.45 and 4.51)

S3 = \{(z0, z1) \ensuremath{\in}C2 : |z0|2 + |z1|2 = 1\} \ensuremath{-}\ensuremath{\to}CP1 \ensuremath{\sim}= S2, (z0, z1) 7\ensuremath{-}\ensuremath{\to}[z0 : z1]. On the chart z0 = 0, with coordinate w = z1/z0, a trivialization is

\ensuremath{\lambda}

\ensuremath{\lambda} \ensuremath{\in}S1. (w, \ensuremath{\lambda}) 7\ensuremath{-}\ensuremath{\to} 1 + |w|2 (1, w), p The other chart is analogous. Thus both maps are locally trivial bundles with fibre S1, and hence fibrations. Since the base is simply connected, their parallel transport determines the same functor from the fundamental groupoid to the homotopy category: a constant local system with value S1. Nevertheless they are not fibre homotopy equivalent. Such an equivalence would induce a weak equivalence of total spaces by the homotopy exact sequences, whereas

H1(S2 \ensuremath{\times} S1; Z) = Z.

H1(S3; Z) = 0, The missing information is the coherence of transport around two-dimensional homotopies, not an additional choice of monodromy around loops.

In contemporary language the appropriate datum is a functor from the entire homotopy type of the base, considered as an \ensuremath{\infty}-groupoid, to spaces. Passing to the fundamental groupoid of the base and the ordinary homotopy category of spaces discards part of that datum. Straightening and unstraightening provide the relevant modern comparison; Lurie (2009, Theorems 2.2.1.2 and 3.2.0.1) treat respectively space-valued and category-valued versions. Here this explains the direction of Grothendieck’s remark about \ensuremath{\infty}-groupoid stacks, without asserting that a single ordinary category HOTW (S) already determines every higher structure he envisages. The commentary on the earlier chapters, on what a diagram theory remembers, remains applicable.

There is also a small logical distinction in the definition on page 6. If the subscript 0 means that all transition functors lie in W\ensuremath{\infty}, their images lie in W once W\ensuremath{\infty}\ensuremath{\subseteq}W. The converse need not hold for a larger W. A W-locally constant diagram and an ordinarily homotopically locally constant diagram must therefore be distinguished whenever W is still variable.

\end{dcomment}


\texttt{[Dans ce diagramme, il y a aussi} $\ul{Liss}_0$ \texttt{et} $\ul{Prop}_0$ \texttt{mais sans connexion aux autres objets.]}

Je suis quasiment sûr du

\textbf{Theorème.} (\footnote[6]{C'est totalement canulé, cf. commentaires page 32 ff., solution des perplexités pages 41 ff. Voir surtout le Théorème-Scholie récapitulatif p. 70-81.}) \textit{Tous les foncteurs précédents induisent des équivalences entre les catégories localisées.} % [sic] Theorème -> Théorème, accent manquant

Comme il y a 16 foncteurs, ça a l'air d'un exercise de patience. Le principe général, c'est qu'on arrive à définir des foncteurs en sens inverse, compatibles avec les $W_S$, de telle façon que les composés soient reliées aux foncteurs identiques par des flèches qui sont dans $W_S$, de sorte que par passage aux localisés ces foncteurs induisent des (équivalences) quasi-inverses des foncteurs déduits de ceux du diagramme. Tous ces foncteurs en sens opposé résultent essentiellement par composition, de trois foncteurs bien familiers. Deux de ceux-ci sont indiqués en pointillés sur le diagramme, comme $\Phi$, $\Psi$. % [sic] exercise -> exercice, orthographe française ; % [sic] reliées -> reliés, accord avec « les composés »

\msp{page 9}
Le troisième est un foncteur canonique
\[
C : \mathrm{Cat}/S \longrightarrow \ul{Parf}\,S,
\]
qu'on va expliciter également.
% foot: Version 9 novembre 2001

\dgrblock{dgrb198}{The walking arrow detects the failed equivalence}
\begin{dcomment}{examples}{p8}{The walking arrow detects the failed equivalence}
Grothendieck’s first comparison theorem would identify all the categories in the large diagram after the specified localizations. The later marginal note withdraws it. The obstruction already occurs over the smallest base with a noninvertible arrow and explains why the distinction between fine and coarse equivalence occupies the rest of the chapter.

Let S = [1] = (0 \ensuremath{\to}1) and take W = W\ensuremath{\infty}. Consider the contravariant diagrams

F(1) = \ensuremath{\varnothing},

F(0) = \{\ensuremath{*}\}, G(0) = G(1) = \{\ensuremath{*}\}. There is a unique natural transformation F \ensuremath{\to}G. Its Grothendieck construction is the functor \{0\} ,\ensuremath{\to}[1] over S. Both source and target are split fibrations, and both total categories have contractible nerve, so this map is a coarse weak equivalence. Its fibre at 1 is \ensuremath{\varnothing}\ensuremath{\to}\{\ensuremath{*}\}, which is not a weak equivalence.

This difference persists after localization; it is not merely a difference between two presenting classes of arrows. Evaluation at 1 descends from the category of strict diagrams localized at objectwise weak equivalences to Hot. The displayed map has image \ensuremath{\varnothing}\ensuremath{\to}\{\ensuremath{*}\} there and therefore cannot become an isomorphism. Its coarse image is invertible by definition. Thus the forgetful comparison from fibrewise localization to coarse localization is not an equivalence. The same argument applies whenever W \ensuremath{\subseteq}W0, since then empty and nonempty categories cannot become isomorphic after applying \ensuremath{\pi}0.

The locally constant subcategory is also genuinely smaller. If a diagram H(1) \ensuremath{\to}H(0) has a weak equivalence as its transition and were isomorphic to F after objectwise localization, evaluation would make H(1) empty and H(0) contractible in Hot, contradicting the invertibility of its transition. Reversing all arrows gives the corresponding example for split opfibrations. This does not invalidate the formal localization lemma on the next pages: the missing hypothesis is that its comparison transformations lie in the particular weak-equivalence classes being used.

\end{dcomment}


\printedpage{7}
% head: A. Grothendieck : Dérivateurs, Ch. VI, Édition des Universités de Montpellier II et Paris VII
Le foncteur $\Phi$ (resp. $\Psi$) appliqué à $\ul{Fib}_0S$ (resp. $\ul{Cofib}_0S$), induit un foncteur
\[
\ul{Fib}_0S \longrightarrow \ul{Fibsc}_0S \qquad\quad \text{resp.} \qquad\quad \ul{Cofib}_0S \longrightarrow \ul{Cofibsc}_0S \tag*{$(*)$}
\]
(ce qu'il faudra vérifier). Cela fournit toutes les flèches `inverses' dont on a besoin, à l'exception seulement de $\alpha$, $\beta$, $\alpha'$, $\beta'$. Mais si on a prouvé, par la méthode indiquée, que toutes les flèches du diagramme induisent des équivalences sur les catégories localisées, pour prouver qu'il en est encore de même pour $\alpha$, $\beta$, $\alpha'$, $\beta'$, il suffit de le prouver pour les deux foncteurs
\[
\ul{Fibsc}_0S \longrightarrow \ul{Fib}\,S, \qquad\qquad \ul{Cofibsc}_0S \longrightarrow \ul{Cofib}\,S,
\]
donc il reste à définir les foncteurs en sens inverse
\[
\ul{Fib}\,S \longrightarrow \ul{Fibsc}_0S, \qquad\qquad \ul{Cofib}\,S \longrightarrow \ul{Cofibsc}_0S,
\]
et comme les sources sont des sous-catégories pleines de $\mathrm{Cat}/S$, il suffira de définir
\[
\Phi_0 : \mathrm{Cat}/S \longrightarrow \ul{Fibsc}_0S, \qquad\qquad \Psi_0 : \mathrm{Cat}/S \longrightarrow \ul{Cofibsc}_0S.
\]
Pour ceci, on posera
\[
\Phi_0 = \Phi r C, \qquad\qquad \Psi_0 = \Psi r C, \tag*{$(**)$}
\]
\msp{page 10}
où $C : \mathrm{Cat}/S \longrightarrow \ul{Parf}\,S$, $\Phi$ et $\Psi$ sont les 3 foncteurs `rétrogrades' inversés plus haut, et $r : \ul{Parf}(S) \hookrightarrow \mathrm{Cat}/S$ est l'inclusion. Il faut donc vérifier que $\Phi_0$, $\Psi_0$ définis par $(**)$, qui a priori ont comme catégories but $\ul{Fibsc}\,S$ et $\ul{Cofibsc}\,S$, prennent leur valeur dans les sous-catégories avec indice 0.

Je vais donner un énoncé d'équivalence unique, qui couvrira (presque) tous les 16 cas à la fois :

\textbf{Corollaire.} \textit{Considérons le diagramme}
\[
\begin{tikzcd}[column sep=small, row sep=small]
 & \ul{Parf}\,S \arrow[dd, "r"'] & \\
\ul{Fibsc}_0S \arrow[dr, "p"'] & & \ul{Cofibsc}_0S \arrow[dl, "q"] \\
 & \mathrm{Cat}/S &
\end{tikzcd}
\]
\textit{et soit $\mathcal{F}$ une sous-catégorie pleine de $\mathrm{Cat}/S$, telle que l'un des foncteurs $p$, $q$, $r$ se factorise par $\mathcal{F}$. Soit $W_{\mathcal{F}}$ induit sur $\mathcal{F}$ par $W_S$ dans $\mathrm{Cat}/S$. Alors $\mathcal{F} \hookrightarrow \mathrm{Cat}/S$ induit une équivalence sur les catégories localisées.}

\msp{page 11}
\textsc{Démonstration}. Je dis qu'il suffit de prouver, pour les trois foncteurs $p$, $q$, $r$, et pour les foncteurs en sens inverse
% foot: Version 9 novembre 2001
\printedpage{8}
% head: A. Grothendieck : Dérivateurs, Ch. VI, Édition des Universités de Montpellier II et Paris VII
\[
\begin{tikzcd}[column sep=small, row sep=small]
 & \ul{Parf}\,S & \\
\ul{Fibsc}_0S & & \ul{Cofibsc}_0S \\
 & \mathrm{Cat}/S \arrow[uu, "C"] \arrow[ul, "\Phi_0"] \arrow[ur, "\Psi_0"'] &
\end{tikzcd}
\]
que les six composés
\[
p\Phi_0, \Phi_0p, q\Psi_0, \Psi_0q, rC, Cr
\]
sont reliés aux foncteurs identiques (dans la catégorie pertinente) par une flèche qui est dans le localisateur correspondant. Pour chacun de ces trois couples de foncteurs $(p,\Phi_0)$, $(q,\Psi_0)$, $(r,C)$, cela prouvera déjà qu'en passant aux localisés, ils induisent des équivalences quasi-inverses l'une de l'autre. Mais on a mieux à l'aide du lemme formel suivant.

\textbf{Lemme.} \textit{Soient $(C,W)$, $(C',W')$ deux catégories avec localisateurs $W \subseteq \mathrm{Fl}\,C$, $W' \subseteq \mathrm{Fl}\,C'$. Soit}
\[
p : C \longrightarrow C'
\]
\textit{un foncteur compatible avec les localisateurs, et supposons qu'il existe un foncteur $\varphi : C' \longrightarrow C$ également compatible, et tel que l'on ait des flèches (à direction non précisée)}
\[
\varphi p \xline{u} id_C, \qquad\qquad p\varphi \xline{v} id_{C'},
\]
\textit{qui soient dans $W$ resp. $W'$ `argument par argument'.}

\msp{page 12}
\textit{Alors $p$ et $\varphi$ induisent des équivalences quasi-inverses l'une de l'autre en passant aux catégories localisées (elles n'ont pas de merite). De plus, pour toute sous-catégorie pleine $C'_0$ de $C'$, telle que $p$ se factorise (en $p_0 : C \longrightarrow C'_0$) par $C'_0$, si on désigne par $W'_0$ le localisateur induit dans $C'_0$ par $C'$, alors $p_0 : C \longrightarrow C'_0$ et $\varphi_0 = \varphi|C'_0 : C' \longrightarrow C$} \texttt{[plutôt} $C'_0 \longrightarrow C$\texttt{]} \textit{induisent également des équivalences quasi-inverses l'une de l'autre sur les localisés.} % [sic] merite -> mérite, accent manquant

En effet, $\varphi_0p_0 = (\varphi i)p_0 = \varphi(ip_0) = \varphi p \xline{u} \mathrm{id}_C$ (ici $i : C'_0 \hookrightarrow C'$ désigne l'inclusion), et $F \overset{\text{déf}}{=} p_0\varphi_0 = p_0\varphi i$ satisfait à
\[
iF = \underbrace{ip_0}_{p}\varphi i = (p\varphi)i.
\]
Pour définir $F \xline{v_0} \mathrm{id}_{C'_0}$, qui soit dans $W'_0$ argument par argument,
\msp{page 13}
il suffit de définir
\[
\underbrace{iF}_{=\,(p\varphi)i} \xline{\bar v_0} i,
\]
qui soit dans $W'$ argument par argument, puisque $W'_0$ induit par $W'$ et que $F \longmapsto iF$
\[
\ul{Hom}(C'_0,C'_0) \longrightarrow \ul{Hom}(C'_0,C')
\]
% foot: Version 9 novembre 2001

\dgrblock{dgrb199}{The localization lemma keeps its full force}
\begin{dcomment}{exposition}{p12}{The localization lemma keeps its full force}
Grothendieck separates a sound formal mechanism from the geometric assertions to which he hopes to apply it. The lemma is useful precisely because its hypotheses record both kinds of information: preservation of weak equivalences by the functors, and membership of the comparison maps in the specified weak-equivalence classes.

Let P : (C, W) \ensuremath{\to}(D, V ) and R : (D, V ) \ensuremath{\to}(C, W) be relative functors. Suppose natural transformations, in either direction, connect RP with idC and PR with idD, with all components respectively in W and V . Their images under localization are natural isomorphisms. The induced functors P\textasciimacron{} and R\textasciimacron{} are therefore quasi-inverse equivalences. One needs neither a calculus of fractions nor triangle identities for the original transformations. Quasi-inverse equivalences, unlike a prescribed adjunction, require only the two natural isomorphisms.

Now suppose P factors through a full subcategory i : D0 ,\ensuremath{\to}D, written P = iP0, and put V0 = V \ensuremath{\cap}Mor(D0) and R0 = Ri. The first comparison restricts to R0P0 = RP. The second restricts to a transformation between iP0R0 and i. Because both of its endpoints lie in D0, fullness gives a unique transformation between P0R0 and idD0; faithfulness verifies its naturality. Every component belongs to V0. Consequently P0 and R0 also give inverse equivalences after localization. Applying the same argument to i shows that D0[V \ensuremath{-}1 ] \ensuremath{\to}D[V \ensuremath{-}1] is an equivalence.

This explains the apparently strong corollary about an arbitrary full subcategory containing the image of a replacement functor. The weak-equivalence class on that full subcategory must be the induced one. For example, a full subcategory of Cat/S containing all perfect replacements may calculate the coarse localization. Nothing in the lemma permits changing its induced coarse class to the fibrewise class without proving their agreement. Nor does fullness apply to the passage from splitting-preserving maps to arbitrary maps over S: that functor is generally faithful but not full. These are the two points at which the first comparison diagram asks for additional mathematics.

\end{dcomment}


% ---- from Der-ChVI-345_p011-015.tex
\printedpage{9}
% head: A. Grothendieck : Dérivateurs, Ch. VI, Édition des Universités de Montpellier II et Paris VII
est pleinement fidèle ($i$ l'étant). Or il suffit de prendre
\[
\bar{v}_0 = v * i : (p\varphi)\circ i \mathrel{\relbar\joinrel\relbar} \mathrm{id}_{C'}\circ i.
\] % ?? uncertain: the arrow between (p\varphi)\circ i and id_{C'}\circ i is printed as a plain line with no visible head -- plain line / arrow

Pour résumer, il faut montrer que le composé des couples de foncteurs quasi-inverses marche dans les cinq cas suivants
\[
\begin{tikzcd}[column sep=large, row sep=huge]
\cat{Fibsc} \arrow[rrd, shift left] & \cat{Fibsc}_0 \arrow[l, hook'] \arrow[rd, shift left] & \cat{Parf} \arrow[d, shift left] & \cat{Cofibsc}_0 \arrow[r, hook] \arrow[ld, shift left] & \cat{Cofibsc} \arrow[lld, shift left] \\
& & \mathrm{Cat}/S \arrow[llu, shift left, "\Phi"] \arrow[lu, shift left, "\Phi_0"] \arrow[u, shift left, "C"] \arrow[ru, shift left, "\Psi_0"] \arrow[rru, shift left, "\Psi"] & &
\end{tikzcd}
\]

D'ailleurs par dualité, on est ramené aux trois cas de $\cat{Fibsc}$, $\cat{Fibsc}_0$, $\cat{Parf}$. (Peut-être le cas $\cat{Fibsc}_0$, avec $\Phi_0$ défini en termes de $\Phi$ et de $C$, se ramènera-t-il aux deux cas qui l'entourent, de $\cat{Fibsc}$ et $\cat{Parf}$.)

Je vois que je dois rectifier le tir, tout d'abord parce que j'ai oublié que les foncteurs
\msp{page 14}
$\Phi$, $\Psi$ que j'avais en vue ne sont \textit{pas} compatibles avec les localiseurs \texttt{[= localisateurs]} - une flèche $f : X \longrightarrow Y$ de $\mathrm{Cat}/S$ telle que la flèche correspondante dans Cat soit dans $W$, ne va pas induire une flèche dans $\cat{Fibsc}\,S$ ou $\cat{Cofibsc}\,S$ qui soit dans $W$ sur chaque fibre. Ce n'est pas grave, car on peut remplacer les foncteurs par $\Phi' = \alpha'\Phi_0$, $\Psi' = \beta'\Psi_0$, qui sauf erreur doivent marcher. Mais en vue du Lemme précédent, il y a une perplexité, car j'ai l'air d'avoir prouvé que si $\mathcal{F}$ est une sous-catégorie pleine de $\mathrm{Cat}/S$ dont je suppose seulement qu'elle contient l'image essentielle d'une des trois catégories $\cat{Fib}_0S$, $\cat{Cofib}_0S$, ou $\cat{Parf}\,S$, alors l'inclusion dans $\mathrm{Cat}/S$ induit une équivalence des catégories localisées, en prenant dans $\mathcal{F}$ comme localisateur celui induit par $\mathrm{Cat}/S$, i.e.
\begin{tikzcd}[column sep=small, row sep=small, baseline=(\tikzcdmatrixname-2-2.base)]
f : X \arrow[rr] \arrow[dr] & & Y \arrow[dl] \\
& S &
\end{tikzcd}
est dans $W_{\mathcal{F}}$ si et seulement si en tant que flèche de Cat, $f$ est dans $W$. Or ceci coïncide bien (via un théorème non trivial sur lequel je devrai revenir dans le cadre d'un localisateur fondamental général $W$) avec l'ensemble de flèches de localisation `fibre par fibre', dans le cas de $\mathcal{F} = \cat{Parf}\,S$, $\cat{Fibhot}\,S$, $\cat{Fib}_0S$, $\cat{Cofib}_0S$.

\msp{page 15}
Par contre, ce n'est nullement le cas en prenant $\mathcal{F} = \cat{Fib}\,S, \cat{Liss}\,S$ ou $\cat{Cofib}\,S$. Ayant choisi dans ces catégories comme localisateurs la notion de $W$-équivalences `fibre par fibre', j'ai donc l'air de dire que je trouve pour ces quatre catégories les \textit{mêmes} localisées (à équivalence près) que si je prenais le localisateur beaucoup plus grossier hérité par $\mathrm{Cat}/S$. Cela semble difficile à croire, et il faudra que je fasse très attention en démontrant les énoncés-clef.

Pour voir que les flèches en traits pleins du grand diagramme p.~6 sont toutes compatibles avec les localis(at)eurs \texttt{[sic]}, donc aussi les flèches composées au moyen de celles-ci, j'ai à utiliser déjà ceci (pour le cas des flèches $\cat{Liss} \longrightarrow \mathrm{Cat}/S$, $\cat{Prop} \longrightarrow \mathrm{Cat}/S$):
% foot: Version 9 novembre 2001

\printedpage{10}
% head: A. Grothendieck : Dérivateurs, Ch. VI, Édition des Universités de Montpellier II et Paris VII
\circled{1}\ \ Si $X, Y$ lisses ou propres sur $S$, et $f : X \longrightarrow Y$ un $S$-foncteur, si $f_s : X_s \longrightarrow Y_s$ est dans $W$ pour tout $s \in S$, alors $f \in W$. (\footnote[7]{OK si on suppose L3 quater et L4.}).

Cela renforce considérablement la condition L5 sur $W$, où on supposait que $f$ est un foncteur \textit{cartésien} entre catégories \textit{fibrées} sur $S$. (À dégager, pour les différentes propriétés dont on a besoin pour $W$, à quelles propriétés standard on peut les ramener. Tout sera OK si $W = W_\infty$.)

\msp{page 16}
I) \textbf{Définition de} $\Phi$, $\Psi$.

\dgrblock{dgrb200}{The precise local-to-global argument}
\begin{dcomment}{exposition}{p15}{The precise local-to-global argument}
Grothendieck needs fibrewise equivalence to imply total equivalence without requiring a map of fibrations to preserve cartesian arrows. Smoothness or properness provides the comparison between strict fibres and comma categories that makes this possible. The treatment of these notions in the commentary on Chapter~I explains the lifting categories behind those comparisons; here their role is descent for a map of total categories.

For a proper functor p : X \ensuremath{\to}S, the canonical functor Xs \ensuremath{\to}X/s is a W\ensuremath{\infty}-equivalence; for a smooth functor, the corresponding functor is Xs \ensuremath{\to}s\textbackslash{}X. These assertions, in the appropriate aspherical form, are the criterion developed in Maltsiniotis (2005, Theorem 1.2 and its dual). Suppose X, Y are proper over S and f : X \ensuremath{\to}Y is any functor over S. Naturality gives

Xs

Ys

X/s

Y/s. Under W\ensuremath{\infty}\ensuremath{\subseteq}W, the vertical maps belong to W. The two-out-of-three property implies that fs \ensuremath{\in}W if and only if f/s \ensuremath{\in}W. If this holds for every s, relative L4 bis gives f \ensuremath{\in}W. For smooth functors the proof uses the undercategories and the dual relative axiom. No cartesian or cocartesian condition on f occurs.

The argument also explains the eventual universal statement on pages 41--43. Proper and smooth functors are stable under change of base. For S\ensuremath{\prime} \ensuremath{\to}S, the strict fibre of the pulled-back map at s\ensuremath{\prime} is fs\ensuremath{\prime} = fp(s\ensuremath{\prime}). Applying the same proof over S\ensuremath{\prime} shows that a fibrewise equivalence between two proper objects, or between two smooth objects, remains a total equivalence after every base change.

The parenthetical assertion here that L3 quater and L4 suffice suppresses one hypothesis: the version of Theorem A actually used is relative L4 bis. Absolute L4 directly covers the case Y = S, since S/s has a terminal object. For general Y it is the stronger formulation that is needed in the displayed proof. Grothendieck explicitly identifies this distinction on manuscript page 44. We follow that correction rather than assigning an unproved consequence to the initial axiom.

\end{dcomment}


On va calquer la construction canonique de $\cat{Ch}(X) \longrightarrow X \times X$ (ou $\cat{Ch}_\infty(X) \longrightarrow X \times X$) factorisant $X \xrightarrow{\mathrm{diag}} X \times X$, et la factorisation canonique correspondante d'un $f : X' \longrightarrow X$ en $X' \xrightarrow{i} \bar{X}' \xrightarrow{\bar{f}} Y$ \texttt{[plutôt} $\bar{X}' \xrightarrow{\bar{f}} X$\texttt{]}. On remplacera simplement $\cat{Ch}(X)$ par la sous-catégorie $\cat{Ch}_+(X)$, laquelle n'est autre que la catégorie des flèches de $X$
\[
\cat{Fl}(X) = \cat{Hom}(\boldsymbol{\Delta}_1, X),
\]
muni des foncteurs
\[
X \xrightarrow{i_X} \cat{Fl}(X) \xrightarrow{p_X=(\sigma_X,\beta_X)} X \times X,
\]
où $i_X$ est la flèche $x \longmapsto \mathrm{id}_x$, et $\sigma_X$, $\beta_X$ sont les foncteurs source et but. La fibre de $p_X$ en $(a,b) \in \mathrm{Ob}(X \times X)$ est la catégorie discrète définie par l'ensemble $\mathrm{Hom}(a,b)$, ça dépend grandement de $(a,b)$ et pas seulement de sa composante connexe dans $X \times X$ - donc $p_X$ n'a rien d'une Hot-fibration. Il n'est en géneral ni propre ni lisse ni rien, par contre les deux flèches composées avec $\mathrm{pr}_1$ et $\mathrm{pr}_2$ % [sic] géneral -> général, missing accent
\[
\begin{tikzcd}[column sep=small]
& \cat{Fl}(X) \arrow[ld, "\sigma_X"'] \arrow[rd, "\beta_X"] & \\
X & & X
\end{tikzcd}
\]
sont sympathiques, on a en effet

\msp{page 17}
\textbf{Lemme.} \textit{On a}
\[
\sigma_X^{-1}(\{x\}) \simeq x\backslash X, \qquad\qquad \beta_X^{-1}(\{x\}) \simeq X/x.
\]
\textit{Le foncteur} $\sigma_X$ \textit{est} Cat-\textit{fibrant à fibres contractiles (donc asphériques, donc} $W$\textit{-asphériques), le foncteur} $\beta_X$ \textit{cofibrant à fibres contractiles (donc asphériques, donc} $W$\textit{-asphériques). Ce sont donc aussi des} Hot-\textit{équivalences, a fortiori des} $W$\textit{-équivalences} (\footnote[8]{En fait, ce sont même des homotopismes, cf. remarque marginale plus bas (plus générale - on retrouve le cas de $\sigma_X$ en faisant $X' = X$.)}).

\textbf{Corollaire 1.} \textit{Le foncteur} $i_X : X \longrightarrow \cat{Fl}(X)$ \textit{est une} Hot-\textit{équivalence} (puisque son composé avec la Hot-équivalence $\sigma_X$ est l'identité).
% foot: Version 9 novembre 2001

\printedpage{11}
% head: A. Grothendieck : Dérivateurs, Ch. VI, Édition des Universités de Montpellier II et Paris VII
Voici la construction générale de la factorisation de $f : X' \longrightarrow X$, correspondant au choix précédent d'une `catégorie de chemins' $\mathcal{C}(X) = \cat{Fl}(X)$
\[
\begin{tikzcd}
\bar{X}' \arrow[d, "\pi_f"] \arrow[rd] \arrow[rrd, "\bar{f}"] & & \\
X' \arrow[u, dotted, bend left=60, "i_f"] \arrow[rd, "f"'] \arrow[r, phantom, "\text{\footnotesize Cart.}"] & \cat{Fl}(X) \arrow[r, "\beta_X"'] \arrow[d, "\sigma_X"'] & X \\
& X &
\end{tikzcd}
\]
où $i_f$ est déduit de la section $i$ \texttt{[plutôt} $i_X$\texttt{]} de $\cat{Fl}(X)$ sur $X$ (\footnote[9]{\textbf{NB} $\pi_f i_f = \mathrm{id}$, et $i_f\pi_f \longrightarrow \mathrm{id}$, donc $i_f$, $\pi_f$ sont des homotopismes inverses l'un de l'autre dans Cat\,hot. Cela précise \texttt{[le]} cor.~3, page 18.}). Ainsi, $\bar{X}'$ est la catégorie des couples
\[
\{(x',u) \mid x' \in \mathrm{Ob}\,X',\ \underbrace{u}_{\in\mathrm{Fl}\,(X)} : f(x') \longrightarrow y\},
\]
le foncteur $\bar{f}$ est défini par
\[
\bar{f}(x',u) = \mathrm{but}(u) = y
\]
et $i_f$ par
\[
i_f(x') = (x', \mathrm{id}_{f(x')} : f(x') \xrightarrow{\mathrm{id}} f(x')),
\]
donc on a bien
\[
f(x') = \bar{f} i_f(x').
\]

\msp{page 18}
\textbf{Corollaire 2.} \textit{Le foncteur}
\[
\bar{f} : \bar{X}' \longrightarrow X
\]
\textit{a comme fibres les catégories $X'/x$, il est cofibrant, les foncteurs de changement de base sont les foncteurs $X'/x \longrightarrow X'/y$ associés aux flèches $x \longrightarrow y$ dans $X$.}

\dgrblock{dgrb201}{The arrow category has two different good endpoints}
\begin{dcomment}{exposition}{p17}{The arrow category has two different good endpoints}
Grothendieck introduces a path construction which remembers the direction of an arrow. Its endpoint maps are excellent replacement tools separately, although the pair of endpoints is not a homotopy-fibration model of the diagonal. The elementary adjunctions make all of these assertions transparent. Write Ar(A) = Fun([1], A), and let i(a) = ida, \ensuremath{\sigma}(a \ensuremath{\to}b) = a, \ensuremath{\beta}(a \ensuremath{\to}b) = b. A commutative square from i(x) to a \ensuremath{\to}b is determined by its arrow x \ensuremath{\to}a. A square from a \ensuremath{\to}b to i(y) is determined by b \ensuremath{\to}y. Therefore

\ensuremath{\beta} \ensuremath{\dashv}i \ensuremath{\dashv}\ensuremath{\sigma}.

In particular \ensuremath{\sigma}i = \ensuremath{\beta}i = id, and the squares associated with an arrow a \ensuremath{\to}b give natural transformations i\ensuremath{\sigma} \ensuremath{\Rightarrow}id and id \ensuremath{\Rightarrow}i\ensuremath{\beta}. Both endpoint projections and i are categorical homotopy equivalences, hence belong to every W satisfying L1 and L2.

The fibration properties can be checked without choosing anything. Pulling a \ensuremath{\to}b back along a\ensuremath{\prime} \ensuremath{\to}a means precomposing to obtain a\ensuremath{\prime} \ensuremath{\to}b; the resulting square is cartesian for \ensuremath{\sigma}. Pushing it forward along b \ensuremath{\to}b\ensuremath{\prime} means postcomposing to obtain a \ensuremath{\to}b\ensuremath{\prime}; that square is cocartesian for \ensuremath{\beta}. The fibres are a\textbackslash{}A and A/b, with initial and terminal identities respectively. This example already belongs to the original descent formalism; see Grothendieck (1971, Exposé VI, \S{}11(a)).

For the combined map (\ensuremath{\sigma}, \ensuremath{\beta}) : Ar(A) \ensuremath{\to}A \ensuremath{\times} A, however, the fibre at (a, b) is the discrete set HomA(a, b). Take A = [1]. Its fibre at (1, 0) is empty, while that at (0, 0) is a point, although the base is connected. Thus it is not a Hot-fibration. It is neither smooth nor proper: at (1, 0) its undercategory has the object arising from (1, 0) \ensuremath{\to}(1, 1), and its overcategory has the object arising from (0, 0) \ensuremath{\to}(1, 0), whereas its strict fibre is empty. In fact each of these comma categories is a point. The necessary fibre-to-comma weak-equivalence condition fails in both directions. This distinction is why the directed arrow construction gives \ensuremath{\Phi} and \ensuremath{\Psi}, while the two-sided replacement C requires the longer path category Ch\ensuremath{\infty}(S).

\end{dcomment}


En fait, on voit que cette catégorie est non seulement cofibrée, mais canoniquement \textit{scindée.}

(\footnote[10]{\textbf{Corollaire 3.} \textit{La flèche $i_f : X' \longrightarrow \bar{X}'$ est un $X$-morphisme (si $X'$ sur $X$ par $f$, $\bar{X}'$ par $\bar{f}$), i.e. $f = \bar{f}\circ i_f$, et c'est une} Hot-\textit{équivalence} (car $\pi_f : \bar{X}' \longrightarrow X'$ est Cat-fibrant, à fibres contractiles, donc une Hot-équivalence).}).

Appliquant ceci à des $f : X \longrightarrow S$, on trouve une construction
\[
\mathrm{Cat}/S \xrightarrow{\Psi} \cat{Cofibsc}\,S,
\]
c'est bel et bien fonctoriel en $X$ variant dans $\mathrm{Cat}/S$. Si
\[
\psi : X \longrightarrow Y
\]
% foot: Version 9 novembre 2001
\printedpage{12}
% head: A. Grothendieck : Dérivateurs, Ch. VI, Édition des Universités de Montpellier II et Paris VII
dans $\mathrm{Cat}/S$, alors $\psi$ induit des
\[
\psi/s : X/s \longrightarrow Y/s\ ,
\]
pour $s \in S$, et ceci de façon compatible avec les
\msp{page 19}
foncteurs `cochangement de base' associés aux $u : s \longrightarrow t$ dans $S$
\[
\begin{tikzcd}[column sep=large, row sep=large]
X/s \arrow[r, "u_*^X"] \arrow[d, "\psi/s"'] & X/t \arrow[d, "\psi/t"] \\
Y/s \arrow[r, "u_*^Y"] & Y/t \quad .
\end{tikzcd}
\]
Cela définit donc un morphisme de catégories scindées sur $S$
\[
\Psi(X) \longrightarrow \Psi(Y),
\]
c'est lui!

On a un morphisme fonctoriel en $X$
\[
X \xrightarrow{i_X} \Psi(X)
\]
dans $\mathrm{Cat}/X$, et c'est une Hot-équivalence, donc $\in W$. Sur les fibres, le morphisme $i_X$ induit
\[
X_s \xrightarrow{i_{X,s}} X/s,
\]
l'inclusion canonique.

\msp{page 20}
\textbf{Corollaire 1.} \textit{Si $X \xrightarrow{f} Y$ dans} $\mathrm{Cat}/S$ \textit{est dans $W$, alors $\Psi(X) \xrightarrow{\Psi(f)} \Psi(Y)$ aussi.} Mais on se gardera de croire que c'est dans $W_S$ au sens de $\cat{Cofib}\,S$ ou $\cat{Cofibsc}\,S$.

\dgrblock{dgrb202}{A canonical opfibration and its two projections}
\begin{dcomment}{exposition}{p19}{A canonical opfibration and its two projections}
Grothendieck’s replacement \ensuremath{\Psi} packages all the overcategories X/s into one category. Writing down its arrows shows both why it is strictly functorial and why a projection used to prove homotopy invariance need not be a morphism over the original base.

For p : X \ensuremath{\to}S, an object of \ensuremath{\Psi}(X) is (x, u : p(x) \ensuremath{\to}s). A morphism to (y, u\ensuremath{\prime} : p(y) \ensuremath{\to}t) is a pair

vu = u\ensuremath{\prime}p(a).

(a : x \ensuremath{\to}y, v : s \ensuremath{\to}t), The structural projection q sends this object to s and this morphism to v. For v : s \ensuremath{\to}t, the morphism (idx, v) : (x, u) \ensuremath{\to}(x, vu) is the specified cocartesian lift. Composition of these lifts is literal composition, so the opfibration is split and its fibre at s is X/s.

There is also the forgetful projection \ensuremath{\pi} : \ensuremath{\Psi}(X) \ensuremath{\to}X, which sends (x, u) to x, and the functor i : X \ensuremath{\to}\ensuremath{\Psi}(X) sending x to (x, idp(x)). A map from i(x) to (y, u\ensuremath{\prime}) is uniquely determined by a : x \ensuremath{\to}y, because its base arrow must be u\ensuremath{\prime}p(a). Hence i \ensuremath{\dashv}\ensuremath{\pi}, with \ensuremath{\pi}i = idX and counit

\ensuremath{-}\ensuremath{-}\ensuremath{-}\ensuremath{-}\ensuremath{\to}(x, u).

i\ensuremath{\pi}(x, u) = (x, idp(x)) This proves that i and \ensuremath{\pi} are inverse up to categorical homotopy. The identity qi = p makes i a functor over S, but p\ensuremath{\pi}(x, u) = p(x) generally differs from q(x, u) = s. Thus \ensuremath{\pi} generally is not a functor over S. The printed statement that i is a morphism in Cat/X is nonetheless true when \ensuremath{\Psi}(X) carries \ensuremath{\pi}: i is a section. Both descriptions of i are legitimate, and neither licenses using \ensuremath{\pi} as an inverse in Cat/S.

For a functor f : X \ensuremath{\to}Y over S, define \ensuremath{\Psi}(f)(x, u) = (f(x), u) and \ensuremath{\Psi}(f)(a, v) = (f(a), v). This preserves the specified cocartesian lifts and is strictly functorial in f. The maps i form a natural transformation in Cat/S. Their being in W shows, by the resulting commutative square, that \ensuremath{\Psi} preserves and reflects coarse weak equivalences. On fibres, however, its effect is f/s. This last observation specifies the stronger class for which \ensuremath{\Psi} becomes a replacement preserving fibrewise weak equivalences.

\end{dcomment}


Car sur les fibres de $\Psi(X)$, $\Psi(Y)$ on trouve les flèches
\[
f/s : X/s \longrightarrow Y/s
\]
induits par $f$, qui n'ont aucune raison d'être dans $W$. Qu'elles le soient signifie que $f : X \longrightarrow Y$ est dans $W$ `localement au dessus de S' (\footnote[11]{cf. \circled{1} page 15.}), ce qui (pour les $W$ qui nous intéressent) est beaucoup plus fort que d'être dans $W$. (Si $Y = S$, cela signifie en fait que $p : X \longrightarrow S$ est $W$-asphérique, beaucoup plus fort que $p \in W$.)

Donc pour prouver que
\[
\cat{Cofib}\,S \longrightarrow \mathrm{Cat}/S \qquad\quad \text{ou} \qquad\quad \cat{Cofibsc}\,S \longrightarrow \mathrm{Cat}/S
\]
% foot: Version 9 novembre 2001
\printedpage{13}
% head: A. Grothendieck : Dérivateurs, Ch. VI, Édition des Universités de Montpellier II et Paris VII
induisent des équivalences pour les catégories localisées, le foncteur $\Psi$ directement ne sera pas utilisable. Cependant :

\textbf{Corollaire 2.} \textit{Supposons $X$, $Y$ dans} $\cat{Prop}\,S$ \textit{ou dans} $\cat{Fibhot}\,S$\textit{, et considérons}
\[
\Psi(f) : \Psi(X) \longrightarrow \Psi(Y),
\]
\textit{pour que ce soit une $W$-équivalence dans} $\cat{Cofibsc}\,S$\textit{,}
\msp{page 21}
\textit{i.e. une $W$-équivalence fibre par fibre, il faut et il suffit que $f$ soit une $W$-équivalence fibre par fibre} (et cela doit impliquer (\footnote[12]{cf. \circled{1} page 15.}) que $f$ est une $W$-équivalence).

\dgrblock{dgrb203}{Which comma construction repairs which variance}
\begin{dcomment}{exposition}{p20}{Which comma construction repairs which variance}
Grothendieck now isolates the extra condition that makes \ensuremath{\Psi} a relative replacement for fibrewise homotopy theory. Its fibre map is exactly f/s, so the condition is that every overcategory map lie in W. When the target is S itself this says that f is W-aspherical, since S/s has a terminal object. For a general target it is a relative condition, later called fine equivalence on the left.

The interval calculation after manuscript page 8 already shows why total equivalence cannot suffice. The complete comparison of both endpoint inclusions, given alongside manuscript page 32, also separates the two comma tests. On a smooth object, the strict fibre compares to the undercategory s\textbackslash{}X; on a proper object, to the overcategory X/s. Consequently the relevant fibrewise replacement for a fibration is \ensuremath{\Phi}, and for an opfibration it is \ensuremath{\Psi}. This is the geometric reason the later fine-right equivalence belongs on the fibration side of the large diagram, while fine-left equivalence belongs on the opfibration side. The direction names refer to the comma constructions, not to the position of the objects in that diagram.

\end{dcomment}


En effet, dans le diagramme commutatif
\[
\begin{tikzcd}[column sep=large, row sep=huge]
X_s \arrow[r, "f_s"] \arrow[d] & Y_s \arrow[d] \\
X/s \arrow[r, "f/s"] & Y/s
\end{tikzcd}
\]
les flèches verticales sont dans $W$, car $X$, $Y$ \texttt{[sont]} propres sur $S$ (ou $X$, $Y$ dans $\cat{Fibhot}\,S$), donc $f_s$ est dans $W$ si et seulement si $f/s$ l'est.

\texttt{[Le foncteur]} $\Phi$ se définit de façon duale à $\Psi$, ce qui revient à échanger les rôles de $\sigma_X$, $\beta_X$ dans le diagramme page 17.

(\footnote[13]{\textbf{Corollaire 3.} \textit{Si $X \in \mathit{Ob}\,$}$\cat{Fibhot}\,S$\textit{, alors $\Psi(X)$ est dans} $\cat{Fibsc}_0S$.\par Cela signifie en effet que les $X/s \longrightarrow X/t$ sont dans $W$, or l'hypothèse sur $X$ implique que c'est dans $W_\infty$. OK.})

\textbf{II) Définition de} $C : \mathrm{Cat}/S \longrightarrow \cat{Parf}\,S$

Cette fois, c'est la factorisation des flèches $p : X \longrightarrow S$ dans Cat, en
\[
X \xrightarrow{j_{X/S}} \bar{X} \xrightarrow{\bar{p}} S,
\]
où $i_X$ \texttt{[plutôt} $j_{X/S}$\texttt{]} est une Hot-équivalence (donc dans $W$), et $\bar{p}$ parfait. Cela se construit formellement comme tantôt, en utilisant $\cat{Ch}_\infty(S)$ au lieu de $\cat{Fl}(S)$, et la factorisation
\[
S \xrightarrow{j_S} \cat{Ch}_\infty(S) \xrightarrow{p_S=(\sigma_S,\beta_S)} S \times S
\]
de $\mathrm{diag}_S : S \longrightarrow S \times S$. La situation est plus jolie, car $p_S$ est déjà `parfait', ce qui implique formellement que le $\bar{p}$ qui s'en déduit, pour factoriser $p : X \longrightarrow S$, est lui-même parfait.
% foot: Version 9 novembre 2001

% ---- from Der-ChVI-345_p016-020.tex
\printedpage{14}
% head: A. Grothendieck : Dérivateurs, Ch. VI, Édition des Universités de Montpellier II et Paris VII
% foot: Version 9 novembre 2001

\msp{page 22}
Ici encore, la construction est fonctorielle, on trouve le foncteur
\[ C : \Cat/S \longrightarrow \und{Parf}\,S \]
qu'il faut composer au foncteur en sens inverse
\[ r : \Cat/S \longleftarrow \und{Parf}\,S. \]
Tout d'abord, vérifions que $C$, tout comme $r$, est compatible avec les localiseurs. Soit $f : X \longrightarrow Y$ dans $\Cat/S$, supposons $f$ est dans $W$, prouvons $C(f)$ est dans $W_S$. On a commutativité dans
\[
\begin{tikzcd}[row sep=large, column sep=large]
X \arrow[r, "f"] \arrow[d, "j_X"'] & Y \arrow[d, "j_Y"] \\
CX \arrow[r, "Cf"] & CY
\end{tikzcd}
\]
et $j_X$, $j_Y$ sont dans $W$, donc $f$ est dans $W$ si et seulement si $Cf$ l'est. Mais dans le cas $W = W_\infty$ on sait que cela équivaut à ce que $Cf$ soit dans $W_\infty$ fibre par fibre. Admettant \circled{1} p. 15 pour $W$, il reste à s'assurer d'une réciproque dans \underline{Parf}, i.e.

\dgrblock{dgrb204}{From comma fibres to homotopy fibres}
\begin{dcomment}{later}{p21}{From comma fibres to homotopy fibres}
Grothendieck introduces the longer path category because he now wants a replacement whose transport is homotopically invertible in both directions. The guiding theorem is Quillen’s Theorem B: if every transition X/s \ensuremath{\to}X/t induced by an arrow s \ensuremath{\to}t is a weak homotopy equivalence, then N(X/s) computes the homotopy fibre of NX \ensuremath{\to}NS over s. A precise restatement appears in Cisinski (2004, Theorem 2.3.2). This is the point at which locally constant comma data become actual fibre data for the homotopy type of a map.

For an object already Hot-fibred over S, the comparison Xs \ensuremath{\to}X/s is a weak equivalence, and the transitions between these comma categories are weak equivalences. Consequently \ensuremath{\Psi}(X) is a split opfibration with homotopically invertible transition functors. The footnote’s printed target Fibsc0S is therefore a slip: the target is Cofibsc0S. Its dual statement concerns \ensuremath{\Phi} and split fibrations. The slip does not affect the argument in the adjacent square.

The functor C uses results about Ch\ensuremath{\infty}(S) established outside these pages. The input here is a natural factorization X \ensuremath{-}\ensuremath{-}\ensuremath{\to}C(X) \ensuremath{\to}S with jX \ensuremath{\in}W\ensuremath{\infty}and C(X) \ensuremath{\to}S perfect. Given this input, all subsequent uses of its functoriality are formal. For instance, the naturality square of j implies that f \ensuremath{\in}W if and only if C(f) \ensuremath{\in}W, provided W\ensuremath{\infty}\ensuremath{\subseteq}W. What is not formal is replacing the last total condition by a condition on fibres for an arbitrary larger W. That is the question on pages 22--24, and the stronger L5 introduced later addresses exactly it. The correct comparison therefore depends on which localizer is being studied, not only on the quality of the replacement object.

\end{dcomment}


\circled{2} Soit
\begin{tikzcd}[baseline=(\tikzcdmatrixname-2-2.base), row sep=small, column sep=small]
X \arrow[rr, "f"] \arrow[dr] & & Y \arrow[dl] \\
& S &
\end{tikzcd}
dans $\und{Parf}\,S$ (ou seulement dans $\und{Fibhot}\,S$). Si $f$ est dans $W$, alors les $f_s : X_s \longrightarrow Y_s$ aussi.

\msp{page 23}
Essayons de le prouver, en calquant la démonstration dans le cas $C = C_\infty$ (XI, th.~2). On peut supposer $S$ connexe (\footnote[14]{ne sert à rien}), en utilisant La (i.e. $W \subseteq W_0$) et en admettant le

\textbf{Lemme.} \textit{Soit $f : X \longrightarrow Y$ dans} Cat. \textit{Alors $f \in W$ si et seulement si $\pi_0(f)$ est bijectif, et pour toute composante connexe $X_i$ de $X$, si $Y_i$ est la composante connexe correspondante de $Y$, $f_i : X_i \longrightarrow Y_i$ induit par $f$ est dans $W$.}

Utilisons alors l'existence de
\[ S' \longrightarrow S \]
Hot-fibration avec $S'$ Hot-asphérique (\footnote[15]{\textbf{NB} Il suffit d'avoir une Hot-fibration $S' \longrightarrow S_0 \hookrightarrow S$, où $S_0$ est la composante connexe de $S$ contenant l'élément $s \in \mathrm{Ob}\,S$ envisagé.}). On peut supposer même que $S' \longrightarrow S$ est parfait. Supposons qu'on sache qu'un changement de base par de telles fibrations parfaites transforme $W$-équivalences en itou. (Ça doit être vrai pour les $W_i$, mais on a vu que c'était faux pour l'équivalence $k$-homologique, cf. I p.~87 ff.) Alors $X' \xrightarrow{f'} Y'$ sur $S'$ est
\printedpage{15}
% head: A. Grothendieck : Dérivateurs, Ch. VI, Édition des Universités de Montpellier II et Paris VII
% foot: Version 9 novembre 2001
une $W$-équivalence. Or pour $s' \in S'$, comme $\{s'\} \longrightarrow S'$ est une Hot-équivalence ($S'$ Hot-asphérique), et $X'$, $Y'$ étant Hot-fibrés sur $S$, il s'ensuit que $X'_{s'} \longrightarrow X'$, $Y'_{s'} \longrightarrow Y'$ sont des Hot-équivalences, donc aussi $X'_{s'} \xrightarrow{f'_{s'}} Y'_{s'}$, comme on voit sur le carré commutatif
\[
\begin{tikzcd}[row sep=large, column sep=large]
X'_{s'} \arrow[r, "f'_{s'}"] \arrow[d, "\wr"'] & Y'_{s'} \arrow[d, "\wr"'] \\
X' \arrow[r, "{f'}"] & Y'.
\end{tikzcd}
\]

\msp{page 24}
Comme $S' \longrightarrow S$ est surjectif, cela montre que les $f_s : X_s \longrightarrow Y_s$ sont dans $W$, qed.

\dgrblock{dgrb205}{The contractible-base argument without a component reduction}
\begin{dcomment}{exposition}{p23}{The contractible-base argument without a component reduction}
Grothendieck’s strategy is to inspect a chosen fibre after making the base contractible. Its force is easiest to see when the required base-change assumption is stated before the reduction. It then becomes clear that the componentwise lemma is unnecessary for this particular argument, as his marginal annotation already observes.

Assume W\ensuremath{\infty}\ensuremath{\subseteq}W, two-out-of-three, and the following additional property: pullback along a perfect functor preserves W-equivalences. Let f : X \ensuremath{\to}Y be a functor over S, with X, Y Hot-fibred over S, and suppose its total functor lies in W. Fix s \ensuremath{\in}S. Use the perfect replacement of the point \{s\} \ensuremath{\to}S to obtain

\ensuremath{-}\ensuremath{\to}jS\ensuremath{\prime} \ensuremath{-}\ensuremath{\to}rS, \{s\} j \ensuremath{\in}W\ensuremath{\infty}, where r is perfect and S\ensuremath{\prime} is weakly contractible. The distinguished object s\ensuremath{\prime} = j(s) lies over s. This construction is allowed even when S is disconnected; its image lies in the relevant component.

Put X\ensuremath{\prime} = X \ensuremath{\times}S S\ensuremath{\prime} and Y \ensuremath{\prime} = Y \ensuremath{\times}S S\ensuremath{\prime}. The extra pullback hypothesis gives f \ensuremath{\prime} : X\ensuremath{\prime} \ensuremath{\to}Y \ensuremath{\prime} in W. Base change preserves the Hot-fibration property. Since S\ensuremath{\prime} is weakly contractible, the inclusions of the fibres at s\ensuremath{\prime} into X\ensuremath{\prime} and Y \ensuremath{\prime} are weak homotopy equivalences. Equivalently, their homotopy-fibre squares are pulled back along the weak equivalence \{s\ensuremath{\prime}\} \ensuremath{\to}S\ensuremath{\prime}. In the commutative square

Xs

Ys

X\ensuremath{\prime}

Y \ensuremath{\prime}, three sides lie in W, so fs \ensuremath{\in}W. This proves the fibre-detection assertion under precisely the stated additional hypothesis. It does not prove fs \ensuremath{\in}W\ensuremath{\infty}from f \ensuremath{\prime} \ensuremath{\in}W.

The proof uses neither La nor a decomposition of f into connected components. Its genuinely substantial input is stability of W under these perfect base changes, later called L6. The speculation here that truncated homotopy equivalences should satisfy that stability is withdrawn in the later annotations; the obstruction is discussed at manuscript pages 43--45. The proof itself is valid as a conditional argument, and its dependence on L6 is the mathematical information to retain.

\end{dcomment}


Ainsi, (2) sera conséquence de \circled{2.} Soit $p : Y' \longrightarrow Y$ une Hot-fibration (ou même un foncteur parfait, si ça peut aider \dots). Soit $f : X \longrightarrow Y$ une $S$-flèche (dans $\Cat/S$), alors $f' : X' \longrightarrow Y'$ par changement de base. Si $f \in W$, $f' \in W'$
\[
\begin{tikzcd}[row sep=large]
X' \arrow[r, "q"] \arrow[d, "{f'}"'] \arrow[dr, phantom, "\scriptstyle\text{cart.}"] & X \arrow[d, "f"] \\
Y' \arrow[r, "p"] & Y
\end{tikzcd}
\qquad\qquad
f \in W \text{ et } p \text{ parfait} \Longrightarrow f' \in W
\]

En tous cas, il nous faut absolument \circled{2} pour pouvoir conclure que $C$ est compatible avec les localiseurs.

Considérons $rC : \Cat/S \longrightarrow \Cat/S$, on a une flèche
\[
\left\{
\begin{array}{l}
\mathrm{id}_{\Cat/S} \xrightarrow{j} rC \\[1.5ex]
j_X : X \longrightarrow CX \quad \text{la flèche canonique } i_{X/S}\ \text{\texttt{[plutôt} } j_{X/S}\text{\texttt{]}} \text{ (p.~21),}
\end{array}
\right.
\]
et $j_X$ est toujours dans $W$.

Considérons $Cr : \und{Parf}\,S \longrightarrow \und{Parf}\,S$, i.e. appliquons la construction $C$ à un $X$ parfait sur $S$. Alors on a encore
\[
\left\{
\begin{array}{l}
\mathrm{id}_{\und{Parf}\,S} \xrightarrow{j'} Cr \\[1ex]
j'_X = j_{r(X)} : X = r(X) \xrightarrow{j_X} CX.
\end{array}
\right.
\]
Ce morphisme est dans $W$, donc aussi une $W$-équivalence fibre par fibre, puisque source et but sont parfaits, donc des Hot-fibrés sur $S$.

\msp{page 25}
Cela prouve déjà que les foncteurs $r$ et $C$ induisent des équivalences quasi-inverses l'une de l'autre sur les catégories localisées. Par le lemme p.~11, cela montre que pour toute sous-catégorie pleine $\mathcal{F}$ de $\Cat/S$ contenant $\und{Parf}\,S$, l'inclusion de $\mathcal{F}$ dans $\Cat/S$ induit une
\printedpage{16}
% head: A. Grothendieck : Dérivateurs, Ch. VI, Édition des Universités de Montpellier II et Paris VII
% foot: Version 9 novembre 2001
équivalence sur les localisés, en prenant comme localiseur dans $\mathcal{F}$ le localiseur `grossier' induit par $\Cat/S$. Cela s'applique notamment aux sous-catégories pleines
\[ \und{Liss}\,S, \quad \und{Fibhot}\,S, \quad \und{Prop}\,S, \]
mais en prenant dans $\und{Liss}\,S$ et $\und{Prop}\,S$ non pas le localiseur fin `fibre par fibre', comme je l'avais annoncé à tort, mais le localiseur grossier induit par $\Cat/S$. (Pour $\und{Fibhot}\,S$, il n'y a pas de risque d'ambiguïté). Il faudrait que je me convainque que si je prends la $W$-équivalence fine (dans le cas $W = W_\infty$, disons), je ne trouve \textit{pas} les `bons' localisés, i.e. le foncteur des localisés, vers le localisé de $\Cat/S$, n'est pas une équivalence. Ceci en fait suppose que dans $\und{Fib}\,S$ et $\und{Cofib}\,S$ également, c'est l'équivalence \textit{grossière}
\msp{page 26}
qu'il convient de prendre. (Dans $\und{Fib}_0S$ et $\und{Cofib}_0S$, ou $\und{Fibsc}_0S$ et $\und{Cofibsc}_0S$, les deux équivalences sont identiques, puisque ces quatre catégories sont contenues dans $\und{Fibhot}\,S$.) La suite nous montrera s'il en est bien ainsi. Par contre, dans $\und{Fibsc}\,S$ et $\und{Cofibsc}\,S$, il n'y a pas de doute, suivant les principes généraux de construction d'un dérivateur à partir d'un $(\mathcal{M}, W)$, qu'il faut prendre l'équivalence fine, fibre par fibre. À ce sujet, il faut faire attention que les foncteurs verticaux du diagramme p.~8
\[ \und{Fibsc}\,S \longrightarrow \und{Fib}\,S, \quad \und{Cofibsc}\,S \longrightarrow \und{Cofib}\,S, \]
quoique essentiellement surjectifs et fidèles, ne sont pas pleinement fidèles. Ils le seraient, et mêmes des équivalences de catégories, si dans $\und{Fib}\,S$ resp. $\und{Cofib}\,S$ je me bornais aux $S$-foncteurs \textit{cartésiens} resp. \textit{cocartésiens}. Donc on peut interpréter, à équivalence près, $\und{Fibsc}\,S$ et $\und{Cofibsc}\,S$ comme $\und{Fib}\,S$ et $\und{Cofib}\,S$, où on aurait restreint les morphismes aux seuls $S$-foncteurs \textit{cartésiens} resp. \textit{cocartésiens.} Il semblerait donc qu'en restreignant ainsi les flèches, mais en restreignant de façon plus draconienne encore le localiseur, on trouve des catégories localisées équivalentes. % [sic] et mêmes des équivalences -> et même des équivalences, adverb "même" printed with plural s

\dgrblock{dgrb206}{What the path replacement actually proves}
\begin{dcomment}{exposition}{p24}{What the path replacement actually proves}
Grothendieck wants a category over S whose strict fibres already have the homotopy types of the homotopy fibres. The replacement jX : X \ensuremath{-}\ensuremath{\to}C(X) achieves this geometrically: jX is a weak homotopy equivalence, is natural over S, and C(X) \ensuremath{-}\ensuremath{\to}S is perfect. The difficulty is to decide which weak equivalences this replacement preserves. Two assertions must be separated, because one is formal and the other detects a substantive property of the coefficient theory.

Write W tot for maps over S whose underlying functor lies in W. Assuming W\ensuremath{\infty}\ensuremath{\subseteq}W and two-outof-three, naturality gives

f \ensuremath{\in}W tot

\ensuremath{\Leftarrow}\ensuremath{\Rightarrow}C(f) \ensuremath{\in}W tot C(f)jX = jY f, S . Thus no base-change axiom is needed to make C a functor between the coarse relative categories. If F \ensuremath{\subseteq} Cat/S is full and contains every perfect object, C factors through F. The two natural transformations j on Cat/S and its restriction to F become invertible after coarse localization. Consequently the inclusion of F induces an equivalence. This proves the assertion of manuscript page 25 under the stated elementary hypotheses, independently of the conjectural pullback assertion on the present page.

The stronger assertion is that C(f) is a W-equivalence on every fibre whenever f \ensuremath{\in}W tot S . This does require the implication later named L5 on manuscript page 43. The preceding argument with a contractible covering base explains a sufficient condition precisely. For a chosen s \ensuremath{\in}S, take a perfect S\ensuremath{\prime} \ensuremath{-}\ensuremath{\to}S with contractible total category and an object over s; it suffices to cover the component containing s. If pullback along this map preserves W, the resulting total map X\ensuremath{\prime} \ensuremath{-}\ensuremath{\to}Y \ensuremath{\prime} belongs to W. For homotopy fibrations the inclusions of the chosen fibres into X\ensuremath{\prime} and Y \ensuremath{\prime} are weak homotopy equivalences. Two-out-of-three then gives the desired fibre map in W. The final comparison gives a Wequivalence, not necessarily a weak homotopy equivalence: the sentence immediately before manuscript page 24 compresses these two conclusions.

For W = W\ensuremath{\infty}this is the familiar invariance of a homotopy fibre under equivalence of total spaces over a fixed base. For a coarser W it is an additional demand. The covering calculation at manuscript page 43 below will show its failure even when both maps to the base are discrete fibrations and discrete opfibrations. In particular the phrase “il nous faut absolument” concerns compatibility with the fibrewise localizer, not the already established compatibility with the coarse one.

\end{dcomment}


\msp{page 27}
C'est assez étrange - on va voir si c'est bien vrai!

\dgrblock{dgrb207}{Strict splittings and Cartesian maps}
\begin{dcomment}{context}{p26}{Strict splittings and Cartesian maps}
The comparison with diagram categories depends on what the morphisms preserve. The definition on manuscript page 4 uses functors preserving the chosen Cartesian lifts strictly. A Cartesian functor between split fibrations need only carry a chosen lift to some Cartesian lift; the latter is isomorphic to the chosen lift at the image object. These conditions give strict and pseudonatural transformations, respectively. They are different conditions on actual functors, before localization.

Here is a concrete failure of fullness. Put S = [1] and let G = C2 = \{1, c\}. The projection S \ensuremath{\times}BG \ensuremath{-}\ensuremath{\to}S has its constant splitting, with chosen lift (u, 1) of u : 0 \ensuremath{-}\ensuremath{\to}1. Define an endofunctor over S which fixes the two fibre groupoids and sends a cross-fibre arrow (u, g) to (u, cg). This is a functor: composition with a vertical arrow is compatible with the formula because G is abelian. Every arrow over u is Cartesian, since its group component is invertible. Our endofunctor is therefore Cartesian, but it does not preserve the splitting, since (u, 1) is sent to (u, c). Restricting the target category to Cartesian functors does not make the forgetful functor from strictly split morphisms full.

There is a separate qualification concerning objects. The functor BC4 \ensuremath{-}\ensuremath{\to}BC2 induced by reduction modulo 2 is a Grothendieck fibration: each base arrow has a lift and every arrow is Cartesian. A splitting would be a homomorphic section C2 \ensuremath{-}\ensuremath{\to}C4. Such a section cannot exist, because each lift of the nonidentity element has order 4. Hence this fibration is not isomorphic over BC2 to an object with a splitting. This is the group-extension example already explained in Grothendieck (1971, Exposé VI, \S{}9). An equivalence over the base to a split replacement is the appropriate statement; essential surjectivity in the ordinary category Cat/S would ask for an isomorphism of categories over the base.

These distinctions do not damage the intended homotopical comparison. Equivalences of fibre categories belong to W\ensuremath{\infty}, and strict replacements can be connected by such equivalences. What needs proof is precisely that replacing coherent choices by strict ones induces the correct localized comparison. Grothendieck encounters that issue directly in the failure of the proposed unit on manuscript page 30 and returns to it on pages 55--60. The underlying classical language comes from the descent formalism of SGA 1; the new task here is to control the weak equivalences during this passage. We retain the printed assertion and read it with this necessary qualification, rather than changing the category of morphisms silently.

\end{dcomment}


\subsection*{III) Construction de $\Phi_0, \Psi_0 :$ \begin{tikzcd}[ampersand replacement=\&, baseline=(\tikzcdmatrixname-2-1.base), row sep=small, column sep=normal]
\& \und{Fibsc}_0S \\
\Cat/S \arrow[ur] \arrow[dr] \& \\
\& \und{Cofibsc}_0S.
\end{tikzcd}}

On doit avoir
\[ \Phi_0(X) = \Phi r C(X) \qquad\qquad\qquad \Psi_0(X) = \Psi r C(X), \]
où
\[ r : \und{Parf}\,S \longrightarrow \Cat/S \]
est l'inclusion canonique. En tant que foncteurs à valeurs dans \underline{Fibsc} resp. \underline{Cofibsc}, cela définit déjà $\Phi_0$, $\Psi_0$ par
\[ \Phi_0 = \Phi r C \qquad\qquad\qquad \Psi_0 = \Psi r C, \]
\printedpage{17}
% head: A. Grothendieck : Dérivateurs, Ch. VI, Édition des Universités de Montpellier II et Paris VII
% foot: Version 9 novembre 2001
mais il faut voir que $\Phi_0$, $\Psi_0$ prennent leurs valeurs dans les sous-catégories pleines à indice 0. Cela équivaut donc au Lemme suivant, appliqué à $Y = rC(X)$:

\textbf{Lemme.} \textit{Soit $Y$ parfait sur $S$ (ou seulement un \textup{Hot}-fibré sur $S$). Alors $\Phi(Y)$ est dans $\und{Fibsc}_0S$, $\Psi(Y)$ dans $\und{Cofibsc}_0S$.}

Il faut prouver (pour $\Psi(Y)$) que si $s \longrightarrow t$ est une flèche dans $S$, alors
\[ \Psi(Y)_s \longrightarrow \Psi(Y)_t \]
est une $W$-équivalence. Or cette flèche s'identifie à
\[ Y/s \longrightarrow Y/t \]
et on sait bien que c'est une Hot-équi\msp{page 28}valence, par définition des Hot-fibrés, donc c'est une $W$-équivalence.

On a donc les foncteurs
\[
\begin{tikzcd}[column sep=large]
\und{Fibsc}_0S \arrow[r, shift left, "p"] & \Cat/S, \arrow[l, shift left, "\Phi_0"]
\end{tikzcd}
\]
il faut voir qu'ils satisfont aux conditions générales du lemme p.~11 (et ce sera pareil pour $q$, $\Psi_0$). On sait déjà que $p$ est compatible avec les localiseurs. Prouvons le pour $\Psi_0$. Soit
\[ f : X \longrightarrow Y \]
dans $\mathrm{Cat/S}$ une $W$-équivalence, alors $rC(f)$ l'est aussi. Il suffit de prouver ceci : % [sic] Cat/S (upright S) -> Cat/$S$, the S printed upright unlike elsewhere

Soit
\[ q : X' \longrightarrow Y' \]
un morphisme dans $\und{Fibhot}\,S$ (par exemple un morphisme de $\und{Parf}\,S$), qui soit $W$-équivalence, donc \texttt{[qui induit]} des $W$-équivalences sur les fibres. Alors $\Psi(q) : \Psi(X') \longrightarrow \Psi(Y')$ est une $W$-équivalence fibre par fibre, i.e.
\[ X'/s \longrightarrow Y'/s \]
une $W$-équivalence pour tout $s \in S$ (a été vu dans le corollaire 2, page 20).

\msp{page 29}
Considérons
\[ p\Phi_0 : \Cat/S \longrightarrow \Cat/S, \]
ça transforme donc $X$ en $\Phi C(X)$, regardé comme objet de $\Cat/S$ sans plus. On a un morphisme fonctoriel
\[ X \xrightarrow{j_X} C(X) \qquad\qquad\qquad j : \mathrm{id}_{\Cat/S} \longrightarrow C, \]
\printedpage{18}
% head: A. Grothendieck : Dérivateurs, Ch. VI, Édition des Universités de Montpellier II et Paris VII
% foot: Version 9 novembre 2001
d'où un morphisme fonctoriel
\[ \Phi(X) \xrightarrow{\Phi(j_X)} \Phi C(X) \qquad\qquad\qquad \Phi * j : \Phi \longrightarrow \Phi_0 \]
et également un morphisme fonctoriel
\[ X \xrightarrow{i_X} \Phi(X), \]
d'où en composant
\[
\begin{tikzcd}
X \arrow[r, "i_X"] \arrow[rr, bend right, "\alpha_X"'] & \Phi(X) \arrow[r, "\Phi(j_X)"] & \Phi C(X),
\end{tikzcd}
\]
i.e.
\[ \alpha : \mathrm{id}_{\Cat/S} \longrightarrow p\Phi_0. \]
Pour $X$ fixé dans $\Cat/S$, on sait que $j_X$ est dans $W$, donc aussi $\Phi(j_X)$ (quoique non fibre par fibre), cf. cor.~1 page 20. On sait aussi que $i_X$ est dans $W$, donc aussi le composé $\alpha_X$.

Considérons
\[ \Phi_0 p : \und{Fibsc}_0\,S \longrightarrow \und{Fibsc}_0\,S\ \ . \]
On part de $X$ dans $\und{Fibsc}_0\,S$, on oublie le scindage et le fait même qu'il soit Cat-fibrant
\msp{page 30}
(on a appliqué $p$), on lui applique $\Phi_0$, i.e. on regarde $\Phi CX$ comme objet de $\und{Fibsc}\,S$ (et même de $\und{Fibsc}_0\,S$). On veut définir encore
\[ X \xrightarrow{\beta_X} \Phi_0 p(X) = \Phi C(X) \qquad\quad \text{\textbf{dans} } \und{Fibsc}_0\text{\textbf{,}} \] % ?? uncertain: weight of "Fibsc_0" after bold "dans" -- regular or bold
fonctoriel en $X$. Or on a comme tantôt
\[ X \xrightarrow{j_X} C(X), \]
d'où en appliquant $\Phi$
\[ \Phi(X) \xrightarrow{\Phi(j_X)} \Phi C(X), \]
mais regardé cette fois comme flèche dans $\und{Fibsc}_0\,S$ (N.B. on a bien $\Phi(X)$ dans $\und{Fibsc}_0\,S$, car $X$, provenant de $\und{Fibsc}_0\,S$, est dans $\und{Fibhot}\,S$, et on applique la remarque `Cor.~3' page 21.) D'autre part, on a une flèche dans $\Cat/S$
\[ X \xrightarrow{i_X} \Phi(X), \]
est-elle une flèche dans $\und{Fibsc}_0(S)$, i.e. est-elle cartésienne? On voit tout de suite que non, que l'on a seulement, pour une flèche $u : s \longrightarrow t$ dans $S$, un homomorphisme de commutation (pas un isomorphisme) dans

% ---- from Der-ChVI-345_p021-025.tex
\printedpage{19}
% head: A. Grothendieck : Dérivateurs, Ch. VI, Édition des Universités de Montpellier II et Paris VII
% foot: Version 9 novembre 2001

\[
\begin{tikzcd}[row sep=large, column sep=large]
X_s \arrow[r, "u_*^X"] \arrow[d, "i_s"] & X_t \arrow[d, "i_t"] \\
X/s \arrow[r, "(\mathrm{id}_X)/u", "= u_*^{\Phi(X)}"'] \arrow[ur, bend left=30, shorten <=28pt, shorten >=12pt, "\alpha"] & X/t
\end{tikzcd}
\]

\msp{page 31}
Donc il me faut remballer! On peut considérer $i_X$ au mieux comme une flèche de $\uu{Fib}_0\, S$, et en prenant les composés on trouve une flèche dans $\uu{Fib}_0\, S$ :
\[
\begin{tikzcd}
X \arrow[r, "i_X"] \arrow[rr, bend right=25, "\beta_X(=\alpha_X)"'] & \Phi(X) \arrow[r, "\Phi(j_X)"] & \Phi C(X),
\end{tikzcd}
\]
i.e. (regardant maintenant $\Phi_0$ comme foncteur à valeurs dans $\uu{Fib}_0\, S$) un homomorphisme de foncteurs
\[
\beta : \mathrm{id}_{\uu{Fib}_0} S \longrightarrow \Phi_0 p_0, % [sic] \mathrm{id}_{\uu{Fib}_0} S -> \mathrm{id}_{\uu{Fib}_0 S}, the S is printed on the baseline, not in the subscript
\]
où $p_0 : \uu{Fib}_0 S \hookrightarrow \mathrm{Cat}/S$ est l'inclusion canonique. Ce $\beta$ est dans $W$ argument par argument, et de même pour $\alpha$ page 29, interprété cette fois comme
\[
\mathrm{id}_{\mathrm{Cat}/S} \longrightarrow p_0\Phi_0.
\]
Ce qu'on a prouvé finalement est ceci:

\dgrblock{dgrb208}{Why the obvious unit fails to be Cartesian}
\begin{dcomment}{exposition}{p30}{Why the obvious unit fails to be Cartesian}
Grothendieck has found a natural weak equivalence of total categories and now asks it to be a morphism of coherent diagrams. The request is stronger. For the fibration construction \ensuremath{\Phi}(X) the fibre at s is s\textbackslash{}X, whereas the obvious map sends x \ensuremath{\in}Xs to (x, 1s). A Cartesian lift u\ensuremath{*}y \ensuremath{-}\ensuremath{\to}y in X, lying over u : s \ensuremath{-}\ensuremath{\to}t, is sent to an arrow whose source is (u\ensuremath{*}y, 1s). The chosen Cartesian lift in \ensuremath{\Phi}(X) ending at (y, 1t) instead starts at (y, u). These two objects are connected by the lift in X, but in general they are not isomorphic in s\textbackslash{}X.

The smallest example is X = S = [1] with the identity projection. Every arrow of X is Cartesian over S. In 0\textbackslash{}[1] the two objects 10 : 0 \ensuremath{-}\ensuremath{\to}0 and u : 0 \ensuremath{-}\ensuremath{\to}1 are joined by a noninvertible arrow. Thus the image of the Cartesian arrow u under iX is not Cartesian in \ensuremath{\Phi}(X). Indeed, in the Grothendieck construction of a strict diagram, an arrow over u is Cartesian exactly when its vertical comparison component is invertible. Here that comparison component is the displayed noninvertible arrow. The obstruction remains even though both comma categories are contractible.

The square printed across manuscript pages 30--31 uses X/s, X/t and covariant symbols u\ensuremath{*}. Those are the symbols of the dual construction \ensuremath{\Psi}, while the surrounding discussion concerns \ensuremath{\Phi}. The intended argument is correct with s\textbackslash{}X, t\textbackslash{}X and contravariant transport, or after dualizing the entire discussion. This variance slip should not be turned into a purported failure of the comma construction itself.

Passing to the full category Fib0S permits the map iX, because its morphisms are arbitrary functors over S. The natural weak equivalence X \ensuremath{-}\ensuremath{\to}\ensuremath{\Phi}C(X) then proves the coarse comparison in the “théorème-remballage”. The theorem does not yet compare strict coherent diagrams by fibrewise weak equivalences. This is the exact boundary between what the argument has proved and what the next attempt must establish. The later Cartesian comparison is directed towards X, rather than away from it; its construction supplies the missing information about coherence.

\end{dcomment}


\textbf{Théorème-remballage.} \textit{Soit $\mathcal{F}$ une sous-catégorie pleine de $\mathrm{Cat}/S$ qui contient une des trois sous-catégories pleines $\uu{Fib}_0 S$, $\uu{Parf}\, S$, $\uu{Cofib}_0 S$, et munissant $\mathcal{F}$ du localiseur $W$ `grossier' (induit par celui de $\mathrm{Cat}/S$). Alors l'inclusion}
\[
\mathcal{F} \longrightarrow \mathrm{Cat}/S
\]
\textit{induit une équivalence sur les localiseurs.}

\msp{page 32}
Par contre, les catégories localisées des quatre catégories de fibrés scindés
\[
\uu{Fibsc}\, S, \quad \uu{Cofibsc}\, S, \quad \uu{Fibsc}_0 S, \quad \uu{Cofibsc}_0 S
\]
échappent à nos arguments. \textit{Et pour cause!} Il est totalement ouvert, sur $S$ connexe disons (p. ex. $S = \mathbf{\Delta}^1$), que pour un type d'homotopie relatif $\xi$ sur $S$, les types induits $\xi_s$ ($s \in S$) forment un système local sur $S$, à valeurs dans $\mathrm{Hot}_W$, comme ce serait le cas si le théorème annoncé au début était vrai. Au contraire, \textit{toute} flèche de $\mathrm{Hot}_W$ doit provenir d'un objet de $\mathrm{Hot}_W(\mathbf{\Delta}^1)$. Les localisés des 9 autres catégories du diagramme p. 8, et de toutes celles dont il est question dans le théorème qu'on vient de prouver, doivent correspondre `visiblement' aux sous-catégories pleines de $\mathrm{HOT}_W(S)$ formées des HOT-coefficients localement constants sur $S$.

Aussi, je me rends compte que c'était absurde de m'attendre à ce que $\uu{Fibsc}\, S$ et $\uu{Cofibsc}\, S$ donnent des catégories localisées canoniquement équivalentes. Car les éléments $\xi$ de la première
\printedpage{20}
% head: A. Grothendieck : Dérivateurs, Ch. VI, Édition des Universités de Montpellier II et Paris VII
% foot: Version 9 novembre 2001
\msp{page 33}
des deux catégories donnent naissance chacun à un foncteur contravariant sur $S$
\[
\bar{\xi} : S^o \longrightarrow \mathrm{Hot}_W,
\]
ceux $\eta$ de la seconde (qui n'est autre que $(\uu{Fibsc}(S^o))^o$) donnent naissance chacun à un foncteur \textit{covariant}
\[
\bar{\eta} : S \longrightarrow \mathrm{Hot}_W,
\]
ce qui n'est pas pareil! Ce n'est `pareil' que si ce sont des coefficients locaux, ce qui signifie exactement que $\xi$ est dans $\uu{Fibsc}_0 S$ resp. $\eta$ dans $\uu{Cofibsc}_0 S$. On s'attend donc bien \circled{$1^o$} que les flèches
\[
\uu{Fibsc}_0 S \longrightarrow \uu{Fib}_0 S, \quad \uu{Cofibsc}_0 S \longrightarrow \uu{Cofib}_0 S
\]
induisent des équivalences sur les localisés (chose que j'ai été incapable de prouver tantôt) (\setcounter{footnote}{15}\footnote{Les quasi-inverses étant bel et bien donnés par $\Phi_0$, $\Psi_0$ - car [?] par quoi d'autre donc ???}). Par contre, il est pire que déraisonnable de s'attendre à ce qu'il en soit de même pour les foncteurs
\[
\uu{Fibsc}_0 S \longrightarrow \uu{Fibsc}\, S, \quad \uu{Cofibsc}_0 S \longrightarrow \uu{Cofibsc}\, S,
\]
ça devrait donner par localisation les foncteurs pleinement fidèles d'inclusion, des catégories d'objets localement constants dans $\mathrm{HOT}_W(S)$ resp. dans $(\mathrm{HOT}_W(S^o))^o$. Et de même pour
\msp{page 34}
les foncteurs induits sur les localisés par
\[
(*) \qquad \uu{Fibsc}\, S \longrightarrow \uu{Fib}\, S, \quad \uu{Cofibsc}\, S \longrightarrow \uu{Cofib}\, S.
\]
Comme les localisés des catégories but de ces deux foncteurs ont tout l'air de s'identifier aux sous-catégories précédentes d'objets localement constants, il semblerait donc que : \circled{$2^o$} Les foncteurs précédents définissent des rétractions de la catégorie $\mathrm{HOT}_W(S)$ resp. $\mathrm{HOT}_W(S^o)^o$ sur la sous-catégorie des objets localement constants sur $S$, les foncteurs localisés des foncteurs d'inclusion
\[
\uu{Fibsc}_0 S \longrightarrow \uu{Fibsc}\, S, \quad \uu{Cofibsc}_0 S \longrightarrow \uu{Cofibsc}\, S
\]
étant pleinement fidèles et ayant comme images essentielles les sous-catégories en question.

\dgrblock{dgrb209}{A relative type which is not locally constant}
\begin{dcomment}{examples}{p32}{A relative type which is not locally constant}
Grothendieck’s reversal becomes transparent on the walking arrow. Let S = [1] = \{0 \ensuremath{-}\ensuremath{\to}1\}, let X be the one-object category over 0, and let Y = S with its identity projection. The inclusion

f : \{0\} \ensuremath{-}\ensuremath{\to}[1]

is a map of split fibrations over S. The corresponding contravariant diagrams are

X1 = \ensuremath{\varnothing}\ensuremath{-}\ensuremath{\to}X0 = \ensuremath{*},

Y1 = \ensuremath{*}\ensuremath{-}\ensuremath{\to}Y0 = \ensuremath{*}. Both total categories are contractible, so f \ensuremath{\in}W\ensuremath{\infty}. But the fibre at 1 is \ensuremath{\varnothing}\ensuremath{-}\ensuremath{\to}\ensuremath{*}, which is not a weak homotopy equivalence. Fibre evaluation is a functor on the localization by objectwise weak equivalences, and sends isomorphisms there to isomorphisms of homotopy types. Therefore the image of f in that localization is not invertible. Since its image in the coarse localization is invertible, the comparison between these two localizations is not an equivalence.

All four tests of manuscript page 41 can be computed here. For the overcategories,

X/0 = \ensuremath{*},

Y/0 = \ensuremath{*}, X/1 = \ensuremath{*}, Y/1 = [1], so f is a local, or left-fine, weak equivalence. For the undercategories,

1\textbackslash{}X = \ensuremath{\varnothing},

0\textbackslash{}X = \ensuremath{*}, 0\textbackslash{}Y = [1], 1\textbackslash{}Y = \ensuremath{*}. Thus it fails the colocal, or right-fine, test. Pulling back along the point 1 \ensuremath{-}\ensuremath{\to}S gives the same failure, so it is not a universal weak equivalence. The example lies inside the category of split fibrations, hence also inside the category of smooth functors. Its opposite supplies the corresponding example for opfibrations and proper functors. This will also resolve the contradictory sentence on manuscript page 47.

An object equivalent to a locally constant diagram must have isomorphic fibre homotopy types along every base arrow: invertibility is invariant under isomorphism in the homotopy category of diagrams. The first diagram above does not have this property. Accordingly, the full subcategory of locally constant objects is proper even when the classifying space of the indexing category is contractible. Contractibility of the base makes local systems constant up to equivalence; it does not force every coherent diagram to be locally constant.

This gives a definite negative answer to the question on manuscript page 25 for the fibrewise localization of smooth or fibred objects. It also explains why the initial thirteen-category comparison could not have been correct with a single choice of weak equivalences. A derived category of diagrams must retain the noninvertible transition \ensuremath{\varnothing}\ensuremath{-}\ensuremath{\to}\ensuremath{*}, whereas the total homotopy type in this example forgets it.

\end{dcomment}


Cette rétraction m'apparait comme une chose tout à fait mystérieuse. C'est lié au fait que pour le dérivateur très spécial $\mathbf{D} = \mathrm{HOT}$, quand on a un $\xi \in \mathbf{D}(S)$, on en déduit (à $\mathrm{HOT}/S$-équivalence près) un type de diagramme $X$ au dessus de $S$,
\msp{page 35}
qu'on peut regarder alors (via la construction $C$) comme définissant un type d'homotopie relatif \textit{localement constant} sur $S$.

\printedpage{21}
% head: A. Grothendieck : Dérivateurs, Ch. VI, Édition des Universités de Montpellier II et Paris VII
% foot: Version 9 novembre 2001

Donc pour élucider la situation, i.e. les relations entre les 13 catégories localisées provenant du diagramme de la page 8, il reste à regarder si les assertions hypothétiques que je viens de faire sont bien valables, savoir

\begin{enumerate}[leftmargin=*]
\item[\circled{$1^o$}] Le foncteur déduit de
\[
\uu{Fibsc}_0 S \longrightarrow \uu{Fib}_0 S
\]
par passage aux localisés est une équivalence de catégories (et on peut dans cette question remplacer $\uu{Fib}_0 S$ par n'importe quelle autre sous-catégorie pleine de $\mathrm{Cat}/S$ la contenant - cela revient au même), et

\item[\circled{$2^o$}] Le foncteur déduit de
\[
\uu{Fibsc}_0 S \longrightarrow \uu{Fibsc}\, S
\]
par localisation est pleinement fidèle. (Que son image soit exactement formée des objets localement constants du premier localisé, interprété comme $\mathrm{HOT}_W(S)$, est tautologique.)
\end{enumerate}

Il n'y a plus lieu de se soucier du côté
\msp{page 36}
cofibrés, puisque tout ce qu'on pourra en dire revient à ce que je viens d'expliciter, pour $X^o$ au lieu de $X$.

Dans la direction de \circled{$1^o$}, nos arguments précédents à coups de $\Phi_0$, $p$ nous donnent déjà un renseignement utile. Mais il sera plus commode à présent de laisser tomber $\uu{Fib}_0 S$ ou plutôt, de le remplacer par $\uu{Fib}\, S$ (ce qui est licite, on l'a dit dans \circled{$1^o$}), de sorte qu'on a maintenant un diagramme commutatif (traits pleins)
\[
\begin{tikzcd}
\uu{Fibsc}\, S \arrow[dr, "j"'] & & \uu{Fibsc}_0 S \arrow[ll, hook', "i"'] \arrow[dl, "k"'] \\
& \uu{Fib}\, S, \arrow[ur, dotted, bend right=40, "\Phi_0=\psi"'] &
\end{tikzcd}
\]
où $i$ est une inclusion (strictement) pleinement fidèle, $j$ (donc aussi $k$) fidèle. Enfin, on a un foncteur canonique $\psi$ ($= \Phi_0$) (en pointillés). Les localiseurs sur la ligne du haut sont les localiseur fins (pour $\uu{Fibsc}_0 S$ c'est pareil que la $W$-équivalence grossière), sur $\uu{Fib}\, S$ % [sic] les localiseur fins -> les localiseurs fins, missing plural s
par contre on prend la $W$-équivalence grossière. (On sait que le localisé est le même, à équivalence près, que si on se restreint à la sous-catégorie pleine $\uu{Fib}_0 S$, ce qui est assez remarquable. % [sic] "(On sait ... remarquable." -> "(On sait ... remarquable.)", the opening parenthesis is never closed

\msp{page 37}
Tous les quatre foncteurs $i, j, k, \Phi_0$ sont compatibles aux localiseurs, donc passent aux localisés
\printedpage{22}
% head: A. Grothendieck : Dérivateurs, Ch. VI, Édition des Universités de Montpellier II et Paris VII
% foot: Version 9 novembre 2001
\[
\begin{tikzcd}
\mathrm{HOT}_W(S) \arrow[dr, "\bar{j}"'] & & \mathrm{HOT}^{?}_W(S) \arrow[ll, "\bar{\imath}"'] \arrow[dl, "\bar{k}"'] \\
& \mathrm{HOT}^{\mathrm{lc}}_W(S), \arrow[ur, bend right=40, "\bar{\psi}"'] &
\end{tikzcd}
\]
où l'exposant ? dans la notation signifie que je ne sais interpréter cette catégorie, sauf si je prouve que $\bar{k}$ est une équivalence, et l'exposant lc dans la catégorie du bas dois suggerer l'intuition qu'il s'agit en tous cas d'une sorte de `types d'homotopie relatifs' (dans un sens modifié) \textit{localement constants} sur $S$. Le problème est de comprendre ce diagramme. La conjecture naturelle est % [sic] dois suggerer -> doit suggérer, verb agreement and missing accent

\begin{enumerate}[leftmargin=*]
\item[$1^o$)] $\bar{k}$ est une équivalence. (Ce qui implique que $\bar{\psi}$ est l'équivalence quasi-inverse, que $\bar{k}$, $\bar{\psi}$ permettent essentiellement d'identifier $\mathrm{HOT}^{?}_W$ et $\mathrm{HOT}^{\mathrm{lc}}_W$, de sorte que $\bar{\imath}$ et $\bar{\jmath}$ deviennent des foncteurs entre $\mathrm{HOT}_W$ et $\mathrm{HOT}^{\mathrm{lc}}_W$.)

\item[$2^o$)] $\bar{\imath}$ est pleinement fidèle (NB l'image essentielle est bien comprise de toutes façons.)
\end{enumerate}

\msp{page 38}
En direction de $1^o$), je sais pour l'instant que l'on a
\[
\bar{k}\bar{\psi} \simeq \mathrm{id}_{\mathrm{HOT}^{\mathrm{lc}}_W(S)},
\]
donc $\bar{k}$ est essentiellement surjectif (sur les objets), et induit des \textit{surjections} sur les $\mathrm{Hom}(-,-)$, tandis que $\bar{\psi}$ est injectif sur les classes d'isomorphismes d'objets, et sur les $\mathrm{Hom}(-,-)$ (i.e. \textit{fidèle}).

Si les relations entre $\mathrm{HOT}^{?}_W$ et $\mathrm{HOT}^{\mathrm{lc}}_W$ s'avèrent imbitables (p. ex. $\bar{\psi}$ pas même pleinement fidèle, i.e. $\bar{k}$ pas même fidèle), alors la catégorie $\mathrm{HOT}^{?}_W$ mérite d'être larguée dans la corbeille, et on ne gardera que
\[
\begin{tikzcd}[row sep=large]
\mathrm{HOT}_W(S) \arrow[d, shift right=1.5, "\bar{j}"'] \\
\mathrm{HOT}^{\mathrm{lc}}_W(S), \arrow[u, shift right=1.5, "\bar{\psi}'"']
\end{tikzcd}
\]
où on définit $\bar{\psi}'$ comme
\[
\bar{\psi}' = \bar{\imath}\bar{\psi} = \overline{i\psi}.
\]
Et la question (plus modeste) devient à présent :

\begin{enumerate}[leftmargin=*]
\item[$3^o$)] $\bar{\psi}'$ est-il pleinement fidèle? (L'image \textit{essentielle} est la même que celle de $\uu{Fib}_0$ ($\approx$ $\uu{Fib}$), elle est formée sauf erreur des objets localement constants - à vérifier au besoin \dots)
\end{enumerate}

\printedpage{23}
% head: A. Grothendieck : Dérivateurs, Ch. VI, Édition des Universités de Montpellier II et Paris VII
% foot: Version 9 novembre 2001

\msp{page 39}
Il est difficile d'imaginer que $\bar{k}$, ou ce qui revient au même $\bar{\psi}$, puisse être une équivalence, sans qu'il existe au moins une `$W_S$-homotopie'
\[
\mathrm{id}_{\uu{Fibsc}_0 S} \ \leftrightarrow\ \leftrightarrow^{h}\!\leftrightarrow\ \leftrightarrow\ \psi k
\]
dans la catégorie des endofoncteurs de $\uu{Fibsc}_0 S$. À vrai dire, je me rends compte maintenant que la description de $\bar{\psi}$ à partir de $\psi : X \longmapsto \Phi_0 X = \Phi C X$ est inutilement compliquée, car on avait fait le détour par $C(X)$ uniquement pour être sûr de tomber dans $\uu{Fibsc}_0 S$, pas seulement dans $\uu{Fibsc}\, S$. Mais si on laisse tomber $\uu{Fibsc}_0 S$, on s'en fout, et le foncteur $\Phi : X \longmapsto \Phi(X)$ est beaucoup plus naturel. Et si on ne laisse pas tomber, il n'y a pas d'inconvénient de remplacer la source $\uu{Fib}\, S$ par $\uu{Fib}_0 S$ (qui a même localisé, à équivalence près), et sur $\uu{Fib}_0 S$ à nouveau il n'y a pas lieu de s'encombrer de $CX$, pour le plaisir de lui appliquer $\Phi$, car $\Phi(X)$ lui-même est déjà dans $\uu{Fibsc}_0 S$.

\dgrblock{dgrb210}{What a one-sided inverse does and does not give}
\begin{dcomment}{examples}{p38}{What a one-sided inverse does and does not give}
Grothendieck has a useful piece of the comparison: a natural isomorphism k\textasciimacron{} \ensuremath{\psi}\textasciimacron{} \ensuremath{\simeq}1. It proves essential surjectivity of k\textasciimacron{} and faithfulness of \ensuremath{\psi}\textasciimacron{}, and it shows that \ensuremath{\psi}\textasciimacron{} is injective on isomorphism classes of objects. The statement about surjectivity on all morphism sets requires an additional hypothesis.

To locate the issue, let A be the discrete category on two objects a0, a1, and let B = \ensuremath{*}be terminal. There is a unique functor k : A \ensuremath{-}\ensuremath{\to}B; let \ensuremath{\psi} pick a0. Then k\ensuremath{\psi} = 1B strictly. Yet

HomA(a0, a1) = \ensuremath{\varnothing}\ensuremath{-}\ensuremath{\to}HomB(\ensuremath{*}, \ensuremath{*}) = \{1\}

is not surjective. Thus a functor with a section need not be full. In this example \ensuremath{\psi} is even fully faithful, so the parenthetical suggestion that failure of faithfulness of k is equivalent to failure of full faithfulness of \ensuremath{\psi} cannot be justified solely from the retract relation. For that converse take instead A = B \ensuremath{\coprod}BM, with B = \ensuremath{*}and a nontrivial monoid M, let k be the unique functor to \ensuremath{*}and let \ensuremath{\psi} select the first summand. Then \ensuremath{\psi} is fully faithful while k is not faithful.

The precise positive statement is easy and sufficient for the progress claimed so far. Choose

\ensuremath{-}\ensuremath{\to}1. For b, b\ensuremath{\prime} \ensuremath{\in}B and h : b \ensuremath{-}\ensuremath{\to}b\ensuremath{\prime}, the morphism \ensuremath{\psi}(h) maps under k, followed by conjugation \ensuremath{\varepsilon} : k\ensuremath{\psi} with \ensuremath{\varepsilon}, to h. Hence the map

HomA(\ensuremath{\psi}b, \ensuremath{\psi}b\ensuremath{\prime}) \ensuremath{-}\ensuremath{\to}HomB(b, b\ensuremath{\prime})

is split surjective. If \ensuremath{\psi}(h) = \ensuremath{\psi}(h\ensuremath{\prime}), applying k and conjugating gives h = h\ensuremath{\prime}, proving faithfulness. If \ensuremath{\psi}b \ensuremath{\simeq}\ensuremath{\psi}b\ensuremath{\prime}, the same argument gives b \ensuremath{\simeq}b\ensuremath{\prime}. None of these statements identifies an arbitrary a \ensuremath{\in}A with \ensuremath{\psi}ka.

A natural zigzag of weak equivalences between 1 and \ensuremath{\psi}k, sought on the following pages, would supply exactly that missing identification after localization. The full comparison theorem can therefore remain correct although the present formal inference is too strong. We record an error in this intermediate inference, with its corrected scope, rather than a gap in the later comparison once that zigzag has been constructed.

\end{dcomment}


Donc, la question est celle-ci : Si $X$ varie
\msp{page 40}
% layout: the arrows nearrow, uparrow, nwarrow are printed under each "Fibsc_0 S", pointing up to it; the arrow labelled c from X to Phi(X) is a long line with its arrowhead printed in the middle
\[
\begin{array}{l}
\text{dans } \underset{\nearrow\uparrow\nwarrow}{\uu{Fibsc}_0 S}, \quad \text{y-a-t'il une homotopie dans } \underset{\nearrow\uparrow\nwarrow}{\uu{Fibsc}_0 S}, \quad \text{`naturelle', i.e. fonctorielle en } X, \\[2ex]
\qquad\qquad X \xrightarrow{\qquad c \qquad} \Phi(X) = \ \text{Catégorie fibrée scindée de fibres les } s\backslash X
\end{array}
\]
Mais un chemin qui soit formé de flèches
\[
X = X_0 \lline X_1 \lline X_2 \cdots \lline X_n = \Phi(X)
\]
qui soient toutes cartésiennes (et de plus, bien sûr, dans $W$!). Cela signifie donc qu'on aurait des catégories fibrées scindées sur $S$, $X_1$, \dots, $X_{n-1}$, dépendant \textit{fonctoriellement} de $X$, qui relient
\[
X = X_0 \qquad \text{à} \qquad \Phi(X) = X_n
\]
i.e.
\[
X_s \qquad \text{à} \qquad s\backslash X \qquad \text{???}
\]
Je ne vois rien à l'horizon.

% ---- from Der-ChVI-345_p026-030.tex
\printedpage{24}
% head: A. Grothendieck : Dérivateurs, Ch. VI, Édition des Universités de Montpellier II et Paris VII

\msp{page 41}

\dgrblock{dgrb211}{The locally constant comparison is now a theorem}
\begin{dcomment}{status}{p40}{The locally constant comparison is now a theorem}
Grothendieck is asking for morphisms that retain coherent transport while forgetting arbitrary choices of Cartesian lifts. His unsuccessful unit is not evidence that the locally constant subcategory fails to embed fully. It identifies a defect of that particular representative of the comparison. The arguments on manuscript pages 55--60 repair the comparison, and pages 65--69 establish the relevant reflection.

The higher coherence retained by local constancy was illustrated at manuscript page 7. The later theory treats the local constancy question for a general derivator. Cisinski (2009, Proposition 3.10 and Theorem 3.13) proves that locally constant objects form a full subderivator and that its inclusion has both a left and a right adjoint. These are the published numbering; the earlier arXiv version numbers the corresponding statements 3.9 and 3.12. The theorem establishes the positive status of the full embedding and the existence of the two adjoints, without identifying all the coarse localizations of this chapter for an arbitrary W.

That last restriction matters. If W is homological or a finite homotopy truncation, a coarse Wequivalence of total spaces need not be a W-equivalence on homotopy fibres. The local constancy construction in a derivator still makes sense, but the additional identification with the coarse localization of Cat/S uses the stronger hypothesis isolated below. For W\ensuremath{\infty}the categorical path model and the modern description by coherent space-valued diagrams fit together; the commentary on manuscript pages 65--69 gives the precise reflector. This question is therefore settled in its intended homotopy-theoretic setting, while its unqualified extension to every localizer is false.

\end{dcomment}


\begin{flushright}
15.11.90
\end{flushright}

Je crois qu j'ai enfin compris, cette nuit, la situation exacte. Il y a lieu, dans Cat$/S$ et dans les autres catégories du diagramme p. 8 (s'envoyant toutes dans Cat$/S$) d'envisager quatre sortes de ‘$W$-équivalences’. % [sic] qu j'ai -> que j'ai, letter missing

\begin{itemize}
\item[a)] La \textit{$W$-équivalence ‘grossière’} : \begin{tikzcd}[baseline=(s.base), row sep=small, column sep=small]
f : X \arrow[rr] \arrow[dr] & & Y \arrow[dl] \\
& |[alias=s]| S &
\end{tikzcd} dans Cat$/S$ est une $W$-équivalence grossière si en tant que flèche de Cat, c'est une $W$-équivalence. (\footnote[17]{stable par changement de base à la fois lisse et propre, plus généralement si $S' \longrightarrow S$ est une Hot-fibration.}).

\item[b)] La \textit{$W$-équivalence ‘par fibres’} (sur $S$) : $\forall\ s \in S$, $f_s : X_s \longrightarrow Y_s$ est une $W$-équivalence. (\footnote[18]{stable par tout chagement de base $S' \longrightarrow S$}). % [sic] chagement -> changement, letter missing

\item[c)] La \textit{$W$-équivalence ‘locale sur $S$’} : $\forall\ s \in S$, $f/s : X/s \longrightarrow Y/s$ est une $W$-équivalence. Signifie que c'est une $W$-équivalence, et le reste par tout changement de base $S' \longrightarrow S$ qui est un isomorphisme local. (\footnote[19]{devrait être stable par changement de base \textit{lisse} $S' \longrightarrow S$}).

\item[c$'$)] La \textit{$W$-équivalence ‘colocale sur $S$’} : $\forall\ s \in S$, $s\backslash f : s\backslash X \longrightarrow s\backslash Y$ est une $W$-équivalence. Signifie que c'est une $W$-équivalence, et le reste par tout changement de base $S' \longrightarrow S$ qui est un co-isomorphisme local. (\footnote[20]{devrait être stable par changement de base \textit{propre} sur $S$}).
\end{itemize}

Je vais résumer dans un énoncé récapitulatif les relations entre ces quatre notions de $W$-équivalence, en pensant au cas $W = W_\infty$. Une fois la chose faite, je me préoccuperai de dégager les conditions minimales sur un $W$ général qui assurent que c'est vrai pour $W$.

\textbf{Proposition 1.} \textit{Soit} \begin{tikzcd}[baseline=(s.base), row sep=small, column sep=small]
f : X \arrow[rr] \arrow[dr] & & Y \arrow[dl] \\
& |[alias=s]| S &
\end{tikzcd} \textit{dans} Cat$/S$.

\begin{itemize}
\item[1)] \textit{La $W$-équivalence locale sur $S$, et la $W$-équivalence colocale sur $S$, impliquent la $W$-équivalence grossière :}
\[
\begin{tikzcd}
\text{c} \arrow[dr, Rightarrow] & & \text{c}' \arrow[dl, Rightarrow] \\
& \text{a} &
\end{tikzcd}
\]

\printedpage{25}
% head: A. Grothendieck : Dérivateurs, Ch. VI, Édition des Universités de Montpellier II et Paris VII

\item[2)] \textit{Si $X$, $Y$ sont propres sur $S$, ou $X$, $Y$ lisses sur $S$, ou $X$, $Y$ des} Hot-\textit{fibrés sur $S$, alors la $W$-équivalence par fibres implique que c'est une $W$-équivalence} universelle \textit{sur $S$ ($W$-équivalence, et le reste par} tout \textit{changement de base $S' \longrightarrow S$), donc dans ces trois cas, on a}
\[
\begin{tikzcd}
& \text{b} \arrow[dl, Rightarrow] \arrow[dr, Rightarrow] & \\
\text{c} \arrow[dr, Rightarrow] & & \text{c}' \arrow[dl, Rightarrow] \\
& \text{a} &
\end{tikzcd}
\]
(\footnote[21]{Il suffit $X$, $Y$ $W$-propres sur $S$, ou $W$-lisses sur $S$, ou $W$-fibrés sur $S$.}).

Mais on peut préciser chacun de ces trois cas :

\item[(2 fib)] \textit{Si $X$, $Y$ sont} Hot-\textit{fibrés sur $S$, alors la $W$-équivalence grossière de $f$ implique que c'est une $W$-équivalence par fibres, donc les quatre notions de $W$-équivalence sont équivalentes pour $f$.}
\[
\text{a} \Longleftrightarrow \text{c} \Longleftrightarrow \text{c}' \Longleftrightarrow \text{b}
\]
(\footnote[22]{a $\Longrightarrow$ b seulement moyennant L5, voir ci-dessous p. 43.})

\item[(2 lisse)] \msp{page 42}\textit{Si $X$, $Y$ sont lisse sur $S$, alors la $W$-équivalence par fibres sur $S$ équivaut à la $W$-équivalence colocale sur $S$ et implique la $W$-équivalence} universelle \textit{sur $S$, donc dans ce cas} % [sic] lisse -> lisses, plural agreement with X, Y
\[
\begin{tikzcd}
\text{b} \arrow[r, Leftrightarrow] \arrow[d, Rightarrow] & \text{c}' \\
\text{c} \arrow[d, Rightarrow] & \\
\text{a} &
\end{tikzcd}
\]

\dgrblock{dgrb212}{Four tests on a map over the base}
\begin{dcomment}{exposition}{p41}{Four tests on a map over the base}
Grothendieck now separates the total homotopy type from the two possible diagram directions. For p : X \ensuremath{-}\ensuremath{\to}S there are three kinds of fibre to consider. The strict fibre Xs remembers only objects and arrows mapping to s and 1s. The comma category X/s also allows arrows p(x) \ensuremath{-}\ensuremath{\to}s in the base; it varies covariantly with s. The comma category s\textbackslash{}X allows arrows s \ensuremath{-}\ensuremath{\to}p(x) and varies contravariantly. Thus for f : X \ensuremath{-}\ensuremath{\to}Y over S the tests are

total

f \ensuremath{\in}W,

fibrewise fs \ensuremath{\in}W (s \ensuremath{\in}S),

local, or left-fine f/s \ensuremath{\in}W (s \ensuremath{\in}S),

colocal, or right-fine s\textbackslash{}f \ensuremath{\in}W (s \ensuremath{\in}S). The word “local” here refers to these particular comma tests, not to a chosen Grothendieck topology. The names “left” and “right” are the manuscript’s conventions. They must not be interchanged merely because the fibrations are drawn on the left side of the thirteen-category diagram.

For the ordinary weak homotopy equivalences, the role of the comma tests is to detect the maps of coherent diagrams produced by \ensuremath{\Psi} and \ensuremath{\Phi}. A map can preserve the total homotopy type while changing one of those diagrams, as the walking-arrow example at manuscript page 32 proves. Conversely, strict fibres alone do not detect a total equivalence without some condition on the maps to S. For example, the functor from the discrete two-object category \{0, 1\} to [1], over [1], is an isomorphism on both strict fibres but changes the total space from two components to a contractible space. It also fails both comma tests: over 1, and under 0, respectively, it compares a discrete two-point category to a contractible one.

The assertion concerning arbitrary base change of a fibrewise equivalence is formal: the fibre of a pullback over s\ensuremath{\prime} \ensuremath{\in}S\ensuremath{\prime} is the original fibre over its image in S. The assertion concerning base change of a coarse equivalence along every homotopy fibration is not formal for a general W; it is exactly the additional stability later called L6 bis. Footnote 17 is valid for W\ensuremath{\infty}, the setting explicitly adopted for the proposition, and must not be read as a theorem about every localizer.

The local-isomorphism formulation has an elementary explanation in the usual categorical convention of discrete fibrations. If T \ensuremath{-}\ensuremath{\to}S is a discrete fibration, then

(X \ensuremath{\times}S T)/t \ensuremath{\sim}= X/p(t).

An object on the right determines the object of T above its source by the unique lifting property. Thus local equivalence is preserved by this pullback, and relative axiom L4 bis gives total equivalence afterward. Conversely the representable discrete fibrations S/s \ensuremath{-}\ensuremath{\to}S recover the comma tests themselves. Dualizing gives the colocal statement. This explains both the strength and the orientation of the local tests without equating them to arbitrary base-change invariance.

\end{dcomment}


\printedpage{26}
% head: A. Grothendieck : Dérivateurs, Ch. VI, Édition des Universités de Montpellier II et Paris VII
(\footnote[23]{il suffit même $X$, $Y$ $W$-lisses sur $S$.}).

\item[(2 prop)] \textit{Si $X$, $Y$ sont propres sur $S$, alors la $W$-équivalence par fibres sur $S$ équivaut à la $W$-équivalence locale sur $S$ :}
\[
\begin{tikzcd}
\text{b} \arrow[r, Leftrightarrow] \arrow[d, Rightarrow] & \text{c} \\
\text{c}' \arrow[d, Rightarrow] & \\
\text{a} &
\end{tikzcd}
\]
(\footnote[24]{Il suffit même $X$, $Y$ $W$-propres sur $S$.}).
\end{itemize}

\textsc{Démonstration} 1) n'est autre que L4 (\footnote[25]{Non, il y a confusion, L4 est plus faible, cf. p. 44 l'axiome L4 bis.}) et son dual L4$'$ pour un localiseur fondamental $W$, et on a vu que si on suppose L1, L3, alors L4 équivaut à L4$'$, et dans la notion de localiseur fondamental on supposera toujours que c'est vérifié. Également, nous supposons que L3 est vrai même sous la forme plus forte L3 quater (p. 1) : $W_\infty \subseteq W$, i.e. toute Hot-équivalence est une $W$-équivalence.

(2 lisse) Si $X$, $Y$ \texttt{[sont]} lisses sur $S$, alors les $X_s \longrightarrow s\backslash X$, $Y_s \longrightarrow s\backslash Y$ sont des $W_\infty$-équivalences, donc des $W$-équivalences (il suffirait même que $X$, $Y$ soient $W$-lisses$/S$), donc le diagramme commutatif
\[
\begin{tikzcd}
X_s \arrow[r, "f_s"] \arrow[d, "i_s^X"'] & Y_s \arrow[d, "i_s^Y"] \\
s\backslash X \arrow[r, "s\backslash f"] & s\backslash Y
\end{tikzcd}
\]
nous montre que $f_s$ est une $W$-équivalence si et seulement si $s\backslash f$ l'est, d'où b $\Longleftrightarrow$ c$'$. Pour voir que la $W$-équivalence par fibres est même une $W$-équivalence \textit{universelle} sur $S$, on est ramené à montrer que c'est une $W$-équivalence grossière (puisque la lissité \texttt{[est]} stable par changement de base sur $S$), et on gagne puisque c$'$ $\Longrightarrow$ a par 1).

Le cas (2 propre) est dual.

(2 fib) Si $X$, $Y$ sont des Hot-fibrés sur $S$, alors
\printedpage{27}
% head: A. Grothendieck : Dérivateurs, Ch. VI, Édition des Universités de Montpellier II et Paris VII
\msp{page 43}
les inclusions
\[
X_s \longrightarrow s\backslash X, \qquad X_s \longrightarrow X/s,
\]
et de même pour $Y$, sont des $W_\infty$-équivalences, donc des $W$-équivalences (il suffirait même que $X$, $Y$ soient des $W$-fibrations sur $S$ (\footnote[26]{comparer déf. 3, p. 117, dans (I).}), i.e. par définition que les foncteurs précédents soient des $W$-équivalences, et que la même condition soit satisfaite après tout changement de base sur $S$), l'argument dans le cas lisse resp. propre montre déjà que
\[
\begin{tikzcd}
\text{c} & \text{b} \arrow[l, Rightarrow] \arrow[r, Leftrightarrow] \arrow[d, Rightarrow] & \text{c}' \\
& \text{a,} &
\end{tikzcd}
\]
et que la $W$-équivalence par fibres (ou la $W$-équivalence locale sur $S$, resp. colocale sur $S$) implique déjà la $W$-équivalence \textit{universelle} sur $S$. (\footnote[27]{valable en supposant seulement $X$, $Y$ des $W$-fibrés sur $S$.})

\dgrblock{dgrb213}{When strict fibres compute a comma fibre}
\begin{dcomment}{exposition}{p42}{When strict fibres compute a comma fibre}
Grothendieck’s proposition is useful because it states exactly when a strict fibre is a valid homotopical substitute for one of the comma categories. The discussion of smoothness and properness in the commentary on Chapter~I, manuscript pages 52--59, supplies the geometric input; here we use that input to compare weak equivalences of objects over the same base.

If X \ensuremath{-}\ensuremath{\to}S is smooth, the canonical inclusion Xs \ensuremath{-}\ensuremath{\to}s\textbackslash{}X is a weak homotopy equivalence. When both X and Y are smooth, the commutative square

Xs

Ys

s\textbackslash{}X

s\textbackslash{}Y has vertical arrows in W if W\ensuremath{\infty}\ensuremath{\subseteq}W. Two-out-of-three proves the equivalence of the fibrewise and right-fine tests. For proper objects the same argument uses Xs \ensuremath{-}\ensuremath{\to}X/s and proves equivalence with the left-fine test. The proof requires no comparison of total spaces at this stage. The same argument works with W-smooth or W-proper objects when the vertical maps are directly known to be in W.

To obtain universality, make any change of base T \ensuremath{-}\ensuremath{\to}S. Smoothness and properness are preserved, and the fibrewise maps remain in W. The preceding comparison gives a colocal or local W-equivalence over T, respectively. Relative axiom L4 bis, or its dual, now makes the total map over T a W-equivalence. The base change was arbitrary, so the map is universally a W-equivalence. This fills the quantifier in the proof; it is not a new appeal to L6. Here base change acts on a map already known to be a fibrewise equivalence between suitably regular objects.

For homotopy fibrations both comma comparisons are weak homotopy equivalences, universally on the base. Hence the fibrewise, local and colocal conditions coincide and imply universal total equivalence. The reverse implication from a total equivalence is the extra statement later named L5. It holds for W\ensuremath{\infty} by invariance of homotopy fibres, and the next examples show why it cannot be appended for every W.

The modern extension to asphericity structures is given by Maltsiniotis (2005, Corollary 4.16 and Theorem 5.4): Grothendieck fibrations are smooth, and smoothness is characterized by the appropriate base-change property. This is the direct continuation of the present comparison argument. In highercategorical straightening one likewise distinguishes equivalence in the diagram direction from equivalence of underlying total spaces; the latter agrees with the former for locally constant fibrations, not for every right or left fibration.

\end{dcomment}


La question a $\Longrightarrow$ b est par contre nettement plus délicate, et on a vu que ce n'est nullement le cas pour tout localisateur fondamental - sans doute pas pour celui associé au dérivateur Der($\mathbf{Z}$-Mod), en tous cas pas à celui associé à Der($k$-Mod), quand 2 est inversible dans $k$. (Il suffirait même que $\exists\ n \in \mathbf{N}$, $n \geq 2$, tel que $n$ inversible dans $k$ - le contrexemple typique est $X = S^1 \xrightarrow{\ n\ } S^1 = Y = S$, le revêtement étale d'ordre $n$ du cercle.)

L'implication a $\Longrightarrow$ b est donc une condition très forte sur $W$, je vais la noter L5 (l'ancienne L5, cf. page 2, ne méritant pas de survivre, car c'est une côte mal taillée).

\textbf{L5} \textit{Soit} \begin{tikzcd}[baseline=(s.base), row sep=small, column sep=small]
X \arrow[rr, "f"] \arrow[dr] & & Y \arrow[dl] \\
& |[alias=s]| S &
\end{tikzcd} \textit{dans} Cat\textit{, avec $X$, $Y$ des} Hot-\textit{fibrés sur $S$. Si $f$ est une $W$-équivalence (‘grossière’ par rapport à $S$), alors c'est une $W$-équivalence par fibres} (\footnote[28]{Dans XII (p. 12) on voit que si $W$ satisfait à L1, L3 quater, L4 bis (cf. page suivante), L5, alors \fbox{$W = W_\infty$}. Donc c'est parfaitement illusoire, si on suppose L5 de continuer à traîner un $W$.}).

On a vu qu'à un $\varepsilon$ près (‘lemme’ p. 23), L5 est une conséquence de

\textbf{L6} \textit{Soit $f : X \longrightarrow Y$ une $W$-équivalence, $Y' \longrightarrow Y$ un morphisme parfait (i.e. lisse et propre), alors $f' : X' = X \times_Y Y' \longrightarrow Y'$ est une $W$-équivalence.}

\msp{page 44}
Notons la variante bien naturelle

\dgrblock{dgrb214}{The circle covering as a categorical counterexample}
\begin{dcomment}{examples}{p43}{The circle covering as a categorical counterexample}
Grothendieck’s covering of the circle gives a complete obstruction to L5 for homology coefficients in which its degree is invertible. It is useful to realize it inside Cat, so that no passage from topological coverings to categorical fibres is left implicit.

Fix n \ensuremath{\geq}2. Let S = BZ, the one-object category with endomorphism group Z. Let X be the action groupoid for translation of Z on Z/n. Its objects are residues a, and a morphism a \ensuremath{-}\ensuremath{\to}b is an integer m with b = a + m modulo n; composition is addition. The projection p : X \ensuremath{-}\ensuremath{\to}S sends that morphism to m. Unique lifting of each integer with specified source or target makes p both a discrete fibration and a discrete opfibration. Its strict fibre is the discrete set Z/n.

The groupoid X is connected. The automorphism group of the object 0 is nZ, so the inclusion B(nZ) \ensuremath{-}\ensuremath{\to}X is an equivalence. Under this identification, p is the homomorphism nZ ,\ensuremath{\to}Z. Thus both classifying spaces are circles, and the induced map on H1(\ensuremath{-}; k) is multiplication by n after choosing the usual generators. It is the identity on H0 and both sides have zero homology in higher degrees. If k is a nonzero commutative ring in which n is invertible, p is a k-homology equivalence. Regard it now as the map from X \ensuremath{-}\ensuremath{\to}S to S \ensuremath{-}\ensuremath{\to}1S. On strict fibres it induces

H0(\ensuremath{-}; k) : kn \ensuremath{-}\ensuremath{\to}k.

Z/n \ensuremath{-}\ensuremath{\to}\ensuremath{*}, The latter map adds coordinates and has nonzero kernel. Hence it is not a fibrewise homology equivalence, although both objects over S are perfect.

The same example disproves L6 itself, not only the fibre-detection formulation. Let E be the translation groupoid of Z acting on itself, and let E \ensuremath{-}\ensuremath{\to}BZ be the evident discrete fibration and opfibration. There is exactly one arrow between each ordered pair of objects of E, so E is contractible. The objects of X \ensuremath{\times}S E are pairs (a, z); simultaneous translation preserves a \ensuremath{-}z modulo n. Each of the resulting n components has exactly one arrow between any two of its objects and is contractible. Therefore

X \ensuremath{\times}S E \ensuremath{-}\ensuremath{\to}E

induces kn \ensuremath{-}\ensuremath{\to}k on H0 and fails to be a homology equivalence. The base-change map is perfect, as required by L6. The argument explains why replacing “homotopy fibration” by “perfect” cannot repair the homological version.

\end{dcomment}


\dgrblock{dgrb215}{An integral-coefficient version of the obstruction}
\begin{dcomment}{development}{p43}{An integral-coefficient version of the obstruction}
Grothendieck also suspects failure for integral homology, where a nontrivial finite covering of the circle no longer gives a homology equivalence. The following explicit construction settles that case. Its ingredients are standard: a semidirect product and the mapping-torus homology sequence. We checked the discussion of homological local constancy in Cisinski (2009, Scholium 1.26) and the mapping-torus calculation in Hatcher (2002, Example 2.48); we did not find this particular categorical application to the sentence here in those sources. The marking records the application and proof supplied by us, not a priority claim for the group or its homology.

Proposition (integral homology does not detect the fibre). \ours There is a perfect functor p : X \ensuremath{-}\ensuremath{\to}S which induces an isomorphism on integral homology in every degree, while its strict fibre has nonzero first integral homology. Consequently the localizer of integral homology equivalences fails L5.

Proof. Put A = Z[1/2], let h : A \ensuremath{-}\ensuremath{\to}A be multiplication by 2, and put G = A \ensuremath{\rtimes}h Z. Take X = BG, S = BZ and let p be induced by the quotient G \ensuremath{\twoheadrightarrow}Z. Since this homomorphism is surjective, every base arrow admits lifts with either prescribed source or prescribed target; all arrows in these groupoids are invertible and hence Cartesian and cocartesian. Thus p is both a Grothendieck fibration and an opfibration, and is perfect. Its strict fibre is BA.

We compute the fibre’s homology first. The group A is the filtered union of the cyclic subgroups 2\ensuremath{-}rZ. The nerve, its free simplicial chain complex and homology all commute with this filtered union: each finite tuple of group elements occurs at a finite stage, and filtered colimits of abelian groups are exact. Since BZ has the homotopy type of a circle, this gives

(i \ensuremath{\geq}2).

H0(BA; Z) = Z, H1(BA; Z) = A, Hi(BA; Z) = 0 The map induced by h acts by 2 on H1 and by the identity on H0.

The group automorphism h induces a homeomorphism of the realization of the nerve of A. Its mapping torus Th is a bundle over the circle. Its fundamental group is A \ensuremath{\rtimes}h Z after choosing the orientation of the circle, and its higher homotopy groups vanish by the homotopy sequence, since the fibre is K(A, 1). Hence Th is a classifying space for G, compatibly with the projection to BZ. The mapping-torus homology sequence has the terms

\textperiodcentered{} \textperiodcentered{} \textperiodcentered{} \ensuremath{-}\ensuremath{\to}Hi(BA)

\ensuremath{-}\ensuremath{-}\ensuremath{-}\ensuremath{\to}Hi(BA) \ensuremath{-}\ensuremath{\to}Hi(Th) \ensuremath{-}\ensuremath{\to}Hi\ensuremath{-}1(BA) \ensuremath{-}\ensuremath{\to}\textperiodcentered{} \textperiodcentered{} \textperiodcentered{} . In degree 1, 1 \ensuremath{-}h\ensuremath{*}= \ensuremath{-}1 on A is invertible; in degree 0 it is zero on Z. Exactness therefore gives Hi(Th) = 0 for i \ensuremath{\geq}2 and an isomorphism H1(Th) \ensuremath{-}\ensuremath{\to}Z. Naturality with respect to the map of fibres BA \ensuremath{-}\ensuremath{\to}\ensuremath{*}identifies this last isomorphism with the homology map induced by the projection to the circle. Both spaces are connected, so the projection is also an isomorphism on H0. This proves that p is an integral homology equivalence. Its map of fibres to the fibre of 1S is BA \ensuremath{-}\ensuremath{\to}\ensuremath{*}, which fails on H1 = A = 0. \ensuremath{\Box}

The example distinguishes absolute homology of the total category from homology after changing the base. The monodromy multiplies the fibre’s homology by 2, and 1 \ensuremath{-}2 is invertible; its contribution disappears from the total homology without disappearing from the fibre. Thus the obstruction is present even with integral coefficients, independently of inverting the degree of a finite covering.

\end{dcomment}


\textbf{L6 bis} \textit{Comme L6, mais en supposant seulement $Y' \longrightarrow Y$ une} Hot-\textit{fibration.}

\printedpage{28}
% head: A. Grothendieck : Dérivateurs, Ch. VI, Édition des Universités de Montpellier II et Paris VII

Il y a des chances que pratiquement, quand L6 (ou seulement L5) est verifié, L6 bis l'est aussi. % [sic] verifié -> vérifié, missing accent

On va introduire un axiome

\textbf{L7} \ a) \textit{Soit $(f_i : X_i \longrightarrow Y_i)_{i \in I}$ une famille de flèches dans} Cat\textit{, soit $f : \coprod X_i \longrightarrow \coprod Y_i$ la flèche correspondante sur les catégories somme. Pour que $f \in W$, il faut et il suffit que $f_i \in W$ pour tout $i \in I$.}

\hspace*{2em}b) \textit{Soit $f : X \longrightarrow Y$ dans $W$. Si $Y$ est} $0$\textit{-connexe, $X$ aussi.}

Cet axiome équivaut à la conjonction de La ($f \in W \Longrightarrow \pi_0 f$ bijectif) et de la validité du ‘Lemme’ page 23, que je remets ici pour mémoire :

\textbf{Lemme.} \textit{L'axiome L7 précédent sur le localiseur fondamental $W$ équivaut à la condition suivante : $f : X \longrightarrow Y$ est dans $W$ si et seulement si $\pi_0(f)$ est bijectif, et $f$ induit des $W$-équivalences $f_i : X_i \longrightarrow Y_i$ sur les composantes connexes.}

Il me faut revenir encore sur le 1) de la proposition, il n'est pas vrai que ‘c'est’ la condition L4 sur $W$ (que j'admets, ai-je dit), c'est L4 dans le cas particulier où $Y = S$. Réflexion faite, je le vois comme une variante renforcée de L4, qui ne semble pas devoir découler de L4 :

\textbf{L4 bis} \textit{Soit} \begin{tikzcd}[baseline=(s.base), row sep=small, column sep=small]
X \arrow[rr, "f"] \arrow[dr] & & Y \arrow[dl] \\
& |[alias=s]| S &
\end{tikzcd} \textit{dans} Cat\textit{, supposons que $f$ soit une $W$-équivalence locale sur $S$ (i.e. $f/s : X/s \longrightarrow Y/s$ une $W$-équivalence $\forall\ s \in S$), alors $f$ est une $W$-équivalence} (\footnote[29]{équivaut à l'axiome dual, en termes des $s\backslash f$}).

\msp{page 45}
Par la suite, par un ‘localiseur fondamental’, nous entendrons un ensemble
\[
W \subseteq \mathrm{Fl}\,\mathrm{Cat}
\]
satisfaisant (tout au moins) les axioms % [sic] axioms -> axiomes, French spelling
\[
\text{L1, L3 quater, L4 bis, L7,}
\]
i.e.

\dgrblock{dgrb216}{Componentwise locality and relative Theorem A}
\begin{dcomment}{exposition}{p44}{Componentwise locality and relative Theorem A}
Grothendieck separates two axiomatic matters here. The first concerns disjoint components; the second concerns descent over a base that need not be the target of the map being tested. Keeping them separate clarifies which part of the preceding argument each axiom repairs.

The component criterion is exactly equivalent to L7. To prove the nontrivial direction, let f : X \ensuremath{-}\ensuremath{\to}Y belong to W and decompose Y = ` j Yj into its nonempty connected components. Put Xj = f \ensuremath{-}1(Yj), a union of components of X. Then f = ` j fj, so L7(a) implies fj \ensuremath{\in}W for every j. L7(b) implies that every Xj is nonempty and connected. Hence \ensuremath{\pi}0(f) is bijective, and the fj are precisely the component maps. Conversely, if \ensuremath{\pi}0(f) is bijective and every component map belongs to W, L7(a) gives f \ensuremath{\in}W. The empty target creates no exception: a functor into the empty category has empty source.

To recover L7 from the criterion, consider a family (fi). The component map of ` i fi is just a component map of one fi; moreover its map on components is bijective exactly when each \ensuremath{\pi}0(fi) is bijective. The criterion therefore gives both directions of L7(a). If the target of a W-equivalence is connected, its component bijection makes the source connected, giving L7(b). This is a complete proof of the lemma, including the reflection direction in the coproduct assertion.

For the second matter, the original L4 applies to a map X \ensuremath{-}\ensuremath{\to}Y when the overcategories of its fibres over objects of Y are aspherical. In Proposition 1, however, one compares X/s and Y/s for an independent base S. These overcategories need not be aspherical; only the map between them is assumed to be a W-equivalence. Thus the printed sentence identifying the proposition directly with the original L4 does not establish it. L4 bis is exactly the required relative formulation. The dual condition uses s\textbackslash{}X \ensuremath{-}\ensuremath{\to}s\textbackslash{}Y .

The accepted modern definition of a fundamental localizer includes the relative formulation; see Cisinski (2004, \S{}1.1). Its minimal member is W\ensuremath{\infty}by that paper’s Theorem 2.2.11. This gives a precise later framework for the corrected argument. We make no claim here that the absolute and relative formulations are logically independent under every possible package of the earlier axioms; what the text has exhibited is the need for the relative assertion in this proof. Opposite-category invariance then transfers the relative assertion to its dual, using the category-of-simplices comparison discussed earlier in the chapter.

\end{dcomment}


\begin{itemize}
\item[a)] \textit{Saturation} de $W$ (conditions a) b) c) p. 1)

\item[b)] $W_\infty \subseteq W$ $[\subseteq W_0]$ (\footnote[30]{Je ne suis plus sûr s'il est vraiment utile d'imposer $W_\infty \subseteq W$. Prendre au lieu de ça l'axiome de l'objet final, soit b$'$). Si on veut a) b$'$) c) + L5, alors on aura $W \subseteq W_\infty$. \textbf{NB} a) b) c) + L5 implique $W = W_\infty$, cf. annotation page 43.}) (\footnote[31]{cf. d)}).

\item[c)] Si $f : X \longrightarrow Y$ sur $S$ est localement sur $S$ une $W$-équivalence, alors $f \in W$ (L4 bis)
\end{itemize}

% ---- from Der-ChVI-345_p031-035.tex
\printedpage{29}
% head: A. Grothendieck : Dérivateurs, Ch. VI, Édition des Universités de Montpellier II et Paris VII
% foot: Version 9 novembre 2001

\begin{enumerate}[leftmargin=*]
\item[d)] Pour que $f : X \longrightarrow Y$ soit dans $W$, il faut et il suffit que $\pi_0(f)$ soit bijectif, et que les $f_i : X_i \longrightarrow Y_i$ induits sur les composantes connexes soient dans $W$.
\end{enumerate}

Ces axiomes sont vérifiés par tous les $W$ qui je connais, et qu' on aurait envie de regarder comme localiseurs. Par contre, L5, L6, L6 bis (p. 43, 44) sont des axiomes de nature beaucoup plus spéciale. J'ai envie de dire qu'ils distinguent les localiseurs de nature proprement \textit{`homotopique'} de ceux qui sont de nature \textit{`homologique'} (ou `cohomologique'). Il faudrait que je m'assure qu'ils sont satisfaits pour les $W_i$ ($i$-équivalence d'homotopie) (\footnote[32]{C'est \textit{faux}, cf. annotation ci-dessus.}). % [sic] qui je connais -> que je connais, relative pronoun object requires "que"

Le fait que L4 bis, i.e. c) ci-dessus, implique l'axiome dual, résulte du fait que L4 bis implique ceci

\textbf{Proposition 2.} \textit{Soit $f : X \longrightarrow Y$ dans} $\Cat$\textit{, alors $f \in W$ si et seulement si $f^o : X^o \longrightarrow Y^o$ est dans $W$.}

Se voit par l'argument de I, page 75 (cor. 2) - n'utilise que L4, pas même L1, L3.

\msp{page 46}

\dgrblock{dgrb217}{Why the new fibration axiom selects full homotopy theory}
\begin{dcomment}{status}{p45}{Why the new fibration axiom selects full homotopy theory}
Grothendieck’s second L5 is much stronger than the first axiom with that label. The first passed from equivalences of fibres to an equivalence of totals. The new one attempts the reverse implication for homotopy fibrations. The annotations correctly recognize that this is incompatible with the proposed range of examples, but the displayed uniqueness claim needs the standing connectedness condition W \ensuremath{\subseteq}W0. Without it the maximal class of all functors satisfies L1, L3 quater, L4 bis and L5, so those four conditions alone do not force W = W\ensuremath{\infty}.

Here is the mechanism of the uniqueness assertion, retaining W\ensuremath{\infty}\ensuremath{\subseteq}W \ensuremath{\subseteq}W0, two-out-of-three, the new L5, and the path replacement already available in the chapter. Suppose first that T \ensuremath{-}\ensuremath{\to}\ensuremath{*} belongs to W. Its classifying space is connected by W \ensuremath{\subseteq}W0. For a chosen object t form the perfect replacement P = CT (\ensuremath{*}) \ensuremath{-}\ensuremath{\to}T of the inclusion of that object. Its total category is weakly contractible. Two-out-of-three implies P \ensuremath{-}\ensuremath{\to}T \ensuremath{\in}W. Applying L5 to this map between the homotopy fibrations P \ensuremath{-}\ensuremath{\to}T and T \ensuremath{-}\ensuremath{\to}1T shows that its fibre F \ensuremath{-}\ensuremath{\to}\ensuremath{*}belongs to W. The fibre models the based loop space of BT. Repeating the argument with F and successive fibres gives categorical models for all iterated loop spaces, each W-aspherical and hence connected. Therefore

\ensuremath{\pi}n(BT, t) = \ensuremath{\pi}0(\ensuremath{\Omega}nBT) = 0

(n \ensuremath{\geq}1). Together with connectedness this says that T is weakly contractible.

Now let f : X \ensuremath{-}\ensuremath{\to}Y belong to W. Replace it by the perfect map CY (X) \ensuremath{-}\ensuremath{\to}Y . Its total map is still in W, and L5 makes every strict fibre W-aspherical. What we have just proved makes those fibres weakly contractible. Since this is a homotopy fibration, its homotopy fibres are weakly contractible, so its total map is a weak homotopy equivalence; the replacement X \ensuremath{-}\ensuremath{\to}CY (X) is also one. Thus f \ensuremath{\in}W\ensuremath{\infty}. This proves W = W\ensuremath{\infty}. The argument uses the previously established homotopy-fibre property of the path replacement, not a new assertion that arbitrary strict fibres compute homotopy fibres. The later abstract characterization in Cisinski (2004, Theorem 2.3.6) likewise includes preservation of connected components.

One can see the failure for every finite truncation directly. Let W\ensuremath{\leq}n mean maps inducing isomorphisms on \ensuremath{\pi}i for 0 \ensuremath{\leq}i \ensuremath{\leq}n, with all basepoints understood. Choose a nonzero abelian group A and a categorical model S of K(A, n + 1). The perfect replacement P = CS(\ensuremath{*}) \ensuremath{-}\ensuremath{\to}S has contractible total space, so belongs to W\ensuremath{\leq}n, while its fibre models K(A, n) and is not W\ensuremath{\leq}n-aspherical. For n = 0 that fibre is the nontrivial discrete group A. This disproves L5. Pulling P \ensuremath{-}\ensuremath{\to}S back along itself yields a total space weakly equivalent to \ensuremath{\Omega}BS, because the base P is contractible and the map is a homotopy fibration. The projection P \ensuremath{\times}S P \ensuremath{-}\ensuremath{\to}P therefore fails to belong to W\ensuremath{\leq}n, disproving L6 even for this perfect base change. The finite truncations are homotopical theories, but they do not satisfy the unrestricted fibre-detection requirement of this L5.

\end{dcomment}


Je vais essayer de résumer ce que je crois avoir compris :

\begin{enumerate}[leftmargin=*]
\item[1)] La $W$-équivalence `par fibres' n'a d'intérêt (pour construire des catégories localisées) que dans le cas où on sait qu'elle implique déjà, soit la $W$-équivalence locale, soit la $W$-équivalence colocale (soit même la $W$-équivalence universelle) sur $S$, donc quand on suppose $X$, $Y$ pour le moins tous deux soit $W$-lisses, soit $W$-propres, soit $W$-fibrés sur $S$.

\item[2)] On pourrait appeler la $W$-équivalence locale sur $S$ la \textit{`$W$-équivalence fine à gauche'} sur $S$, et la $W$-équivalence colocale sur $S$ la \textit{`$W$-équivalence fine à droite'} sur $S$. Elles coïncident dans $\uc{Fibhot}\,S$, a fortiori dans $\uc{Parf}\,S$ et dans les quatre catégories (cf. diagramme p. 8) $\uc{Fib}_0S$, $\uc{Fibsc}_0S$, $\uc{Cofib}_0S$, $\uc{Cofibsc}_0S$ (\footnote[33]{On pourrait les appeler les six catégories \textit{internes} du diagramme.}). Par contre, elles ne coïncident pas sur les catégories \textit{bordantes} du diagramme, notamment sur $\uc{Liss}\,S$, $\uc{Prop}\,S$, $\uc{Fib}\,S$, $\uc{Cofib}\,S$, $\uc{Fibsc}\,S$, $\uc{Cofibsc}\,S$. À cet égard, c'est la \textit{$W$-équivalence fine} \textbf{droite} (en termes des $s\backslash X$) qui est intéressante sur le côté (hélas) gauche du diagramme, i.e. pour $\uc{Liss}\,S$, $\uc{Fib}\,S$, $\uc{Fibsc}\,S$, alors que c'est la \textit{$W$-équivalence fine} \textbf{gauche} (en termes des $X/s$) qui est pertinente du côté droit, pour $\uc{Prop}\,S$, $\uc{Cofib}\,S$, $\uc{Cofibsc}\,S$. On n'aura jamais à regarder simultanément les deux équivalences fines sur une même catégorie du diagramme, à la seule exception de $\Cat/S$ tout entier.
\end{enumerate}

\msp{page 47}

\begin{enumerate}[leftmargin=*]
\item[3)] La $W$-équivalence grossière, elle, est `bonne partout' - mais ça dépend pour qui. Pour les 6 catégories qui sont au dessus de $\uc{Fibhot}\,S$, i.e.
\[
\uc{Fibhot}\,S,\quad \uc{Parf}\,S,\quad \uc{Fib}_0S,\quad \uc{Fibsc}_0S,\quad \uc{Cofib}_0S,\quad \uc{Cofibsc}_0S,
\]
\printedpage{30}
% head: A. Grothendieck : Dérivateurs, Ch. VI, Édition des Universités de Montpellier II et Paris VII
% foot: Version 9 novembre 2001
elle coïncide avec toutes les autres équivalences (moyennnant l'axiome L5), pour les trois catégories bordantes supérieures gauches % [sic] moyennnant -> moyennant, triple n
\[
\uc{Liss}\,S,\quad \uc{Fib}\,S,\quad \uc{Fibsc}\,S
\]
du diagramme, elle coïncide avec la $W$-équivalence fine droite, et dualement sur
\[
\uc{Prop}\,S,\quad \uc{Cofib}\,S,\quad \uc{Cofibsc}\,S
\]
elle coïncide avec la $W$-équivalence fine gauche. Quant à $\Cat/S$ tout entier, c'est la seule des catégories du diagramme où il y ait \textit{trois} notions intéressantes différentes de $W$-équivalence, à savoir : la grossière, et les deux fines (à gauche, à droite). Sur les autres six catégories bordantes du diagramme, il y en a en principe deux : la grossière, et une spécifiée des deux fines (défendu de regarder l'autre - à ses risques et périls \ldots) : la droite du côté gauche, la gauche du côté droit. (NB Je devrais refaire le diagramme, décidément, en mettant les fibrations à droite, pas à gauche.) Quant aux six catégories internes au diagramme, dans elles il n'y a plus qu'une seule notion de $W$-équivalence (puisque sur elles les quatre notions coïncident).
\end{enumerate}

\msp{page 48 (vide)}

\msp{page 49} (\textbf{NB} il n'y a pas de page 48)

\dgrblock{dgrb218}{Reading the corrected classification}
\begin{dcomment}{exposition}{p47}{Reading the corrected classification}
Grothendieck is now trying to separate the homotopy theory of a diagram from the homotopy type of its total category. This separation is the reason for retaining three localizations. Write WS for equivalences of total categories, W d S for the condition that every s\textbackslash{}X \ensuremath{\to}s\textbackslash{}Y belongs to W, and W g S for the analogous condition on X/s \ensuremath{\to}Y/s. The superscripts record the two comma directions, not the placement of a category in the drawing. One sentence in the present paragraph must be corrected mathematically. The assertion that WS coincides with W d S on the smooth or fibred side, and dually with W g S on the proper or cofibred side, is false. The edition prints it without a negation. On the fibred side it is fibrewise equivalence that coincides with W d S; total-category equivalence remains a weaker condition. The walking-arrow example in the commentary at manuscript p. 32 proves the distinction: \{0\} ,\ensuremath{\to}[1], viewed over [1], is an equivalence of total homotopy types, but its fibres over 1 are respectively empty and nonempty. The same paragraph subsequently says that there are two notions on each border, and p. 49 uses the correct distinction.

On the six internal categories the equality with WS uses the new L5 of p. 43. This is different from the earlier axiom carrying that name. Under L5, a map between homotopy fibrations which is an equivalence on total categories is an equivalence on fibres. The ordinary replacement arguments compare all the coarse localizations without this additional axiom, and comma resolutions independently compare the fine theories. L5 is needed when we identify the coarse class on internal objects with their fine class, and thereby realize the coarse theory as a fully faithful locally constant part of the fine theory. This distinction prevents the corrected theorem from inheriting the false equality in its provisional summary.

\end{dcomment}


Je résume maintenant ce que je crois avoir compris sur les catégories localisées correspondantes. Il y en a, à équivalence près, \textit{trois} différentes seulement en tout et pour tout, quand on prend les localiseurs `admissibles' dans les catégories, comme je viens de le dire (c'étant ça en somme l'intérêt objectif des prescriptions!) De façon précise :

\begin{enumerate}[leftmargin=*]
\item[$1^o$)] La \textit{$W$-équivalence grossière} donne `la même' catégorie localisée pour les 13 catégories du diagramme, i.e. toutes les flèches de transition du diagramme induisent des équivalences pour cette $W$-équivalence grossière. Le localisé sera noté \fbox{$\HOT^{\mathrm{lc}}_W(S)$}, où l'exposant lc signifie: ($W$-types d'homotopie relatif sur $S$) \textit{`localement constants.'} (Il ne me reste qu'un petit doute pour $\uc{Fibsc}\,S$ et $\uc{Cofibsc(S)}$ (\footnote[34]{Ce doute n'est pas fondé, tout marche parfaitement!}), si l'équivalence grossière y est bel et bien raisonnable. On verra en faisant la démonstration.)

\item[$2^o$)] La \textit{$W$-équivalence fine à droite} donne le même résultat pour toutes les catégories du diagramme où elle est considérée comme admise, \textit{et où elle ne coïncide pas avec la $W$-équivalence grossière} déjà considérée, donc sur les quatre catégories bordantes gauches
\[
\Cat/S,\quad \uc{Liss}\,S,\quad \uc{Fib}\,S,\quad \uc{Fibsc}\,S
\]
(\footnote[35]{\textbf{NB} Sur les trois dernières cette équivalence est la $W$-équivalence par fibres.}). Le localisé est noté \fbox{$\HOT_W(S)$}.
\end{enumerate}

\msp{page 50}

\printedpage{31}
% head: A. Grothendieck : Dérivateurs, Ch. VI, Édition des Universités de Montpellier II et Paris VII
% foot: Version 9 novembre 2001

\begin{enumerate}[leftmargin=*]
\item[$2'^o$)] Dualement, la \textit{$W$-équivalence fine gauche} donne le même résultat sur toutes les catégories du diagramme où elle est admise, et où elle ne coïncide pas avec la grossière, i.e. sur les 4 catégories
\[
\Cat/S,\quad \uc{Prop}\,S,\quad \uc{Cofib}\,S,\quad \uc{Cofibsc}\,S.
\]
Il y a lieu de noter la catégorie des fractions commune par le sigle
\[
\HOT_W(S^o)^o \qquad \text{(ou aussi $\HOT^o_W(S)$)},
\]
puisqu'elle est isomorphe à celle notée par ce sigle.
\end{enumerate}

Si on veut des notations sans ambiguïté, désignant des catégories bien définies, et non pas seulement `à équivalence près', je choisirai
\[
\HOT_W(S) \defeq \uc{Fibsc}(S)(W^{\mathrm{d}}_S)^{-1}
\]
d'où
\[
\HOT_W(S^o) \qquad \underbrace{\simeq}_{\text{\footnotesize isomorphisme de catégories}} \qquad \uc{Cofibsc}(S)(W^{\mathrm{g}}_S)^{-1},
\]
où $W^{\mathrm{d}}_S$, $W^{\mathrm{g}}_S$ désignent les ensembles de flèches correspondants à la $W$-équivalence fine droite resp. gauche. Pour $\HOT^{\mathrm{lc}}_W(S)$, je préfère une définition aussi symétrique que possible en $S$, $S^o$, et de plus une catégorie aussi petite que possible, ce qui me conduit à choisir
\[
\HOT^{\mathrm{lc}}_W(S) \defeq (\uc{Parf}\,S)W^{-1}_S,
\]
où on a noté par $W_S$ les $W$-équivalences dans $\uc{Parf}\,S$. (Inutile cette fois de préciser lesquelles - toutes les quatre \texttt{[notions]} coïncident dans $\uc{Parf}\,S$ !)

\msp{page 51}

\dgrblock{dgrb219}{The covariant theory and the extra opposite}
\begin{dcomment}{exposition}{p50}{The covariant theory and the extra opposite}
Grothendieck wants a notation which makes the symmetry between fibrations and opfibrations immediate. There is a genuine symmetry, but the edition alternates between two incompatible formulas. With the definition fixed on manuscript p. 3, the precise identities are

[Cofibsc(S)][(W g

HOTW (S) = [Sop, Cat][W \ensuremath{-}1 S)\ensuremath{-}1] \ensuremath{\simeq}[S, Cat][W \ensuremath{-}1 obj] = HOTW (Sop). obj], Here Wobj means objectwise membership in W. There is no opposite on the value category in the last expression. The displayed isomorphism in the middle of p. 50 has this correct form; the preceding prose and the editorial correction at p. 51 insert an additional opposite. We retain those readings as part of the source, while using the comma-category models to determine the variance of every assertion.

Indeed, taking opposite total categories sends a fibration over Sop to an opfibration over S. It sends a functor X \ensuremath{\to}Y to a functor Xop \ensuremath{\to}Y op in the same direction. Thus it defines a covariant isomorphism of the categories of structured objects; it does not reverse their morphisms. The same observation justifies the correctly stated symmetry of the perfect-fibration localization at p. 51.

The extra opposite cannot be repaired by an unspecified equivalence. Take S = e and W = W\ensuremath{\infty}. The category of homotopy types is not equivalent to its opposite. Its initial object \ensuremath{\varnothing}is strict: a morphism T \ensuremath{\to}\ensuremath{\varnothing}exists only when T \ensuremath{\simeq}\ensuremath{\varnothing}. Such an equivalence would force the initial object of the opposite category, namely the one-point homotopy type, to be strict there. This would say that every map \ensuremath{*}\ensuremath{\to}T is an isomorphism, contradicted by the inclusion of one point into a two-point set. The external opposite is therefore a substantive variance slip, even though both fine localizations themselves are correctly constructed.

The notation HOTo W (S) may still be used as a label for the covariant theory, provided it is defined by the left-fine localization. It should not be interpreted here as taking the opposite of HOTW (Sop). This convention also gives the stated direction to the functor j in the diagram which follows.

\end{dcomment}


Notons l'isomorphisme
\[
\uc{Parf}(S^o) \simeq \uc{Parf}\,S
\]
(provenant du fait que $f : X \longrightarrow Y$ dans $\Cat$ est une fibration parfaite si et seulement si $f^o$ l'est). Cet isomorphisme $X \longmapsto X^o$ est compatible avec les localiseurs (car $f \in W \Longleftrightarrow f^o \in W$), d'où
\[
\HOT^{\mathrm{lc}}_W(S^o) \simeq \HOT^{\mathrm{lc}}_W(S)
\]
(isomorphisme de catégories). On aura deux foncteurs canoniques entre les trois catégories de $W$-types d'homotopie relatifs sur $S$
\[
\begin{tikzcd}[column sep=small]
(\uc{Fibsc}_0S)W_S^{-1} \arrow[rr, "\approx"] \arrow[dr, dotted, "i'"'] & & \HOT^{\mathrm{lc}}_W(S) \arrow[dl, "i"] \arrow[dr, "j"'] & & (\uc{Cofibsc}_0S)W_S^{-1} \arrow[ll, "\approx"'] \arrow[dl, dotted, "j'"] \\
& \HOT_W(S) & & \HOT^o_W(S) \defeq \HOT_W(S^o). &
\end{tikzcd}
\]
\printedpage{32}
% head: A. Grothendieck : Dérivateurs, Ch. VI, Édition des Universités de Montpellier II et Paris VII
% foot: Version 9 novembre 2001
\texttt{[plutôt} $\defeq \HOT_W(S^o)^o$\texttt{]} Ces flèches $i$, $j$ sont suggerées par les flèches en pointillés $i'$, $j'$, provenant des inclusions % [sic] suggerées -> suggérées, missing accent
\[
\uc{Fibsc}_0S \lhook\joinrel\longrightarrow \uc{Fibsc}\,S,\quad \uc{Cofibsc}_0S \lhook\joinrel\longrightarrow \uc{Cofibsc}\,S,
\]
qui définiraient $i$, $j$ à l'aide du choix de foncteurs quasi-inverses pour les flèches horizontales. Mais $i$, $j$ ne seraient définis qu'à isomorphisme (unique) près. Il sera plus joli de choisir les quasi-inverses de façon canonique, à l'aide des flèches
\[
\uc{Fibsc}_0S \xleftarrow{\Phi_0} \uc{Parf}\,S \xrightarrow{\Psi_0} \uc{Cofibsc}_0S
\]
(où $\Phi_0 = \Phi|\uc{Parf}\,S$, $\Psi_0 = \Psi|\uc{Parf}\,S$) du diagramme p. 8. La question principale
\msp{page 52}
qui reste non élucidée, c'est si ces foncteurs $i$, $j$ sont bien \textit{pleinement fidèles,} comme je m'y attends certes.

Je veux maintenant prouver les affirmations précédentes sur les localisés.

\textcircled{1} Le localisé pour la $W$-équivalence grossière. Il suffira, à un $\varepsilon$ près, d'établir les conditions du lemme p. 11, dans les trois cas suivants.
\begin{tikzcd}[baseline=(\tikzcdmatrixname-1-1.base)]
\uc{Parf}\,S \arrow[r, "r"] & \Cat/S \arrow[l, dotted, bend left, "C"]
\end{tikzcd}
(ça donnera le cas de $\uc{Fibhot}\,S$, $\uc{Liss}\,S$ et $\uc{Prop}\,S$ sur $\Cat/S$) et
\begin{tikzcd}[baseline=(\tikzcdmatrixname-1-1.base)]
\uc{Fibsc}_0S \arrow[r, "p"] & \Cat/S \arrow[l, dotted, bend left]
\end{tikzcd},
\begin{tikzcd}[baseline=(\tikzcdmatrixname-1-1.base)]
\uc{Cofibsc}_0S \arrow[r, "q"] & \Cat/S \arrow[l, dotted, bend left]
\end{tikzcd},
qui donneront respectivement les cas de
\[
\uc{Fib}_0S,\ \uc{Fib}\,S,\ \uc{Liss}\,S \qquad \text{et} \qquad \uc{Cofib}_0S,\ \uc{Cofib}\,S,\ \uc{Prop}\,S \qquad \text{sur } S.
\]
Mais l'ennui pour $p$ et $q$, c'est que les foncteurs en sens inverse que j'avais construit (et noté encore $\Phi_0$, $\Psi_0$ (mais c'étaient essentiellement $\Phi\circ r\circ C$ et $\Psi\circ r\circ c$, à ne pas confondre avec les foncteurs $\Phi_0$, $\Psi_0$ en bas de la page précédente, et marqués aussi dans le diagramme p. 8), notons les plutôt $\Phi'$, $\Psi'$, marchent bien du côté $\Cat/S$, i.e. $p\Phi'$, $q\Psi'$ reçoivent le foncteur $\mathrm{id}_{\Cat/S}$ par une $W$-flèche, mais pas du côté des catégories but - on n'arrive pas à trouver une flèche naturelle entre $\Phi'p$ resp. $\Psi'q$, et les foncteurs identiques de $\uc{Fibsc}_0S$ resp. $\uc{Cofibsc}_0S$. Je ne vois
\msp{page 53}
pas, en effet, qu'il y en ait - la difficulté, par cette voie, parait irrémédiable. Par contre, si on regarde $\Phi'$, $\Psi'$ comme foncteurs à valeurs dans $\uc{Fib}_0S$, $\uc{Cofib}_0S$ respectivement (donc dans les catégories où on n'exige pas que les flèches soient catésiennes resp. cocartésiennes), on trouve encore des $W$-morphismes % [sic] catésiennes -> cartésiennes, missing r
\[
\mathrm{id} \longrightarrow \Phi'p, \qquad \mathrm{id} \longrightarrow \psi'q, % [sic] \psi'q -> \Psi'q, lowercase psi where \Psi' is used everywhere else
\]
ce qui liquide au moins déjà le cas de toutes les catégories des diagrammes p. 8 qui sont des
\printedpage{33}
% head: A. Grothendieck : Dérivateurs, Ch. VI, Édition des Universités de Montpellier II et Paris VII
% foot: Version 9 novembre 2001
sous-catégories pleines de $\Cat/S$ (\footnote[36]{Mais c'est idiot de prendre des foncteurs si compliqués, parce qu'on veut tout ramener à $\Cat/S$. Il est de plus simple de comparer $\uc{Fibhot}\,S$ directement à $\uc{Fibsc}_0S$ et à $\uc{Cofibsc}_0S$. $\Phi_0$ et $\Psi_0$ (qui sont en fait définis sur $\uc{Fibhot}$, pas seulement sur $\uc{Parf}$), et on trouve comme ci-dessous $\mathrm{id} \xrightarrow{W} \pi_0\Phi_0$ et $\Phi_0\pi_0 \xrightarrow{W} \mathrm{id}$ (itou pour $\Psi_0$ et $\uc{Cofibsc}_0S$), où $\pi_0 : \uc{Fibsc}_0S \longrightarrow \uc{Fibhot}\,S$.}). Il ne reste à traiter que le cas des quatre catégories de structures `scindées'
\[
\uc{Fibsc}\,S,\quad \uc{Fibsc}_0S \qquad \text{et} \qquad \uc{Cofibsc}\,S,\quad \uc{Cofibsc}_0S.
\]
Bien sûr il suffit de le faire d'un côté, disons du côté des fibrés, par dualité ($S \longmapsto S^o$) cela établira le résultat symétrique. On va donc regarder d'un peu plus près le carré de foncteurs
\[
\begin{tikzcd}[column sep=huge, row sep=large]
\uc{Fibsc}\,S \arrow[d, "p"] & \uc{Fibsc}_0S \arrow[l, hook, "\alpha'"'] \arrow[d, "p_0"'] \\
\uc{Fib}\,S \arrow[u, dotted, bend left, "\varphi"] & \uc{Fib}_0S \arrow[l, hook, "\alpha"'] \arrow[u, dotted, bend right, "\varphi_0"']
\end{tikzcd}
\]
(\footnote[37]{NB Il s'agit de prouver seulement que $p$, $p_0$ induisent des équivalences pour les localisés (pour la $W$-équivalence grossière). On sait déjà qu'il en est ainsi pour $\alpha$, et ça en résultera pour $\alpha'$.}, \footnote[38]{inutile de regarder ce carré, cf. annotations marginales ci-dessus (haut de la page)}). Le foncteur
\[
\Phi : \Cat/S \longrightarrow \uc{Fibsc}\,S
\]
induit des foncteurs en sens inverse de $p$, $p_0$ sur les sous-catégories pleines $\uc{Fib}\,S$, $\uc{Fib}_0S$ de $\Cat/S$, foncteurs marqués en pointillés comme $\varphi$, $\varphi_0$.

\msp{page 54}

\dgrblock{dgrb220}{Why the free split fibration is the right comparison}
\begin{dcomment}{context}{p53}{Why the free split fibration is the right comparison}
The difficulty with the first attempt is a difficulty about morphisms, not about the existence of enough good objects. A path completion puts every object into a particularly well-behaved class, but its unit need not preserve a chosen splitting. The comma construction \ensuremath{\Phi} has a stronger universal property: it freely supplies the selected Cartesian lifts. That property gives the reverse comparison map with exactly the morphism convention required by Fibsc(S).

This construction was already part of the theory of fibred categories before the present manuscript. Gray’s treatment constructs the free split fibration, and the modern formulation identifies split fibrations and splitting-preserving functors with strict algebras for the comma-category 2-monad on Cat/S (Gray, 1966; Emmenegger et al., 2024). In the latter reference, Theorem 2.6 separates strict algebras from pseudo-algebras; Theorem 5.4 treats the two adjoints used to split an already fibred category. These are distinct operations. The present \ensuremath{\Phi} is the free construction on an arbitrary category over S.

Its importance for localization is elementary. Let U : Fibsc(S) \ensuremath{\to}Cat/S forget the splitting. There are transformations

\ensuremath{\eta} : id \ensuremath{-}\ensuremath{\to}U\ensuremath{\Phi},

\ensuremath{\gamma} : \ensuremath{\Phi}U \ensuremath{-}\ensuremath{\to}id. The first lives in Cat/S; the second lives in the category with splitting-preserving morphisms. Both become isomorphisms after the relevant localization. Consequently U and \ensuremath{\Phi} induce inverse equivalences. One does not need a comparison transformation in a category with an incorrectly enlarged or restricted class of arrows.

For coarse equivalences, this conclusion requires only that W contain categorical homotopy equivalences and satisfy two-out-of-three. The ordinary functor forgetting the initial base object gives U\ensuremath{\Phi}(X) \ensuremath{\to}X, with a section \ensuremath{\eta}X and a natural contraction. Although the forgetful map is usually not over S, the contraction proves that \ensuremath{\eta}X is an equivalence of total homotopy types. The counit is then a coarse equivalence because its composite with the unit on UY is the identity. The more informative fibrewise verification comes on pp. 58--60.

\end{dcomment}


Il faut seulement vérifier que si $X/S$ est dans $\uc{Fib}_0S$, alors $\varphi(X) = \Phi(X)$ est dans la sous-catégorie pleine $\uc{Fibsc}_0S$ de $\uc{Fibsc}\,S$. Mais on a vu que c'est le cas, pour peu seulement que $X/S$ soit dans $\uc{Fibhot}\,S$, car les morphismes de changement de base dans $\Phi(X)$, associés aux flèches $u : s \longrightarrow t$ dans $S$
\[
u^*_{\Phi_X} : t\backslash X \longrightarrow s\backslash X
\]
sont des $W_\infty$-équivalences si $X$ est Hot-fibré sur $S$. D'autre part, pour $X$ quelconque dans $\Cat/S$, on a une flèche, fonctorielle en $X$,
\[
X \longrightarrow \Phi(X),
\]
i.e. $\mathrm{id}_{\Cat/S} \longrightarrow \pi\Phi$ où $\pi : \uc{Fibsc}\,S \longrightarrow \Cat/S$ est le foncteur canonique, qui est une $W_\infty$-équivalence (même un homotopisme), donc une $W$-équivalence. A fortiori les foncteurs $p\varphi$, $p_0\varphi_0$ donnent lieu ($\uc{Fib}\,S$ et $\uc{Fib}_0S$ étant des sous-catégories pleines de $\Cat/S$) à des flèches
\[
\mathrm{id} \longrightarrow p\varphi, \qquad \mathrm{id} \longrightarrow p_0\varphi_0.
\]

% ---- from Der-ChVI-345_p036-040.tex
\printedpage{34}
% head: A. Grothendieck : Dérivateurs, Ch. VI, Édition des Universités de Montpellier II et Paris VII

Il reste à trouver une flèche (et $W$-équivalence grossière)
\[\varphi p \longrightarrow \mathrm{id},\]
\msp{page 55}
laquelle induira automatiquement
\[\varphi_0 p_0 \longrightarrow \mathrm{id}\]
puisque $\underline{\mathrm{Fibsc}}_0 S$ est \texttt{[une]} sous-catégorie pleine de $\ul{Fibsc}\,S$. En fait, on va définir $\varphi p \longrightarrow \mathrm{id}$, donc
\[\boxed{\gamma_X : \Phi(X) \longrightarrow X}\]
quand $X$ est fibrée scindée sur $X$ - une flèche \textit{cartésienne,} donc (si on veut) compatible avec les scindages de $\Phi(X)$, et de $X$. En fait, on aura une telle flèche, fonctorielle en $X$, pour $X$ variable dans $\ul{Fib}\,S$ (donc sans se borner à des $S$-morphismes cartésiens entre catégories fibrées sur $S$), mais la flèche $\longrightarrow$ étant bien sûr \textit{cartésienne.} (Ça a un sens, sans choisir de scindage sur $X$, et sans qu'il en existe forcément.) La définition de cette flèche $\gamma$ fait appel de façon essentielle à l'hypothèse que $X$ soit Cat-fibrant sur $S$. Bien sûr, composé à la flèche d'inclusion évidente % [sic] fibrée scindée sur X -> fibrée scindée sur S, the base category is S
\[X \xrightarrow{i_X} \Phi(X)\]
on va avoir
\[\gamma_X i_X = \mathrm{id}_X\]
(\footnote[39]{NB Comme on sait que $i_X \in W$, il en résultera bien $\gamma_X \in W$ comme on le veut.}). On définit $\gamma$ fibre par fibre,
\[\gamma_{X,s} : \Phi(X)_s \overset{\text{déf}}{=} s\backslash X \longrightarrow X_s,\]
donc par une \textit{rétraction} de $s\backslash X$ sur la sous-catégorie pleine $X_s$. Un objet de $s\backslash X$ est un couple $(x,u)$, $x \in \mathrm{Ob}\, X$, $u : s \longrightarrow p_X(x)$ ($p_X : X \longrightarrow S$ le morphisme structural de $X$). On pose
\[\gamma(x,u) = u_X^*(x) \in \mathrm{Ob}\, X_s,\]
\msp{page 56}
ce qui suppose donc qu'on a choisi des foncteurs image inverse $u_X^* : X_t \longrightarrow X_s$, pour la catégorie Cat-fibrée $X/S$. (Mais à isomorphisme unique près, $\gamma$ ne dépend pas de ces choix.) Je me dispense de la vérification fastidieuse que cette loi est bien ‘fonctorielle en $s$’, i.e. définit bien une flèche cartésienne
\[\gamma_X : \Phi(X) \longrightarrow X,\]
et de vérifier que celle-ci est fonctorielle en $X$.

\printedpage{35}
% head: A. Grothendieck : Dérivateurs, Ch. VI, Édition des Universités de Montpellier II et Paris VII

Peut-être est-il plus élégant de montrer que $\pi : \ul{Fibsc}\,S \longrightarrow \mathrm{Cat}/S$ et $\Phi : \mathrm{Cat}/S \longrightarrow \ul{Fibsc}\,S$ sont adjoints l'un de l'autre :
\[\boxed{\mathrm{Homcart}_S(\Phi(\underbrace{\;X\;}_{\in\mathrm{Cat}/S}),\underbrace{\;Y\;}_{\in\ul{Fibsc}}) \simeq \mathrm{Hom}_S(X,\pi Y),}\]
d'où résulteront des flèches d'adjonction
\[\Phi\pi \longrightarrow \mathrm{id}_{\ul{Fibsc}}, \qquad\qquad \pi\Phi \longleftarrow \mathrm{id}_{\ul{Fib}},\]
\texttt{[plutôt} $\pi\Phi \longleftarrow \mathrm{id}_{\mathrm{Cat}/S}$\texttt{]} et on aura gagné. Mais pour prouver cette formule d'adjonction, j'ai l'impression qu'il faut utiliser $\varphi p \longrightarrow \mathrm{id}$ \ldots\ Il y a des vérifications fastidieuses \ldots

Finalement, je vois maintenant qu'on a cinq cas à traiter, et trois si on joue sur les symétries :

\msp{page 57}

\dgrblock{dgrb221}{The free split adjunction, with its arrows specified}
\begin{dcomment}{exposition}{p56}{The free split adjunction, with its arrows specified}
Grothendieck’s proposed adjunction is the efficient way to finish the comparison, once “Cartesian” in the Hom-set on the left is read as preserving the selected Cartesian lifts. We give its construction because the same calculation resolves the naturality issue in the preceding paragraph.

For p : X \ensuremath{\to}S, an object of \ensuremath{\Phi}X is (s, x, u) with u : s \ensuremath{\to}p(x). A morphism

(s, x, u) \ensuremath{-}\ensuremath{\to}(t, y, v)

is a pair (a, f) with a : s \ensuremath{\to}t, f : x \ensuremath{\to}y, and p(f)u = va. The chosen lift along a with target (t, y, v) is

\ensuremath{-}\ensuremath{-}\ensuremath{-}\ensuremath{-}\ensuremath{\to}(t, y, v).

(s, y, va) Identity and composition laws for these chosen lifts hold strictly. The unit sends x to (p(x), x, id). Let q : Y \ensuremath{\to}S be split, and let h : X \ensuremath{\to}UY be any functor over S. Write cu,z : u\ensuremath{*}z \ensuremath{\to}z for the chosen lift. Define

bh(s, x, u) = u\ensuremath{*}h(x).

For (a, f) as above, define bh(a, f) as the unique arrow over a such that

cv,h(y) bh(a, f) = h(f)cu,h(x).

The Cartesian property of cv,h(y) gives existence and uniqueness. These same uniqueness statements prove the identity and composition laws: both possible images of a composite have the required base arrow and the same composite with the terminal chosen lift. The strict splitting laws show that a selected lift (a, id) is taken to the selected lift in Y . Thus bh is splitting-preserving.

Conversely, every splitting-preserving H : \ensuremath{\Phi}X \ensuremath{\to}Y is forced by h = H\ensuremath{\eta}X. The chosen arrow (s, x, u) \ensuremath{\to}(p(x), x, id) forces its object value to be u\ensuremath{*}h(x), and Cartesian uniqueness forces its value on every morphism. The two constructions are inverse and natural, proving the actual adjunction

Homsplit/S(\ensuremath{\Phi}X, Y ) \ensuremath{\sim}= HomS(X, UY ).

Its counit is \ensuremath{\gamma}Y (s, y, u) = u\ensuremath{*}y, and \ensuremath{\gamma}Y \ensuremath{\eta}UY = idUY . Its naturality is strict for splitting-preserving functors, which are precisely the morphisms needed on this side of the comparison. This is the free-algebra adjunction underlying the 2-monad description cited above.

For a merely cloven fibration the same formula still gives a Cartesian functor \ensuremath{\Phi}X \ensuremath{\to}X, using the coherence isomorphisms of the cleavage. Different choices yield canonically isomorphic functors compatible with their Cartesian comparison arrows. But an arbitrary functor over S need not preserve Cartesian arrows, so the formula is not strictly natural on the full category Fib(S). Nor can “Cartesian” alone imply preservation of the selected splitting. Those stronger assertions in the source are unnecessary; the strict adjunction just proved gives the comparison in Fibsc(S), and the next block explains what survives after localization for all functors over S.

\end{dcomment}


\[
\begin{tikzcd}[row sep=large]
\ul{Parf}\,S \arrow[d, "r"'] \\
\mathrm{Cat}/S \arrow[u, dotted, bend right, "C"']
\end{tikzcd}
\qquad\qquad\qquad
\left\{\begin{array}{lcl}
rC & \longleftarrow & \mathrm{id} \\
Cr & \longleftarrow & \mathrm{id}
\end{array}\right.
\]

\[
\begin{tikzcd}
\ul{Fibsc}\,S \arrow[dr, "\lambda"] & & \ul{Cofibsc}\,S \arrow[dl, "\mu"'] \\
& \mathrm{Cat}/S \arrow[ul, dotted, bend left, "\Phi"] \arrow[ur, dotted, bend right, "\Psi"'] &
\end{tikzcd}
\qquad
\begin{array}{l}
\left\{\begin{array}{lcl}
\lambda\Phi & \longleftarrow & \mathrm{id} \\
\Phi\lambda & \longrightarrow & \mathrm{id}
\end{array}\right. \\[2ex]
\left\{\begin{array}{lcl}
\mu\Psi & \longleftarrow & \mathrm{id} \\
\Psi\mu & \longrightarrow & \mathrm{id}
\end{array}\right.
\end{array}
\]

\[
\begin{tikzcd}
\underline{\mathrm{Fibsc}}_0 S \arrow[dr, "\lambda_0"] & & \underline{\mathrm{Cofibsc}}_0 S \arrow[dl, "\mu_0"'] \\
& \ul{Fibhot}\,S \arrow[ul, dotted, bend left, "\Phi_0"] \arrow[ur, dotted, bend right, "\Psi_0"'] &
\end{tikzcd}
\qquad
\begin{array}{l}
\left\{\begin{array}{lcl}
\lambda_0\Phi_0 & \longleftarrow & \mathrm{id} \\
\Phi_0\lambda_0 & \longrightarrow & \mathrm{id}
\end{array}\right. \\[2ex]
\left\{\begin{array}{lcl}
\mu_0\Psi_0 & \longleftarrow & \mathrm{id} \\
\Psi_0\mu_0 & \longrightarrow & \mathrm{id}
\end{array}\right.
\end{array}
\]

(\footnote[40]{\textbf{NB} Tous les foncteurs marqués sont compatibles avec la $W$-équivalence grossière, et des dix homomorphismes fonctoriels sont des $W$-équivalences grossières.})

NB Les quatre derniers homomorphismes de foncteurs, relatifs au troisième diagramme (avec les indices 0) sont simplement induits par ceux relatifs au diagramme qui le précède. Donc les vérifications-clef sont relatives aux deux premiers diagrammes, et par symétrie, on n'a finalement qu'à regarder les \textit{deux} couples de foncteurs $(r,C)$, $(\lambda,\Phi)$, et pour chacun de ces couples, exhiber les deux homomorphismes fonctoriels, des deux composés avec les foncteurs identiques.

\printedpage{36}
% head: A. Grothendieck : Dérivateurs, Ch. VI, Édition des Universités de Montpellier II et Paris VII

\msp{page 58}

\dgrblock{dgrb222}{Naturality recovered after localization}
\begin{dcomment}{development}{p57}{Naturality recovered after localization}
Grothendieck wants the comparison to accommodate all functors over the base, since the categories Fib(S) and Liss(S) are full subcategories of Cat/S. The raw pullback construction does not have that naturality. The localization nevertheless supplies exactly the required repair. We spell out the following elementary elaboration of his argument; the “ours” label concerns this formulation and proof, not priority. We checked the free-fibration and strictification formulations in Emmenegger et al. (2024, Sections 2 and 5). They distinguish strict and pseudo preservation of cleavages but do not state the following repair for arbitrary functors over the base. We make no broader claim about its absence from the literature.

Localized naturality of pullback retractions. \ours Suppose W\ensuremath{\infty}\ensuremath{\subseteq}W. Choose normalized cleavages on fibrations X \ensuremath{\to}S. Then the functors \ensuremath{\gamma}X : \ensuremath{\Phi}X \ensuremath{\to}X determine a natural isomorphism between the functors induced by U\ensuremath{\Phi} and the identity on the right-fine localization of the full category Fib(S). The same conclusion holds for coarse localization.

Proof. First, any natural transformation over S between functors F, G : A \ensuremath{\to}B makes their images equal in (Cat/S)[(W d S)\ensuremath{-}1]. It is represented by a functor H : A \ensuremath{\times} [1] \ensuremath{\to}B over S. The projection r : A \ensuremath{\times} [1] \ensuremath{\to}A is right-fine: for every s its comma functor is the projection

(s\textbackslash{}A) \ensuremath{\times} [1] \ensuremath{-}\ensuremath{\to}s\textbackslash{}A,

a categorical homotopy equivalence. Its two sections become the same inverse of r, hence F = Hj0 and G = Hj1 become equal. The identical proof applies to coarse localization. If A is fibred, A \ensuremath{\times} [1] is fibred, so this argument also takes place within the indicated full subcategory.

Now let f : X \ensuremath{\to}Y be any functor over S, without a Cartesian hypothesis. For (s, x, u) define a vertical arrow

\ensuremath{\theta}f,(s,x,u) : f(u\ensuremath{*}

Xx) \ensuremath{-}\ensuremath{\to}u\ensuremath{*} Y f(x) by the equation

cY

u,f(x)\ensuremath{\theta}f,(s,x,u) = f(cX u,x). Existence and uniqueness follow from the Cartesian property of the lift in Y . This defines a natural transformation f\ensuremath{\gamma}X \ensuremath{\Rightarrow}\ensuremath{\gamma}Y \ensuremath{\Phi}(f) over S. To check naturality, compose the two candidate arrows with the chosen lift to the terminal object of a morphism (a, g) in \ensuremath{\Phi}X. Both composites are f(g)f(cX u,x), u,x) = f(gcX and both have the same base arrow; Cartesian uniqueness identifies them. The first paragraph of the proof therefore gives naturality after localization.

The comparison maps are right-fine equivalences. On a fibre, \ensuremath{\gamma}X,s is right adjoint to Xs ,\ensuremath{\to}s\textbackslash{}X: an arrow y \ensuremath{\to}x over u : s \ensuremath{\to}p(x) factors uniquely through u\ensuremath{*} Xx. Every adjunction gives a homotopy equivalence of nerves. Since source and target of \ensuremath{\gamma}X are fibrations, right-fine equivalence is equivalent to fibrewise equivalence. Thus the localized natural transformation is an isomorphism. Its coarse version follows either from the same argument or from \ensuremath{\gamma}X\ensuremath{\eta}X = id and coarse invertibility of \ensuremath{\eta}X. \ensuremath{\Box}

The defect before localization is real. Take S = [1], X = S, Y = S \ensuremath{\times} [1], and f(s) = (s, s). At the object (0, 1, 0 \ensuremath{\to}1) of \ensuremath{\Phi}X, the two functors in the comparison square have values (0, 0) and (0, 1). These objects are not isomorphic. The comparison \ensuremath{\theta} is the noninvertible arrow between them. Thus the source’s strict naturality assertion cannot simply be replaced by a claim of naturality up to isomorphism; the localization, or the restriction to splitting-preserving morphisms, performs an essential part of the repair.

\end{dcomment}


\cercle{2} Les localisés pour la \textit{$W$-équivalence fine droite.} Il n'y a plus lieu, on l'a dit, de regarder les catégories internes du diagramme p. 8, où ladite équivalence coïncide avec la $W$-équivalence grossière, pour laquelle on connait déjà les localisés. Il reste à regarder les catégories
\[\mathrm{Cat}/S,\ \ \ul{Liss}\,S,\ \ \ul{Fib}\,S,\ \ \ul{Fibsc}\,S.\]
Le lemme de la page 11 nous dit qu'il suffit de montrer que dans la situation
\[\ul{Fibsc}\,S \xrightarrow{\lambda} \mathrm{Cat}/S,\]
il y a un foncteur en sens inverse, avec les propriétés specifiées. Il en résultera que toute sous-catégorie pleine $\mathcal{F}$ de $\mathrm{Cat}/S$ qui contient $\ul{Fib}\,S$, quand on la munit du localiseur de la $W$-équivalence fine droite, donne un localisé équivalent, par le foncteur induit par l'inclusion, au localisé correspondant de $\mathrm{Cat}/S$, i.e. essentiellement à $\mathrm{HOT}_W(S)$. % [sic] specifiées -> spécifiées, missing accent

Comme foncteur en sens inverse, il s'impose de prendre encore $\Phi$. Il y a à vérifier trois choses :

a) $\Phi$ \texttt{[est]} compatible avec les localiseurs, i.e. si
\begin{tikzcd}[baseline=(\tikzcdmatrixname-2-2.base), row sep=small, column sep=small]
X \arrow[rr, "f"] \arrow[dr] & & Y \arrow[dl] \\
& S &
\end{tikzcd}
sur $S$ est une $W$-équivalence fine droite, i.e. les $s\backslash X \longrightarrow s\backslash Y$ sont des $W$-équivalences, alors $\Phi(f)$ aussi. Mais ladite équivalence coïncide
\msp{page 59}
dans $\ul{Fibsc}\,S$ (comme aussi dans $\ul{Fib}\,S$ et $\ul{Liss}\,S$) avec la $W$-équivalence par fibres, et d'autre part $\Phi(f)_s$ n'est autre que la flèche précédente $s\backslash X \longrightarrow s\backslash Y$. Donc il est évident que
\[f \in \underbrace{W_S^{\mathrm{d}}}_{W\text{-équivalence droite}} \iff \Phi(f) \in \underbrace{W_S^{\mathrm{f}}}_{W\text{-équivalence par fibres}} .\]

b) Que $X \xrightarrow{i_X} \lambda\Phi(X)$ est dans $W_S^{\mathrm{d}}$, pour tout $X$ dans $\mathrm{Cat}/S$. Cela signifie donc que $i_X$ induit
\[s\backslash X \longrightarrow s\backslash\Phi(X)\]
qui est une $W$-équivalence. Ce serait O.K. s'il était clair que $X \longmapsto \Phi(X)$ commute à tout changement de base propre sur $S$ (donc en particulier à $s\backslash S \longrightarrow S$), puisqu'on sait que $X \longrightarrow \Phi(X)$ est une $W$-équivalence. C'est sans doute trop demander, mais du moins que ça commute au changement de base qui est un co-isomorphisme local (telles les $s\backslash S \longrightarrow S$).

Pour un tel changement de base $S' \xrightarrow{p} S$, vérifier que la fibre de $\Phi(X')$ en $s'$ est isomorphe à celle de $\Phi(X)$ en $s = p(s')$, je m'en contenterai. Or cette fibre est $s'\backslash X'$, déduit de $X'$ sur $S'$ par le changement de base $s'\backslash S' \longrightarrow S'$, donc déduit de $X$ par $s'\backslash S' \longrightarrow S$. Or $s'\backslash S' \longrightarrow s\backslash S$ est un isomorphisme par hypothèse
\printedpage{37}
% head: A. Grothendieck : Dérivateurs, Ch. VI, Édition des Universités de Montpellier II et Paris VII
\msp{page 60}
sur $p$ (co-isomorphisme local), donc le changement de base $s'\backslash S' \longrightarrow S$ s'identifie à $s\backslash S \longrightarrow S$, et l'image inverse est bien isomorphe à $s\backslash X$, OK.

\dgrblock{dgrb223}{The comma resolution is right-fine}
\begin{dcomment}{exposition}{p59}{The comma resolution is right-fine}
Grothendieck has reduced the fine comparison to checking that the unit and counit are fine equivalences. This is the decisive point at which the comma categories preserve the relative diagram, rather than merely its total homotopy type. A direct contraction proves the assertion for every X \ensuremath{\to}S and makes the proposed base-change verification unnecessary.

Fix s \ensuremath{\in}S. An object of s\textbackslash{}\ensuremath{\Phi}X consists of

\ensuremath{-}\ensuremath{\to}vt

\ensuremath{-}\ensuremath{\to}up(x). s Define R : s\textbackslash{}\ensuremath{\Phi}X \ensuremath{\to}s\textbackslash{}X by composition, R(v, u, x) = (x, uv). A morphism represented by (a, f) satisfies av = v\ensuremath{\prime} and p(f)u = u\ensuremath{\prime}a; it is sent to f, which satisfies p(f)uv = u\ensuremath{\prime}v\ensuremath{\prime}. Thus R is a functor. The functor J = s\textbackslash{}\ensuremath{\eta}X sends (x, w) to the factorization

\ensuremath{-}\ensuremath{\to}p(x)

\ensuremath{-}\ensuremath{\to}p(x). s We have RJ = id. The arrow (u, idx) gives a natural transformation

ids\textbackslash{}\ensuremath{\Phi}X \ensuremath{-}\ensuremath{\to}JR.

It is over s, since uv is the composite arrow of the target factorization. These identities prove that J and R are inverse up to categorical homotopy. Hence \ensuremath{\eta}X belongs to W d S whenever W contains W\ensuremath{\infty}.

For a split fibration Y , the fibre of the counit is

(y, u) 7\ensuremath{-}\ensuremath{\to}u\ensuremath{*}y.

\ensuremath{\gamma}Y,s : s\textbackslash{}Y \ensuremath{-}\ensuremath{\to}Ys, If K : Ys ,\ensuremath{\to}s\textbackslash{}Y is the inclusion, Cartesian factorization gives

Homs\textbackslash{}Y (Kz, (y, u)) \ensuremath{\sim}= HomYs(z, u\ensuremath{*}y).

Thus K \ensuremath{\dashv}\ensuremath{\gamma}Y,s, with identity unit for a normalized splitting. In particular both functors are W\ensuremath{\infty}equivalences. The counit is therefore fibrewise, and hence right-fine, an equivalence. Finally, \ensuremath{\Phi}(f)s is literally s\textbackslash{}f; this verifies compatibility of \ensuremath{\Phi} with the two localizers without an additional argument. All parts of the comparison now fit together. The unit and counit induce inverse natural isomorphisms between the localized composites. We obtain

S)\ensuremath{-}1] \ensuremath{\simeq}Fibsc(S)[W \ensuremath{-}1

(Cat/S)[(W d obj]. If F is any full subcategory of Cat/S containing the fibrations, the resolution U\ensuremath{\Phi} lands in F and its unit is a morphism of F on restriction. The same argument proves the asserted equivalence with F[(W d S)\ensuremath{-}1]. Reversing the comma directions gives the corresponding statement for opfibrations and W g S.

\end{dcomment}


c) Que de même
\[\Phi\lambda \longrightarrow \mathrm{id}\]
est non seulement dans $W_S$ (grossier), mais que c'est une $W$-équivalence par fibres. Donc on part de $X$ fibrée sur $S$. On regarde
\[\Phi(X) \longrightarrow X,\]
pour voir que c'est une $W$-équivalence par fibres, il suffit de voir qu'il en est ainsi de $X \longrightarrow \Phi(X)$,
\[X_s \longrightarrow s\backslash X\]
est une $W$-équivalence pour tout $s \in S$. Mais ceci est déjà vrai dès que $X$ est lisse sur $S$.

Comme le cas de la $W$-équivalence fine gauche se ramène au cas actuel par $S \longmapsto S^{o}$, on a donc prouvé tout ce qu'était attendu, à la seule exception de la pleine fidèlité des foncteurs $i$, $j$ de la page 51. % [sic] ce qu'était -> ce qui était, grammar; fidèlité -> fidélité, accent

Mais dans le résumé prospectif de la situation pour les localiseurs, p. 49-51, j'ai oublié d'autres foncteurs que $i$, $j$, de nature si possible encore plus triviale. Ce sont
\msp{page 61}
ceux qu' on peut décrire à partir du foncteur identique, sur une des catégories du diagramme où aient cours deux, voire trois, notions de $W$-équivalence différentes. Ce sont, je le rappelle

\dgrblock{dgrb224}{The hypotheses of the fine comparison}
\begin{dcomment}{exposition}{p60}{The hypotheses of the fine comparison}
The comparison just established has survived the earlier changes of axioms. Its proof uses the inclusion W\ensuremath{\infty}\ensuremath{\subseteq}W, the elementary saturation properties, and the smooth-fibre criterion already checked on p. 42. In particular it does not use the new L5. On split fibrations, right-fine equivalences are exactly objectwise W-equivalences of the associated strict diagrams; the analogous statement holds on split opfibrations. This is the part of the construction which remains useful for homological and truncated equivalences.

The pure coarse comparison with the internal categories also needs no new L5. The path replacement C produces homotopy fibrations, its unit X \ensuremath{\to}C(X) is a total homotopy equivalence, and naturality together with two-out-of-three makes C preserve coarse equivalences. L5 enters in the additional identification of coarse equivalences between these replacements with fibrewise, hence fine, equivalences. This is what makes C induce the functor back into the fine category used for the reflection. The later annotation correctly says that this requirement is highly restrictive. The commentary on pp. 43--45 explains why it must not be silently transferred to arbitrary homological localizers.

For W = W\ensuremath{\infty}, the modern literature gives a precise model for the coarse part: small categories, nerve weak equivalences, and smooth proper functors form a category of fibrant objects; perfect objects over a base suffice to compute the corresponding localization. See Cisinski (2009, Proposition 4.4 and Corollary 4.5). This places Grothendieck’s path replacement inside a developed homotopical framework. The explicit comma proof above is still useful, since it shows exactly which comparison remains valid when the coefficient equivalences are enlarged.

\end{dcomment}


\begin{enumerate}[label={}]
\item[$1^{o}$)] Les quatre catégories
\[(*) \qquad\qquad \ul{Fibsc}\,S \longrightarrow \ul{Fib}\,S \longrightarrow \ul{Liss}\,S \longrightarrow \mathrm{Cat}/S, \qquad\qquad W_S^{\mathrm{d}} \subseteq W\]
avec la $W$-équivalence fine droite d'une part, grossière de l'autre, et dualement
\[(*') \qquad\qquad \ul{Cofibsc}\,S \longrightarrow \ul{Cofib}\,S \longrightarrow \ul{Prop}\,S \longrightarrow \mathrm{Cat}/S, \qquad\qquad W_S^{\mathrm{g}} \subseteq W\]
pour la $W$-équivalence fine gauche d'une part, la grossière de l'autre, et de plus
\item[$2^{o}$)] $\mathrm{Cat}/S$, qui était inclue dans les deux suites précédentes, et qui est la seule catégorie du diagramme donnant lieu à trois notions de $W$-équivalence différentes % [sic] inclue -> incluse, past participle of inclure
\end{enumerate}

\[
\begin{tikzcd}[row sep=small, column sep=large]
W_S^{\mathrm{d}} \arrow[dr, hook] & \\
& W_S \\
W_S^{\mathrm{g}} \arrow[ur, hook] &
\end{tikzcd}
\qquad\qquad .
\]

\printedpage{38}
% head: A. Grothendieck : Dérivateurs, Ch. VI, Édition des Universités de Montpellier II et Paris VII

Dans chacun des quatre cas ($C$, munie de deux localiseurs $\Sigma \subseteq \Sigma'$) de la suite $(*)$, le foncteur
\[C\Sigma^{-1} \longrightarrow C\Sigma'^{-1}\]
est essentiellement ‘le même’, à cause des théorèmes de comparaison. Il suffit de regarder ce foncteur dans un seul das, p. ex. un des deux extrêmes $\ul{Fibsc}\,S$ ou $\mathrm{Cat}/S$. Même remarque pour la suite $(*')$. Donc finalement il reste deux foncteurs remarquables % [sic] das -> cas, misprint
\msp{page 62}
à examiner, qui n'ont pas l'air du tout de provenir des foncteurs du diagramme p. 8 ou de quelque composé (ça ne donne que des équivalences canoniques ou bien les foncteurs $i$, $j$ de page 51, quand on va des catégories internes au diagramme vers les catégories bordantes). Le plus joli est sans doute de les regarder à partir de $\mathrm{Cat}/S$, puisque c'est là seulement qu'on les a tous les deux à la fois
\[
\begin{tikzcd}[row sep=small, column sep=tiny]
(\mathrm{Cat}/S)(W_S^{\mathrm{d}})^{-1} \approx \mathrm{HOT}_W(S) \arrow[dr, "\rho"] & \\
& (\mathrm{Cat}/S)W_S^{-1} \approx \mathrm{HOT}_W^{\mathrm{lc}}(S). \\
(\mathrm{Cat}/S)(W_S^{\mathrm{g}})^{-1} \approx \mathrm{HOT}_W^{o}(S) \arrow[ur, "\sigma"'] &
\end{tikzcd}
\]
Par contre, les foncteurs (qu'on a envie d'appeler ‘foncteurs d'inclusion’) $i$, $j$ de la page 51, en sens inverse de $\rho$, $\sigma$, sont ceux qui proviennent, par localisation $W_S^{\mathrm{d}}$ ou $W_S^{\mathrm{g}}$ selon le cas, d'un des foncteurs du diagramme page 8 (ou d'un composé de tels foncteurs), allant d'une des six catégories internes (où toutes les $W$-équivalences envisagées coïncident) vers l'une des sept catégories bordantes, sauf quand cette dernière est $\mathrm{Cat}/S$, est déterminé déjà avec quelle $W$-équivalence, $W_S^{\mathrm{d}}$ ou $W_S^{\mathrm{g}}$, on va travailler \texttt{[structure de cette phrase?]}. La détermi-\msp{page 63}% hyphen kept as printed: the word is divided at the manuscript page marker
nation peut-être la plus jolie est par
\[(*) \qquad\qquad \ul{Parf}\,S \longrightarrow \mathrm{Cat}/S\]
qui donne le foncteur $i$ quand on travaille avec $W_S^{\mathrm{d}}$, le foncteur $j$ quand on travaille avec $W_S^{\mathrm{g}}$. On voit alors tout de suite que si on compose avec le foncteur pertinent, soit $\rho$ soit $\sigma$, défini par l'inclusion $W_S^{\mathrm{d}} \subseteq W_S$ resp. $W_S^{\mathrm{g}} \subseteq W_S$ (comme parties de $\mathrm{Fl}(\mathrm{Cat}/S)$), on trouve le foncteur
\[(**) \qquad\qquad (\ul{Parf}\,S)W_S^{-1} \xrightarrow{\approx} (\mathrm{Cat}/S)W_S^{-1}\]
qui est l'\textit{équivalence canonique} entre ces deux catégories. Ainsi le foncteur correspondant à $i$, $j$ page 51, exprimé en termes de $\mathrm{Cat}/S$ comme foncteur en sens inverse de $\rho$, $\sigma$ (p. 62), est-il le composé de l'équivalence quasi-inverse de $(**)$, laquelle s'explicite par le foncteur $C$, et du foncteur sur le localisé $W_S^{\mathrm{d}}$ ou $W_S^{\mathrm{g}}$ déduit de $(*)$, i.e.

% ---- from Der-ChVI-345_p041-045.tex
\printedpage{39}
% head: A. Grothendieck : Dérivateurs, Ch. VI, Édition des Universités de Montpellier II et Paris VII

\[
\begin{tikzcd}[column sep=large, row sep=large]
 & \HOT_W(S)\approx(\Cat/S)(W_S^{\mathrm{d}})^{-1} \\
\HOT^{\mathrm{lc}}_W(S)\approx(\Cat/S)W_S^{-1} \arrow[ur, "i"] \arrow[r, "C"] \arrow[dr, "j"'] & (\ul{Parf}\,S)W_S^{-1} \arrow[u, "\text{\footnotesize car $W_S=W_S^{\mathrm{d}}$ sur $\ul{Parf}$}"'] \arrow[d, "\text{\footnotesize car $W_S=W_S^{\mathrm{g}}$ sur $\ul{Parf}$}"] \\
 & \HOT^{o}_W(S)\approx(\Cat/S)(W_S^{\mathrm{g}})^{-1}.
\end{tikzcd}
\]

\msp{page 64}
On voit par là qu'on a des isomorphismes canoniques
\[
\rho i\simeq\id,\qquad\sigma j\simeq\id
\]
dans la catégorie des endofoncteurs de $(\Cat/S)W_S^{-1}$ ($\approx\HOT^{\mathrm{lc}}_W(S)$), isomorphisme qui s'explicite ici par le fait que $\rho i$, $\sigma j$ sont déduits du même foncteur $C$ par passage aux localisés $W_S^{-1}$ (donc en fait $\rho i=\sigma j$ !), et qu'on a une flèche fonctorielle en $X$
\[
X\xrightarrow{i_X}C(X),
\]
qui est dans $W$ pour tout $X$. Ainsi, $\rho$ et $\sigma$ sont à isomorphisme près des \textit{rétractions} des catégories $\HOT_W(S)$, $\HOT^{o}_W(S)$ sur la catégorie $\HOT^{\mathrm{lc}}_W(S)$. De plus, ces rétractions sont ce qu'on appelle des foncteurs de localisation (comme on voit sur le diagramme de la page 62). (On dit que $C\xrightarrow{F}C'$ est un \textit{foncteur de localisation} si, désignant par $\Sigma\subseteq\Fl\,C$ l'ensemble des flèches de $C$ transformées en des isomorphismes par $F$, le foncteur
\[
C\Sigma^{-1}\xrightarrow{\bar{F}}C'
\]
qui factorise $F$ est une équivalence de catégories.)

\msp{page 65}
On sait d'ailleurs que si un foncteur $f$ a un adjoint $g$ (à gauche ou à droite), alors $f$ est un foncteur \texttt{[de]} localisation si et seulement si l'adjoint $g$ est pleinement fidèle. Pour prouver que $i$, $j$ sont pleinement fidèles, il serait donc tentant de prouver qu'ils sont adjoints des foncteurs $\rho$, $\sigma$.

\dgrblock{dgrb225}{What the section of localization still has to prove}
\begin{dcomment}{exposition}{p64}{What the section of localization still has to prove}
Grothendieck has now produced the same locally constant category from both comma directions. The identity of Cat/S induces

S)\ensuremath{-}1] \ensuremath{-}\ensuremath{\to}(Cat/S)[W \ensuremath{-}1

\ensuremath{\rho} : (Cat/S)[(W d S ], and path replacement induces a functor in the reverse direction. The equation \ensuremath{\rho}i \ensuremath{\simeq}id says that the coarse category is a retract. It does not yet say that i is fully faithful; this is the distinction already encountered at manuscript p. 38.

The required additional datum is now visible. There is a natural arrow X \ensuremath{\to}C(X) in the fine category, and it becomes invertible in the coarse category. Its direction suggests a unit for \ensuremath{\rho} \ensuremath{\dashv}i, not for the opposite adjunction. The inverse coarse arrow C(X) \ensuremath{\to}X then supplies a counit. Grothendieck’s change of orientation on p. 67 is consequently forced by where the unit is invertible.

A localization with a fully faithful right adjoint is a reflection onto a full subcategory. In this case its image consists of the relative homotopy types whose transport along every base arrow is invertible. The adjective “locally constant” therefore describes an intrinsic condition on a diagram, not constancy of a chosen presentation of its total category. The path construction supplies a representative satisfying that condition and a universal map to it. The argument on the next pages establishes that universality through the two triangle identities, so that no fullness conclusion is extracted merely from a retraction.

\end{dcomment}


Pour nous y reconnaître, je vais poser
\[
\left\{
\begin{aligned}
&\cA=\Cat/S,\ \ \Sigma_0=W_S^{\mathrm{d}},\ \ \Sigma=W_S\qquad\quad(\subseteq\Fl(\cA))\\
&C:\cA\longrightarrow\cA\ \text{ le foncteur familier, qui a la vertu}\\
&\qquad\qquad\qquad\quad C(\Sigma)\subseteq\Sigma_0\quad.
\end{aligned}
\right.
\]

On a donc la situation
\printedpage{40}
% head: A. Grothendieck : Dérivateurs, Ch. VI, Édition des Universités de Montpellier II et Paris VII
\[
(*)\qquad\qquad(\cA,\Sigma_0)\ \underset{C}{\overset{\rho_0\,=\,\id_{\cA}}{\rightleftarrows}}\ (\cA,\Sigma),
\]
induisant par passage aux localisés
\[
\cA\Sigma_0^{-1}\quad\underset{\bar{C}}{\overset{\rho\,=\,\bar{\rho}_0}{\rightleftarrows}}\ \cA\Sigma^{-1},
\]
on veut vérifier que $\rho$ et $C$ \texttt{[plutôt $\bar{C}$]} sont adjoints l'un de l'autre, $C$ \texttt{[plutôt $\bar{C}$]} \textit{sûrement à droite,} ce qui signifie que l'on doit avoir des homomorphismes fonctoriels
\[
\left\{
\begin{aligned}
&\id_{\cA\Sigma^{-1}}&&\rightleftarrows&&\rho\,\bar{C}\quad(=\overline{\rho_0C})\\
&\id_{\cA\Sigma_0^{-1}}&&\rightleftarrows&&\bar{C}\rho\quad(=\overline{C\rho_0})
\end{aligned}
\right.
\]
(\footnote[41]{mais si $C$ est à \textit{droite,} les flèches ici doivent aller en sens inverse et c'est bien à quoi je dois me résoudre p.~67.}) \texttt{[les flèches en bas vont en direction corrigée, mais les flèches en haut ne sont pas effacées]} satisfaisant deux relations de compatibilité sur lesquelles on reviendra. On peut espérer trouver ces homomorphismes fonctoriels à partir de
\[
\left\{
\begin{aligned}
&\id_{\cA}&&\xrightarrow{\ i\ }&&\rho_0C&&=&&C\\
&\id_{\cA}&&\xleftarrow{\ p\ }&&C\rho_0&&=&&C
\end{aligned}
\right.
\]

\msp{page 66}
et dans notre cas on dispose bel et bien de
\[
i:\id_{\cA}\longrightarrow C\qquad\text{i.e. de }\ i_X:X\longrightarrow C(X)
\]
fonctoriel en $X\in\Ob\,\cA$ (et de plus $i_X$ est dans $\Sigma$). Il n'y a pas cependant de flèche naturelle en sens inverse $p_X$
\[
C(X)\xrightarrow{\ ?\ }X.
\]
J'avais songé à la projection $\pi_X:C(X)\longrightarrow X$, déduite de la construction $C(X)\simeq X\times_S(\ul{Ch}_\infty(S),\sigma)$, mais ce n'est malheureusement pas un $S$-morphisme. Comme on cherche une flèche dans $\cA\Sigma_0^{-1}$, on pourrait se contenter d'avoir une flèche fonctorielle en sens inverse, du moment qu'elle soit dans $\Sigma_0$, donc puisse s'inverser quand on travaille dans $\cA\Sigma_0^{-1}$. Le candidat evident est $i_X:X\longrightarrow C(X)$ dont on vient de parler. L'ennui, c'est que $i_X$ est seulement dans $\Sigma$, pas dans $\Sigma_0$ (\footnote[42]{S'il n'en était pas ainsi, le type d'homotopie relatif sur $S$ défini par $X/S$ serait localement constant (puisqu'il en est ainsi pour celui défini par $C(X)$), ce qui n'est pas toujours le cas.}). Pourtant, quand $X$ est lui-même de la forme $C(Y)$ (plus généralement, s'il est un Hot-fibré sur $S$), alors le fait que $X\longrightarrow C(X)$ soit dans $\Sigma$ implique déjà qu'il soit dans $\Sigma_0$. Notons donc ceci % [sic] evident -> évident, accent missing
\[
\left\{
\begin{array}{lrl}
i:\id_{\cA}\longrightarrow C\ \text{tel que} & 1^o) & i_X:X\longrightarrow C(X)\ \textit{dans}\ \Sigma\\[1ex]
 & & \text{pour tout }X\in\Ob\,\cA\\[1ex]
 & 2^o) & i_X\ \text{dans }\Sigma_0\text{ si }X\text{ de la}\\[1ex]
 & & \text{forme }C(Y)
\end{array}
\right.
\]

\printedpage{41}
% head: A. Grothendieck : Dérivateurs, Ch. VI, Édition des Universités de Montpellier II et Paris VII

\msp{page 67}
L'idée me vient, après des essais merdouillants \texttt{[sic]}, que $\bar{C}$ pourrait être l'adjoint à droite, et non à gauche \texttt{[c'est la version corrigée (droite et gauche échangées) cf.~p.~65]}, i.e. que les flèches cherchées pourraient aller en sens inverse
\[
\left\{
\begin{aligned}
&\id_{\cA\Sigma^{-1}}&&\xleftarrow{\ p\ }&&\rho\,\bar{C}\quad(=\overline{\rho_0C})\\
&\id_{\cA\Sigma_0^{-1}}&&\xrightarrow{\ i\ }&&\bar{C}\rho\quad(=\overline{C\rho_0})
\end{aligned}
\right.
\]
(j'inverse du coup les notations pour $p$ et $i$), donc on est amené si possible à chercher des flèches
\[
\begin{aligned}
&\id_{\cA}&&\xleftarrow{\ p\ }&&\rho_0C&&=&&C\\
&\id_{\cA}&&\xrightarrow{\ i\ }&&C\rho_0&&=&&C,
\end{aligned}
\]
donc des flèches fonctorielles en $X$
\[
X\xrightarrow{i_X}C(X)\xrightarrow{p_X}X.
\]
On tient bel et bien une $i_X$ fonctorielle - toujours la même! Pour $p_X$, on n'en a pas plus qu'avant, mais se rappelant qu'il suffit de le trouver dans $\cA\Sigma^{-1}$ (et non plus $\cA\Sigma_0^{-1}$!), il suffit d'inverser la flèche $i_X$, qui est en effet un \textit{isomorphisme} dans $\cA\Sigma^{-1}$ (mais non dans $\cA\Sigma_0^{-1}$). On posera donc
\[
p_X=i_X^{-1}:C(X)\longrightarrow X\qquad\quad\text{dans }\cA\Sigma^{-1}.
\]
Il faut prouver deux compatibilités seulement, pour les foncteurs $\bar{C}$ et $\rho$ entre $\cA\Sigma_0^{-1}$ et $\cA\Sigma^{-1}$ (et j'écrirai $C$ au lieu de $\bar{C}$) : les deux composés
\[
\left\{
\begin{array}{l}
\begin{tikzcd}[ampersand replacement=\&]
\rho \arrow[r, "\rho*i"] \arrow[rr, bend left=20, dotted, "\id"] \& \rho\,C\rho \arrow[r, "p*\rho"] \& \rho
\end{tikzcd}
\qquad\quad\text{\footnotesize (relation pour foncteurs à valeurs dans $\cA\Sigma^{-1}$)}\\[2ex]
\begin{tikzcd}[ampersand replacement=\&]
C \arrow[r, "i*C"] \arrow[rr, bend right=20, dotted, "\id"'] \& C\rho\,C \arrow[r, "C*p"] \& C
\end{tikzcd}
\qquad\quad\text{\footnotesize (relation pour foncteurs à valeurs dans $\cA\Sigma_0^{-1}$)}
\end{array}
\right.
\]
doivent être l'identité $\id_\rho$ resp. $\id_C$.

\msp{page 68}
La première relation se vérifie tautologiquement, elle revient à dire que dans $\cA\Sigma^{-1}$, on a un composé identité des
\[
X\xrightarrow{i_X}C(X)\xrightarrow{i_X^{-1}}X\qquad\quad!
\]
La seconde signifie que pour tout $X$, on a un composé identité dans $\cA\Sigma_0^{-1}$
\[
C(X)\xrightarrow{i_{C(X)}}C(C(X))\xrightarrow{C(p_X)}C(X).
\]
Or $p_X:C(X)\longrightarrow X$ avait été défini dans $A\Sigma^{-1}$ comme inverse $i_X^{-1}$ de $i_X$, d'autre part le foncteur $C$ sur $\cA$ définit par hypothèse $\cA\Sigma^{-1}\longrightarrow\cA\Sigma_0^{-1}$, plus précisément $C(\Sigma)\subseteq\Sigma_0$ % [sic] $A\Sigma^{-1}$ (italic A) -> $\mathcal{A}\Sigma^{-1}$, elsewhere script A
\printedpage{42}
% head: A. Grothendieck : Dérivateurs, Ch. VI, Édition des Universités de Montpellier II et Paris VII
ainsi $C(i_X)$ est inversible dans $\cA\Sigma_0^{-1}$, et $C(p_X)=C(i_X)^{-1}$. Donc la compatibilité revient à la relation
\[
C(i_X)^{-1}i_{C(X)}=\id\qquad\quad\textit{dans}\ \ \cA\Sigma_0^{-1},
\]
i.e.
\[
\boxed{\,i_{C(X)}=C(i_X)\,}:C(X)\longrightarrow C^2(X).
\]
En d'autre termes, il faut prouver le : % [sic] d'autre termes -> d'autres termes, plural

\dgrblock{dgrb226}{A precise reflection lemma}
\begin{dcomment}{exposition}{p67}{A precise reflection lemma}
The remaining argument is formal once the two classes of arrows and the two uses of the unit are distinguished. We state the lemma in the form needed here; it also shows that the iterated-unit equality can be established without further geometry. This is the ordinary categorical form of the recognition principle for reflective localizations; compare Lurie (2009, Proposition 5.2.7.4) for its higher categorical formulation.

Lemma. Let \ensuremath{\Sigma}0 \ensuremath{\subseteq}\ensuremath{\Sigma} be classes of arrows of a category A, and suppose the localizations below exist. Let C : A \ensuremath{\to}A and \ensuremath{\eta} : id \ensuremath{\to}C satisfy

C(\ensuremath{\Sigma}) \ensuremath{\subseteq}\ensuremath{\Sigma}0,

\ensuremath{\eta}X \ensuremath{\in}\ensuremath{\Sigma}, \ensuremath{\eta}CX \ensuremath{\in}\ensuremath{\Sigma}0 (X \ensuremath{\in}A). Then the identity and C induce functors

P : A0 = A[\ensuremath{\Sigma}\ensuremath{-}1

0 ] \ensuremath{-}\ensuremath{\to}A1 = A[\ensuremath{\Sigma}\ensuremath{-}1], R : A1 \ensuremath{-}\ensuremath{\to}A0, with P \ensuremath{\dashv}R and R fully faithful. The unit is induced by \ensuremath{\eta}, and the counit is the inverse of its image in A1.

Since C(\ensuremath{\Sigma}0) \ensuremath{\subseteq}\ensuremath{\Sigma}0, the endofunctor C and its natural transformation descend to an Proof. endofunctor T = RP and a transformation e : idA0 \ensuremath{\to}T. Naturality extends across the formal inverses because both functors invert the required arrows. At every object x, both

a = eT x : Tx \ensuremath{\to}T 2x,

b = T(ex) : Tx \ensuremath{\to}T 2x are isomorphisms: the first by the third hypothesis, the second by the first two. We claim that a = b.

Naturality at ex gives bex = aex. Apply T and cancel the isomorphism b on the right to obtain T(b) = T(a). Naturality at a gives

T(a)a = eT 2xa,

so T(a) = eT 2x. Naturality at b gives

T(b)a = eT 2xb.

Substitute the preceding equalities and cancel the isomorphism T(a). This proves a = b. Notice that we have never cancelled ex, which need not be invertible in A0.

Let \ensuremath{\epsilon} : PR \ensuremath{\to}idA1 be the inverse of the transformation induced by \ensuremath{\eta} in A1. Its first triangle identity is \ensuremath{\epsilon}P xP(ex) = idP x, true by construction. For the second, every object of A1 is represented by an object X of A. On its image under R, the composite R(\ensuremath{\epsilon}X)eRX is b\ensuremath{-}1a, which is the identity by the calculation above. Hence P \ensuremath{\dashv}R. The counit is an isomorphism, so R is fully faithful. \ensuremath{\Box}

The full image is exactly the class of x for which ex is invertible. Objects in the image satisfy this by the third hypothesis; conversely, invertibility identifies x with RPx. Applied to \ensuremath{\Sigma}0 = W d S and \ensuremath{\Sigma} = WS, this proves the desired full faithfulness of i. The left-comma version proves the corresponding assertion for j, with the covariant convention explained at p. 50. Grothendieck’s concatenation argument is a particularly transparent geometric verification of the same identity, and remains useful for understanding the replacement itself.

\end{dcomment}


\textbf{Lemme.} \textit{Soit $X$ dans $\Cat/S$, considérons les deux flèches}
\[
i_{C(X)},\ C(i_X):C(X)\rightrightarrows C(C(X)).
\]
\textit{Ces flèches sont égales, sinon au sens absolu, du moins dans $(\Cat/S)(W_S^{\mathrm{d}})^{-1}$.}

Ici $f_X:X\longrightarrow S$ est la projection standard et
\[
\begin{aligned}
C(X)\ &=\ \quad\{(x,c)\ \mid\ x\in\Ob\,X,\ c:f_X(x)\ \underbrace{\longrightarrow}_{\text{chemin dans }S}\ t\}\\
C^2(X)\ &=\ \{(x,c,c')\ \mid\ x\in\Ob\,X,\ \underbrace{f_X(x)\xrightarrow{c}t\xrightarrow{c'}t'}_{c,\ c'\text{ chemins dans }S}\},
\end{aligned}
\]
les projections dans $S$ sont données par
\[
\begin{aligned}
(x,c)&&\longmapsto&&t&&=&&\text{but }c\\
(x,c,c')&&\longmapsto&&t'&&=&&\text{but }c'
\end{aligned}
\]

\msp{page 69}
On trouve
\[
i_{C(X)}(x,c)=(x,c,1_t):
\begin{tikzcd}
x \arrow[d] & & \\
f_X(x) \arrow[r, "c"] & t \arrow[r, "1_t", "\text{chemin vide}"'] & t
\end{tikzcd}
\]
\[
C(i_X)(x,c)=(x,1_{f_X(x)},c):
\begin{tikzcd}
x \arrow[d] & & \\
f_X(x) \arrow[r, "1_{f_X(x)}", "\text{chemin vide}"'] & f_X(x) \arrow[r, "c"] & t
\end{tikzcd}
\]
Mais on a le choix entre deux théories de chemins, pour construire $C(X)$ : soit avec les $\ul{Ch}_\infty(S)$, soit avec les $\ul{Ch}(S)$. ‘Ça donne les mêmes résultats’, mais suivant les situations l'un est plus commode que l'autre. Je vais travailler avec les $\ul{Ch}(S)$, auquel cas les $1_S$ sont des unités à gauche et à droite, et on a une loi de composition pour les chemins. Cette loi nous permet de définir
\[
\begin{aligned}
C^2(X)&&\xrightarrow{\ \gamma\ }&&C(X)\\
(x,c,c')&&\longmapsto&&(x,c'\circ c),
\end{aligned}
\]
et les composés de $i_{C(X)}$ et de $C_{i(X)}$ avec cette flèche sont $\id_{C(X)}$ dans les deux cas % [sic] $C_{i(X)}$ -> $C(i_X)$, the arrow considered is C(i_X)
\printedpage{43}
% head: A. Grothendieck : Dérivateurs, Ch. VI, Édition des Universités de Montpellier II et Paris VII
\[
\begin{tikzcd}
C(X) \arrow[r, shift left] \arrow[r, shift right] \arrow[rr, bend right=15, "\id"'] & C^2(X) \arrow[r, "\gamma"] & C(X).
\end{tikzcd}
\]

Donc pour montrer que $i_{C(X)}=C(i_X)$ dans $(\Cat/S)(W_S^{\mathrm{d}})^{-1}$, il suffit de prouver que $\gamma$ est dans $W_S^{\mathrm{d}}$. Mais on sait que le foncteur C transforme $W_S$ en $W_S^{\mathrm{d}}$, donc $C(i_X)$ est dans $W_S^{\mathrm{d}}$, et comme $\gamma\circ C(i_X)=\id$, $\gamma$ est dans $W_S^{\mathrm{d}}$, OK! % [sic] foncteur C (upright roman C) -> foncteur $C$, math italic elsewhere

\msp{page 70}
J'ai envie de grouper ce qui a été prouvé dans un énoncé récapitulatif.

\dgrblock{dgrb227}{Concatenation and the two iterated units}
\begin{dcomment}{exposition}{p69}{Concatenation and the two iterated units}
Grothendieck’s path model makes the reflection tangible. An object of C(X) is an object x together with a path from p(x) to its new base point. Iterating C adds a second path. The two units place the empty path in different positions, so there is no reason for the resulting functors to be literally equal before localization.

Using the finite-path version with strict empty-path units, concatenation gives a functor over S

(x, c, c\ensuremath{\prime}) 7\ensuremath{-}\ensuremath{\to}(x, c\ensuremath{\prime}c).

\ensuremath{\mu}X : C2(X) \ensuremath{-}\ensuremath{\to}C(X), On morphisms it concatenates the two compatible path diagrams; compatibility of their endpoints makes this well-defined, and composition is respected because it is respected in each path segment. The two empty-path laws give the strict equations

\ensuremath{\mu}X\ensuremath{\eta}C(X) = idC(X),

\ensuremath{\mu}XC(\ensuremath{\eta}X) = idC(X). The relevant homotopical input is C(WS) \ensuremath{\subseteq}W d S. Since \ensuremath{\eta}X \ensuremath{\in}WS, we have C(\ensuremath{\eta}X) \ensuremath{\in}W d S. Two-out-ofthree applied to the second equation yields \ensuremath{\mu}X \ensuremath{\in}W d S. After fine localization both \ensuremath{\eta}C(X) and C(\ensuremath{\eta}X) are therefore the inverse of the same arrow \ensuremath{\mu}X, which proves their equality.

This proof does not cancel the unit X \ensuremath{\to}C(X) in the fine category. That would erase precisely the distinction the chapter has worked to establish. Instead it cancels a map between already completed objects. Nor does the argument require the associativity of concatenation: the two unit identities and two-out-of-three suffice for the triangle equation. Associativity equips C with the familiar monadic structure, but the adjunction needs only the part displayed here.

The formal lemma in the preceding block shows that the equality also follows from the stated hypotheses on pp. 65--66. Thus there is no outstanding coherence gap in this part of the proof. The choice of finite paths supplies an explicit common inverse and exhibits why adjoining paths becomes idempotent after the appropriate localization.

\end{dcomment}


\dgrblock{dgrb228}{Path completion and the modern category of local systems}
\begin{dcomment}{later}{p69}{Path completion and the modern category of local systems}
The theorem reaches its most familiar modern form when W = W\ensuremath{\infty}. Let S denote the infinitycategory of spaces, and let

\ensuremath{\ell}: N(S)op \ensuremath{-}\ensuremath{\to}\ensuremath{\Pi}\ensuremath{\infty}N(S) be groupoid completion; we identify the completion of N(S)op with that of N(S) by reversing paths. Then the fine theory models homotopy coherent contravariant space-valued diagrams, and the locally constant theory models diagrams which carry every base arrow to an equivalence:

HOTW\ensuremath{\infty}(S) \ensuremath{\simeq}Ho Fun(N(S)op, S),

HOTlc

W\ensuremath{\infty}(S) \ensuremath{\simeq}Ho Fun(\ensuremath{\Pi}\ensuremath{\infty}N(S), S). The first comparison can be justified through an actual model equivalence. Thomason’s adjunction

c Sd2 : sSet \ensuremath{\rightleftarrows}Cat : Ex2 N

is a Quillen equivalence for nerve weak equivalences (Thomason, 1980, Theorem 4.9 and the introduction). Apply both functors objectwise to Sop-diagrams with their projective model structures. The right adjoint preserves fibrations and trivial fibrations objectwise. If F is projectively cofibrant and G projectively fibrant, the criterion for an adjoint pair of maps to be weak equivalences is checked at each object using the original Quillen equivalence; projective cofibrancy implies objectwise cofibrancy. Thus the diagram adjunction is again a Quillen equivalence. The simplicial-set diagram model then presents the coherent functor category by Lurie (2009, Proposition 4.2.4.4). This explains why strict category-valued diagrams, after objectwise localization, retain coherent space-valued diagrams.

Here Ho is taken after forming the category of coherent diagrams. It is not the category of ordinary functors into Ho(S). The higher transport information discussed earlier in the chapter is retained by this order of operations.

Restriction \ensuremath{\ell}\ensuremath{*}is fully faithful by the universal property of groupoid completion. Its left adjoint is left Kan extension \ensuremath{\ell}!, and the reflector on diagrams is \ensuremath{\ell}\ensuremath{*}\ensuremath{\ell}!. This identifies the universal property of Grothendieck’s C, rather than asserting that the particular path-category construction is literally a chosen infinity-categorical Kan extension. The standard straightening comparison and the adjunction for Kan extension are given in Lurie (2009, Theorem 2.2.1.2 and Proposition 4.3.3.7); the reflection for a left proper combinatorial model category of coefficients is Cisinski (2009, Corollary 2.5).

A local system of spaces on the infinity-groupoid \ensuremath{\Pi}\ensuremath{\infty}N(S) can equivalently be viewed as a space over |N(S)|, by taking its total space and homotopy fibres. Consequently the coarse theory has the interpretation

HOTlc

W\ensuremath{\infty}(S) \ensuremath{\simeq}Ho(S/|N(S)|). For a direct model-categorical comparison see Cisinski (2009, Example 1.7, equation (1.7.1)). This is a homotopy category of a slice, and not generally the ordinary slice of a homotopy category. The distinction already matters for a base with nontrivial fundamental group, as the following example shows.

These interpretations explain what became of the full-faithfulness question: for spaces it is settled by the reflector, and more general coefficient theories have their own locally constant subcategories and reflection results. They do not justify identifying an arbitrary W-coarse localization with that subcategory. Such an identification requires the additional exactness properties discussed on pp. 43--45. The existence of a locally constant reflection and the adequacy of a particular coarse test are separate assertions.

\end{dcomment}


\dgrblock{dgrb229}{Three theories within the comparison diagram}
\begin{dcomment}{exposition}{p70:J'ai envie de grouper ce qui a été prouvé}{Three theories within the comparison diagram}
Grothendieck now wants the many presentations to express one intelligible distinction. The proof that supplies it has already been explained in the commentary at manuscript pp.~64, 67, and 69: the path replacement defines a section of the coarse localization, and the two iterated units agree after fine localization. Here the new work is to organize the consequences and identify the image of that section.

For a map $v:X\to Y$ over $S$, the three tests are
\[
\begin{array}{c|l}
 W_S&v\in W\text{ on total categories},\\
 W_S^{\mathrm d}&s\backslash X\longrightarrow s\backslash Y\in W
                    \text{ for every }s,\\
 W_S^{\mathrm g}&X/s\longrightarrow Y/s\in W
                    \text{ for every }s.
\end{array}
\]
The relative Theorem A supplies the inclusions of both fine classes in $W_S$. Thus the identity on $\Cat/S$ gives $\rho$ and $\sigma$; these are quotient functors between homotopy categories, before any assertion about adjoints. In the other direction one must first replace $X$: the functor $C$ sends a coarse equivalence to both kinds of fine equivalence. This is why the inclusion of locally constant objects is represented by a nontrivial construction on total categories.

The replacement by a split fibration does something different. Its fibre at $s$ is $s\backslash X$, so it preserves the right-fine information already present in $X$. Dually, the split opfibration has fibre $X/s$. Neither construction forces its transition maps to become equivalences. The path replacement does force this, and consequently belongs to the upper part of the diagram. Confusing these two purposes would make the distinction between $\HOT_W(S)$ and $\HOT_W^{\mathrm{lc}}(S)$ disappear before it has been proved.

Part~I is therefore a statement about coarse localization of all the presentations. Part~II identifies the appropriate fine localization along each side. Part~III says when a strict fibre is a valid test and when the two fine tests agree. The hypothesis L5 is used at precisely the passage from a coarse equivalence between homotopy fibrations to a fibrewise equivalence. It should not be silently dropped when the diagram is reused. The subsequent remarks at manuscript pp.~81--82 explicitly reopen this dependency.
\end{dcomment}


\textbf{Théorème-Scholie.} On considère le diagramme de catégories de la page 8 (\footnote[43]{On suppose sur $W$ les conditions a) b) c) de p.~45, de plus L5 (p.~43). Cf. commentaires p.~81.}).

\circled{I} Quand on munit toutes les 13 catégories du diagramme du localiseur $W_S$ ($W$-équivalence grossière, i.e. $W$-équivalence dans Cat, sans plus), \textit{tous} les foncteurs du diagramme induisent des équivalences sur les catégories localisées (\footnote[44]{De plus, si $\mathcal{F}$ est une sous-catégorie pleine de $\Cat/S$ qui contient une des deux sous-catégories $\ul{Fib}_0S$, $\ul{Cofib}_0S$ (ce qui inclut le cas des 8 catégories des diagrammes, autres que $\Cat/S$ elle-même), l'inclusion $(\mathcal{F},W_S|\mathcal{F})$ dans $(\Cat/S,W_S)$ induit une équivalence pour les localisés.}). Désignons par
\[
\HOT^{\mathrm{lc}}_W(S)
\]
‘la’ catégorie localisée commune.

\circled{II} Munissons les quatre catégories bordantes gauches
\[
(*)\qquad\qquad
\begin{tikzcd}
\Cat/S & \ul{Liss}\,S \arrow[l, hook'] & \ul{Fib}\,S \arrow[l, hook'] & \ul{Fibsc}\,S \arrow[l, hook']
\end{tikzcd}
\]
du localiseur $W_S^{\mathrm{d}}$ de la $W$-équivalence fine \textit{droite} sur $S$ ( $\begin{tikzcd}[row sep=small, column sep=small]
X \arrow[rr, "f"] \arrow[dr] & & Y \arrow[dl]\\
 & S &
\end{tikzcd}$ est dans $W_S^{\mathrm{d}}$ si et seulement si les $s\backslash X\longrightarrow s\backslash Y$ le sont \texttt{[plutôt, sss ils sont dans $W$]}, pour tout $s\in\Ob\,S$). Les foncteurs du diagramme, i.e. ceux de $(*)$, entre ces catégories, donc aussi leurs composés, induisent des équivalences entre les localisées (\footnote[45]{Plus généralement, ça marche pour toute sous-catégorie pleine de $\Cat/S$ qui contient $\ul{Fib}\,S$ (et dualement, avec $W_S^{\mathrm{g}}$ au lieu de $W_S^{\mathrm{d}}$, si $\mathcal{F}$ contient $\ul{Cofib}\,S$).}). Désignons par
\[
\HOT_W(S)
\]
‘la’ catégorie localisée commune.

Enoncé dual pour
\[
(*')\qquad\qquad
\begin{tikzcd}
\Cat/S & \ul{Prop}\,S \arrow[l, hook'] & \ul{Cofib}\,S \arrow[l, hook'] & \ul{Cofibsc}\,S, \arrow[l, hook']
\end{tikzcd}
\]

% foot (each page): Version 9 novembre 2001


\dgrblock{dgrb230}{Two computations of the reflection}
\begin{dcomment}{examples}{p70}{Two computations of the reflection}
The reflection records how the base arrows must become invertible. For the interval S = [1], a contravariant diagram is an arrow B \ensuremath{-}\ensuremath{\to}uA. Since |N[1]| is contractible, locally constant diagrams are equivalent to constant diagrams. The reflector is

\ensuremath{-}\ensuremath{\to}uA) 7\ensuremath{-}\ensuremath{\to}(A

\ensuremath{-}\ensuremath{\to}A), (B with unit given by u at 1 and the identity at 0. To verify its universal property at the coherent level, a map from B \ensuremath{\to}A to the constant diagram on K consists of maps A \ensuremath{\to}K, B \ensuremath{\to}K, and a specified homotopy relating the second to the composite through A. The space of choices of the second map together with such a homotopy is contractible over the first. Thus the mapping space is naturally equivalent to Map(A, K), exactly as required. This argument also includes empty values.

For example, the diagram \ensuremath{\varnothing}\ensuremath{\to}\ensuremath{*}becomes the constant point diagram. Its total Grothendieck construction is the category supported at 0, while that of the reflected diagram is [1]. Their total nerves are equivalent, although the unit fails to be an objectwise equivalence at 1. This is the earlier counterexample viewed as the unit of an adjunction. The inclusion of constant diagrams also has a right adjoint, sending B \ensuremath{\to}A to B: the dual mapping-space computation gives Map(K, B). The left and right adjoints are distinct constructions even for this smallest base.

For a second computation take the one-object groupoid S = BZ. Every arrow of the base is already invertible, so every coherent diagram on S is locally constant and the reflection is an equivalence. Nevertheless the coarse slice contains more morphisms than the ordinary slice of Ho(S). Consider the object z : \ensuremath{*}\ensuremath{\to}BZ given by a base point. In the infinity-categorical slice,

EndHo(S/BZ)(z) \ensuremath{\sim}= Z.

Map/BZ(z, z) \ensuremath{\simeq}\ensuremath{\Omega}BZ, Indeed, the underlying map \ensuremath{*}\ensuremath{\to}\ensuremath{*}is unique, but being a map over BZ includes a path from the chosen base point to itself. Homotopy classes of such paths are the integers. In the ordinary category Ho(S)/BZ, the same object has only its identity endomorphism, because the over-condition is imposed merely as an equality of homotopy classes.

The associated local system has fibre Z with its regular translation action. A translation-equivariant map Z \ensuremath{\to}Z is determined by the image n of 0 and is k 7\ensuremath{\to}k + n; composition adds the integers. This recovers the endomorphism calculation directly from local systems. Path completion of the point over BZ supplies the categorical analogue of its universal-cover replacement, recovering exactly the transport that the ordinary slice would forget.

\end{dcomment}
\begin{otherlanguage}{english}
\dgroverview{VI-a}{Chapter VI \textperiodcentered\ \textup{VI-a} \textperiodcentered\ Overview}
% Text recovered from pages 57--59 of the original complete reader PDF; mathematical notation flattened by extraction.
The problem of the chapter. Grothendieck starts from a category of coefficients, here the category of small categories itself, and asks how its homotopy theory should vary with a diagram shape. The value at a small category S is to be


HOTW (S) = Fun(Sop, Cat)[W \ensuremath{-}1


obj], where a natural transformation belongs to Wobj when all its components belong to W. The essential operation is to localize the category of diagrams before forgetting their coherence. A functor from Sop to the already localized category of coefficients generally retains less information. This is the same motivation that led to the diagram formulation in Grothendieck (1983, \S{}69), and that the letter to Thomason subsequently places at the centre of the derivator programme (Grothendieck, 1991). The discussions of coherence and evaluation in the commentary on Chapter~I, manuscript pages 7, 17--18 and 27, remain the background; the present chapter supplies categorical constructions rather than repeating those objections.


The first seventy manuscript pages, dated 14--15 November 1990, record the discovery of the correct comparison theorem. Thirteen categories of objects over S initially appear to offer interchangeable models. They include split fibrations, arbitrary fibrations, smooth and proper functors, homotopy fibrations, perfect functors, and all of Cat/S. The objects alone do not decide the homotopy theory. The permitted morphisms and the chosen weak equivalences matter just as much. The first theorem overlooks that distinction, is expressly withdrawn, and is replaced by a comparison between three theories. The manuscript’s revisions are part of the proof’s development, not noise to remove from it. The axioms and the four tests. The axiom list also changes. The original L4 is an absolute Theorem A: an aspherical functor is a weak equivalence. L4 bis, introduced on manuscript page 44, is the relative assertion about a map X \ensuremath{\to}Y over a third category S, tested on every map X/s \ensuremath{\to}Y/s. It is the latter formulation that supports the relative descent arguments. The initial L5 passes from fibrewise weak equivalences of Cartesian maps between fibrations to weak equivalences of total categories. The new L5 on page 43 goes in the opposite direction for maps between homotopy fibrations. Their common label must not obscure their different content. The commentary proves the initial claim that L1, L2 bis and the original L5 already imply the required elementary axioms, including relative descent, by a twisted-arrow argument for opposite invariance.


For a map f : X \ensuremath{\to}Y over S, the decisive tests are


coarse


f \ensuremath{\in}W,


fibrewise fs : Xs \ensuremath{\to}Ys belongs to W for every s,


left-fine


f/s : X/s \ensuremath{\to}Y/s belongs to W for every s,


right-fine


s\textbackslash{}f : s\textbackslash{}X \ensuremath{\to}s\textbackslash{}Y belongs to W for every s. Over a smooth object, the strict fibre compares to the undercategory; over a proper object, it compares to the overcategory. These comparisons and relative descent explain when fibrewise equivalence implies a universal equivalence of totals. The modern asphericity formulation retains precisely this mechanism (Maltsiniotis, 2005, Theorem 1.2, Corollary 4.16 and Theorem 5.4). Coarse equivalence alone does not give either fine condition on the external borders of the diagram. Over [1], the inclusion \{0\} \ensuremath{\to}[1] is a coarse equivalence of split fibrations, but its fibre at 1 is the map from the empty category to a point. Evaluation at that fibre proves that the two localizations themselves differ. The two comma resolutions. The functor \ensuremath{\Phi} replaces X \ensuremath{\to}S by the split fibration whose fibre at s is s\textbackslash{}X; \ensuremath{\Psi} gives the split opfibration with fibre X/s. Their units are categorical homotopy equivalences on total categories. More precisely, the unit of \ensuremath{\Phi} is right-fine, and the unit of \ensuremath{\Psi} is left-fine. Composition of the two arrows in a comma object gives an explicit contraction proving this stronger fact. Consequently the right-fine localization of Cat/S computes the contravariant diagram theory, while the left-fine localization computes its covariant counterpart.


The comparison with chosen splittings needs a genuine repair of the argument, but not a change in the desired theorem. A Cartesian functor need not preserve the selected Cartesian lifts strictly, and an arbitrary functor over the base need not preserve Cartesian arrows at all. The actual free split adjunction is


Homsplit/S(\ensuremath{\Phi}X, Y ) \ensuremath{\sim}= HomS(X, UY ).


Its counit supplies the comparison in the strict category. For arbitrary functors between fibrations the pullback retraction has a canonical, possibly noninvertible natural comparison. Natural transformations over S induce equality after fine localization, since the cylinder projection A \ensuremath{\times} [1] \ensuremath{\to}A is fine. This yields the natural isomorphism needed after localization. The exposition constructs the adjunction and the localized comparison in full; the historical and modern strictification references are separated from this particular elaboration (Gray, 1966; Emmenegger et al., 2024).


The covariant theory is HOTW (Sop). The additional outer opposite that appears in several passages of the edition is incorrect with the initial definition above: taking opposite total categories preserves the direction of functors between them. The corrected variance is used throughout the commentary, while both the original statements and the editorial brackets are preserved in the source. The strength of path completion. The two-sided path replacement gives a natural map X \ensuremath{\to}C(X) which is a weak homotopy equivalence, with C(X) \ensuremath{\to}S perfect. This immediately proves comparisons for the coarse localizations, without the new L5. To identify coarse and fibrewise equivalences on the internal categories, and then identify their common theory with a full subcategory of the fine theory, is stronger. That step is what the new L5 supplies.


Its restriction on W is severe. A degree-n covering of the circle is an equivalence for homology with coefficients in a ring where n is invertible, but its fibre has n components. The commentary realizes the covering as an action groupoid, so that perfectness and the strict fibre can be checked


directly. For integral homology it gives a different explicit example:


Z[1/2] \ensuremath{\rtimes}2 Z


  \ensuremath{-}\ensuremath{\to}BZ. B The fibre has first homology Z[1/2]. The mapping-torus sequence shows that the total projection is nevertheless an integral homology equivalence, because 1 \ensuremath{-}2 acts invertibly on that fibre homology. This calculation settles the tentative integral-coefficient assertion on manuscript page 43. The development label marks the supplied application and proof, with a bounded account of the literature checked; it claims no priority for the group or the homology calculation.


Finite homotopy truncations fail as well: a path fibration over K(A, n + 1) has contractible total space and fibre K(A, n). Retaining the standing condition W\ensuremath{\infty}\ensuremath{\subseteq}W \ensuremath{\subseteq}W0, the new L5 actually forces W = W\ensuremath{\infty}. Iterating path fibres shows that every W-aspherical category has connected iterated loop spaces, hence is weakly contractible; applying this to the fibres of a replacement of an arbitrary W-map proves the conclusion. The shortened annotation omitting the component condition cannot stand literally, since the maximal class of all functors is a counterexample. Cisinski’s later characterization explicitly includes the corresponding condition on components (Cisinski, 2004, Theorem 2.3.6). The reflection and its modern meaning. The final argument proves full faithfulness through an adjunction. Put


\ensuremath{\Sigma}0 = W d


A = Cat/S, S \ensuremath{\subseteq}\ensuremath{\Sigma} = WS. Path completion carries \ensuremath{\Sigma} into \ensuremath{\Sigma}0, its unit \ensuremath{\eta}X belongs to \ensuremath{\Sigma}, and \ensuremath{\eta}CX belongs to \ensuremath{\Sigma}0. These conditions give


P : A[\ensuremath{\Sigma}\ensuremath{-}1


0 ] \ensuremath{\rightleftarrows}A[\ensuremath{\Sigma}\ensuremath{-}1] : R, P \ensuremath{\dashv}R, with R fully faithful. The unit is induced by \ensuremath{\eta}, and the counit is its inverse in the coarse category. The commentary verifies both triangle identities. In particular, the equality C(\ensuremath{\eta}X) = \ensuremath{\eta}CX after fine localization follows formally from naturality and the stated invertibilities. Grothendieck’s concatenation of paths gives an additional geometric proof by providing a common inverse. There is no need to cancel \ensuremath{\eta}X in the fine category, where it need not be invertible.


For W = W\ensuremath{\infty}, the resulting interpretation is


HOTW\ensuremath{\infty}(S) \ensuremath{\simeq}Ho Fun(N(S)op, S),


HOTlc


W\ensuremath{\infty}(S) \ensuremath{\simeq}Ho Fun(\ensuremath{\Pi}\ensuremath{\infty}N(S), S) \ensuremath{\simeq}Ho(S/|N(S)|). Here S is the infinity-category of spaces, and the homotopy category is taken after forming coherent diagrams or the slice. The reflector is restriction after left Kan extension along groupoid completion. The commentary explains the comparison through Thomason’s model and straightening (Thomason, 1980; Lurie, 2009). It also computes the interval reflector and the endomorphisms of the point over BZ, which distinguish the homotopy category of the slice from the ordinary slice of the homotopy category.


The locally constant embedding and its two adjoints are now theorems in the derivator setting (Cisinski, 2009, Proposition 3.10 and Theorem 3.13). Their existence does not identify an arbitrary coarse W-localization with that subcategory. The surviving qualification is about hypotheses, not an unresolved version of the full-faithfulness problem for spaces. We also do not label the


relationship between the weakest absolute and relative axiom packages as a current open problem: this audit establishes the implications used here and records where a weaker package is insufficient for the printed proof. The construction of the path categories themselves is an input from outside these pages. The first part of the chapter ends in the opening clauses of the theorem--scholium on manuscript page 70; the continuation on pages 71--81 is commented in the second part.
\end{otherlanguage}
\dgrchunkend
